\documentclass[11pt,a4paper,oneside]{book}
% ==================================================================
% Preâmbulo comum a todos os capítulos
% ==================================================================
\usepackage{fix-cm}
\usepackage[utf8]{inputenc}
\usepackage{babel}
\babelprovide[import,main]{brazilian}
\usepackage{amsmath,amssymb,amsthm,mathtools}
\usepackage[a4paper,margin=2.5cm]{geometry}
\usepackage{microtype}
\usepackage{anyfontsize}
\usepackage{enumitem}
\usepackage{booktabs,array}
\usepackage{xcolor}
\usepackage{pifont}
\usepackage{listings}
\usepackage{pgfplots}
\pgfplotsset{compat=1.16}
\usetikzlibrary{shapes.geometric,arrows.meta,positioning,fadings}
\usepackage[most]{tcolorbox}
\usepackage[colorlinks=true,linkcolor=azul,urlcolor=laranja,citecolor=azul,filecolor=laranja,
  bookmarksnumbered=true,pdftitle={Cálculo Numérico Interativo},pdfauthor={Osmar Alessio}]{hyperref}

\definecolor{azul}{RGB}{20,60,150}
\definecolor{azulescuro}{RGB}{10,30,80}
\definecolor{laranja}{RGB}{200,90,0}
\definecolor{verde}{RGB}{0,120,60}
\definecolor{vermelho}{RGB}{200,30,30}
\definecolor{cL1}{RGB}{31,99,180}
\definecolor{cL2}{RGB}{27,138,59}
\definecolor{cL3}{RGB}{230,126,34}
\definecolor{cL4}{RGB}{148,103,189}

\sloppy
\setlength{\parskip}{0.4em}
\renewcommand{\arraystretch}{1.15}

\lstset{basicstyle=\ttfamily\small,keywordstyle=\color{azul}\bfseries,commentstyle=\color{gray},
  stringstyle=\color{verde},frame=single,rulecolor=\color{gray!50},backgroundcolor=\color{gray!4},
  columns=fullflexible,keepspaces=true,showstringspaces=false,breaklines=true}

% ------------------------------------------------------------------
% Links para os aplicativos do GeoGebra e para os textos complementares.
% Opção 1 (padrão): arquivos .ggb na pasta ggb/ e PDFs na pasta complementos/,
% ao lado do PDF do livro.
% Opção 2: depois de publicar os arquivos no geogebra.org, troque
% \ggbabrir por \ggbonline e informe o endereço de cada material.
% Os nomes dos arquivos de cada capítulo (\arq...) são definidos no início
% de cada arquivo da pasta capitulos/.
% ------------------------------------------------------------------
\newcommand{\ggbpasta}{ggb/}
\newcommand{\ggbabrir}[2]{\href{run:\ggbpasta #1}{#2}}
\newcommand{\ggbonline}[2]{\href{#1}{#2}}
\newcommand{\pdfabrir}[2]{\href{complementos/#1}{#2}}
\newcommand{\arqtxt}[1]{\texttt{\expandafter\detokenize\expandafter{#1}}}

\newtcolorbox{atividade}[2]{enhanced,breakable,colback=orange!6,colframe=laranja,
  fonttitle=\bfseries,coltitle=white,
  title={Atividade no GeoGebra: #1\hfill\ggbabrir{#2}{\textcolor{white}{[abrir o aplicativo]}}},
  after upper={\par\smallskip\small\textbf{Arquivo:}\\ \ggbabrir{#2}{\texttt{ggb/\arqtxt{#2}}}}}

\newtcolorbox{leitura}[2]{enhanced,breakable,colback=teal!4,colframe=teal!70!black,
  fonttitle=\bfseries,coltitle=white,
  title={Leitura complementar: #1\hfill\pdfabrir{#2}{\textcolor{white}{[abrir o texto]}}},
  after upper={\par\smallskip\small\textbf{Arquivo:} \pdfabrir{#2}{\texttt{complementos/\arqtxt{#2}}}}}

\newtcolorbox{exemplo}[1][]{enhanced,breakable,colback=blue!3,colframe=azul!70,
  fonttitle=\bfseries,title={Exemplo#1}}

\newtcolorbox{observacao}{enhanced,breakable,colback=green!4,colframe=verde!70,
  fonttitle=\bfseries,title={Observação}}

\theoremstyle{definition}
\newtheorem{definicao}{Definição}[chapter]
\theoremstyle{plain}
\newtheorem{teorema}[definicao]{Teorema}
\newtheorem{proposicao}[definicao]{Proposição}
\newtheorem{corolario}[definicao]{Corolário}

\tcolorboxenvironment{teorema}{colback=azul!4,colframe=azul,boxrule=0.6pt,arc=1mm,breakable}
\tcolorboxenvironment{proposicao}{colback=azul!4,colframe=azul,boxrule=0.6pt,arc=1mm,breakable}
\tcolorboxenvironment{definicao}{colback=gray!5,colframe=gray!60,boxrule=0.6pt,arc=1mm,breakable}

\renewcommand{\qedsymbol}{$\blacksquare$}
\newcommand{\x}{\mathbf{x}}
\newcommand{\bb}{\mathbf{b}}
\newcommand{\y}{\mathbf{y}}

% Dados do livro (edite aqui)
\newcommand{\livrotitulo}{Cálculo Numérico Interativo}
\newcommand{\livrosubtitulo}{Aritmética de ponto flutuante, zeros de funções,\\sistemas lineares, interpolação, ajuste de curvas\\e integração numérica com aplicativos no GeoGebra}
\newcommand{\livroautor}{Osmar Alessio}
\newcommand{\livroinstituicao}{Universidade Federal do Triângulo Mineiro}
\newcommand{\livrocurso}{Licenciatura em Matemática -- ICENE}
\newcommand{\livrolocal}{Uberaba -- MG}
\newcommand{\livroano}{2026}


\begin{document}

\frontmatter
% ==================================================================
% Capa
% ==================================================================
\begin{titlepage}
\thispagestyle{empty}
\begin{tikzpicture}[remember picture,overlay]
  % fundo
  \shade[top color=azulescuro,bottom color=azul] (current page.north west) rectangle (current page.south east);
  % malha sutil
  \begin{scope}[shift={(current page.center)}]
    \draw[white,opacity=0.07,step=0.8cm] (-11,-15) grid (11,15);
  \end{scope}
  % faixa superior
  \fill[laranja] ([yshift=-1.1cm]current page.north west) rectangle ([yshift=-1.25cm]current page.north east);
  \node[anchor=north west,text=white,font=\sffamily\large] at ([xshift=2cm,yshift=-1.6cm]current page.north west)
    {\textsc{\livroinstituicao}};
  \node[anchor=north west,text=white!80,font=\sffamily] at ([xshift=2cm,yshift=-2.3cm]current page.north west)
    {\livrocurso};
  % título
  \node[anchor=north west,text=white,font=\sffamily\bfseries\fontsize{50}{56}\selectfont,align=left]
    at ([xshift=2cm,yshift=-4.2cm]current page.north west) {Cálculo\\Numérico\\Interativo};
  \node[anchor=north west,text=orange!60!yellow,font=\sffamily\Large,align=left]
    at ([xshift=2cm,yshift=-11.6cm]current page.north west) {\livrosubtitulo};
  % gráfico: polinômio interpolador, tangente de Newton e bisseção
  \begin{scope}[shift={([xshift=11.8cm,yshift=-19.6cm]current page.north west)},x=0.55cm,y=0.55cm]
    \draw[white,opacity=0.5,-{Stealth}] (-6,0)--(10.3,0);
    \draw[white,opacity=0.5,-{Stealth}] (0,-5.3)--(0,7.3);
    \draw[orange!70!yellow,line width=2.2pt,domain=-5.1:9.2,samples=120,smooth]
      plot (\x,{-9/5+97/120*\x+43/160*\x*\x-41/960*\x*\x*\x});
    \draw[white,line width=1pt,opacity=0.85] (5.3,{-9/5+97/120*5.3+43/160*5.3^2-41/960*5.3^3}) ++(-3.2,{-3.2*(97/120+43/80*5.3-41/320*5.3^2)}) -- ++(6.4,{6.4*(97/120+43/80*5.3-41/320*5.3^2)});
    \foreach \p/\q in {-4/2,-2/-2,4/3,8/0} \fill[white] (\p,\q) circle (4.5pt);
    \foreach \p/\q in {-4/2,-2/-2,4/3,8/0} \fill[azul] (\p,\q) circle (2.3pt);
    \draw[cyan!60,line width=1.4pt,|-|] (6.4,-1.0)--(9.6,-1.0);
    \draw[cyan!60,line width=1.4pt,|-|] (6.4,-1.8)--(8.0,-1.8);
    \draw[cyan!60,line width=1.4pt,|-|] (7.2,-2.6)--(8.0,-2.6);
  \end{scope}
  % matriz aumentada estilizada
  \node[text=white,opacity=0.6,font=\sffamily] at ([xshift=4.4cm,yshift=-23.2cm]current page.north west)
    {$\left(\begin{array}{rrrr|r}1&-4&16&-64&2\\0&2&-12&56&-4\\0&0&48&-96&17\\0&0&0&480&-\frac{41}{2}\end{array}\right)$};
  % rodapé
  \fill[white,opacity=0.08] ([yshift=4.2cm]current page.south west) rectangle (current page.south east);
  \node[anchor=south west,text=white,font=\sffamily\LARGE\bfseries] at ([xshift=2cm,yshift=2.4cm]current page.south west) {\livroautor};
  \node[anchor=south west,text=white!80,font=\sffamily\large] at ([xshift=2cm,yshift=1.4cm]current page.south west) {\livrolocal\quad\textbullet\quad\livroano};
  \node[anchor=south east,text=white,font=\sffamily\small,align=right] at ([xshift=-2cm,yshift=1.5cm]current page.south east)
    {Livro com aplicativos interativos:\\clique nos links para abrir no GeoGebra};
\end{tikzpicture}
\end{titlepage}

% ==================================================================
% Folha de rosto
% ==================================================================
\begin{titlepage}
\centering
{\large\textsc{\livroinstituicao}\par}
{\livrocurso\par}
\vspace{5cm}
{\sffamily\bfseries\Huge\color{azul}\livrotitulo\par}
\vspace{0.8cm}
{\Large \livrosubtitulo\par}
\vspace{3cm}
{\Large\bfseries \livroautor\par}
\vfill
{\large \livrolocal\par \livroano\par}
\end{titlepage}


\tableofcontents

\chapter{Apresentação}

Este livro reúne textos de Cálculo Numérico preparados para a Licenciatura em Matemática, com
uma característica central: cada método é acompanhado de um ou mais aplicativos no GeoGebra, em que
o aluno executa o método passo a passo, confere as contas e observa a construção geométrica.

\paragraph{Como usar os links.} Os quadros laranja \emph{Atividade no GeoGebra} trazem o link
\textbf{[abrir o aplicativo]}, que abre o arquivo correspondente no GeoGebra Classic. Os quadros
verdes \emph{Leitura complementar} abrem textos de apoio em PDF. Para que os links funcionem:
\begin{itemize}[nosep]
  \item mantenha o arquivo \texttt{livro.pdf} e as pastas \texttt{ggb} e \texttt{complementos} juntos,
        na mesma pasta;
  \item abra o PDF num leitor que permita abrir arquivos externos (Adobe Acrobat Reader, Okular,
        SumatraPDF, Foxit). Os leitores de PDF dos navegadores costumam bloquear esses links; nesse
        caso, abra os arquivos diretamente das pastas.
\end{itemize}

\paragraph{Organização.} O Capítulo~\ref{cap:pontoflutuante} trata da mudança de base e da
representação dos números em ponto flutuante (padrão IEEE~754), com máquinas de 8, 16, 32 e 64 bits
no GeoGebra, e dos erros de arredondamento e truncamento que acompanham todos os métodos seguintes. O
Capítulo~\ref{cap:zeros} apresenta os métodos para zeros de funções (bisseção, posição falsa, ponto
fixo, Newton-Raphson e secante) e sua relação com os laços de repetição. O
Capítulo~\ref{cap:sistemas} trata dos sistemas lineares, com métodos
diretos (eliminação de Gauss e fatoração $LU$) e iterativos (Gauss-Jacobi e Gauss-Seidel).  O
Capítulo~\ref{cap:interp} trata da interpolação polinomial por três caminhos (sistema linear,
Lagrange e Newton), que levam ao mesmo polinômio. O Capítulo~\ref{cap:ajuste} trata do ajuste de curvas
pelo método dos mínimos quadrados, resolvido de forma direta (sistema normal) e iterativa (gradiente
descendente). O Capítulo~\ref{cap:integracao} trata da integração numérica: regras dos trapézios
e de Simpson (simples e repetidas), extrapolação de Richardson e método de Romberg, e quadratura de
Gauss-Legendre. O Capítulo~\ref{cap:edo} trata das soluções numéricas de equações diferenciais
ordinárias: problemas de valor inicial (Euler, Runge-Kutta, sistemas e estabilidade) e problemas de
valor de contorno (método do tiro e diferenças finitas). Cada capítulo termina com exercícios e respostas.

\mainmatter
% Capítulo gerado a partir de ponto_flutuante.tex (arquivos usados neste capítulo)
\def\arqBase{MudancaDeBase.ggb}
\def\arqSisF{SistemaPontoFlutuante_reta.ggb}
\def\arqMoito{Maquina_PontoFlutuante_8bits.ggb}
\def\arqMdezesseis{Maquina_PontoFlutuante_16bits.ggb}
\def\arqMtrintadois{Maquina_PontoFlutuante_32bits.ggb}
\def\arqMsessentaquatro{Maquina_PontoFlutuante_64bits.ggb}
\providecommand{\bits}[1]{\texttt{#1}}
\chapter{Mudança de base e aritmética de ponto flutuante}\label{cap:pontoflutuante}

\begin{tcolorbox}[colback=gray!4,colframe=gray!50,title={Como usar os aplicativos e os textos deste capítulo},fonttitle=\bfseries]
Os quadros laranja \emph{Atividade no GeoGebra} trazem o link \textbf{[abrir o aplicativo]}.
São seis aplicativos. Os quatro primeiros substituem as planilhas de Excel das
\emph{máquinas} de 8, 16, 32 e 64 bits:
\begin{itemize}[nosep]
  \item máquinas de ponto flutuante: \ggbabrir{\arqMoito}{8 bits}, \ggbabrir{\arqMdezesseis}{16 bits},
        \ggbabrir{\arqMtrintadois}{32 bits (\emph{float})} e \ggbabrir{\arqMsessentaquatro}{64 bits (\emph{double})}.
        Cada bit é uma caixa de seleção; o aplicativo decodifica a palavra, codifica um número
        digitado (por arredondamento ou truncamento) e calcula os erros;
  \item \ggbabrir{\arqBase}{mudança de base}, com as divisões e as multiplicações sucessivas passo a passo;
  \item \ggbabrir{\arqSisF}{sistema $F(\beta,t,m,M)$ na reta}, que mostra todos os números de
        uma máquina pequena e o valor $fl(x)$ de qualquer $x$.
\end{itemize}

\end{tcolorbox}

% ==================================================================
\section{Por que estudar a representação dos números}

A resolução de um problema por métodos numéricos passa por várias fases: o problema real,
o levantamento de dados, a construção do modelo matemático, a escolha do método, a
implementação computacional e a análise dos resultados. Não é raro que o resultado final
fique distante do esperado, mesmo com todas as fases feitas corretamente. Os resultados
dependem também
\begin{itemize}[nosep]
  \item da precisão dos dados de entrada;
  \item da forma como esses dados são representados no computador;
  \item das operações numéricas efetuadas.
\end{itemize}

\begin{exemplo}[{: a área de um círculo}]
A área de um círculo de raio $100$\,m, calculada com $\pi\approx3{,}14$, $\pi\approx3{,}1416$ e
$\pi\approx3{,}141592654$, dá
\[
A=31\,400\ \text{m}^2,\qquad A=31\,416\ \text{m}^2,\qquad A=31\,415{,}9265\ \text{m}^2 .
\]
As diferenças vêm da aproximação usada para $\pi$, que é irracional. Nenhum computador obtém
a área ``exatamente''.
\end{exemplo}

\begin{exemplo}[{: um somatório}]
Calcule $S=\sum_{i=1}^{30\,000}x_i$ com $x_i=0{,}5$ e com $x_i=0{,}11$. As respostas exatas são
$15\,000$ e $3\,300$. Somando num programa:

\begin{center}
\begin{tabular}{lcc}
\toprule
 & $x_i=0{,}5$ & $x_i=0{,}11$\\
\midrule
precisão simples (32 bits) & $15\,000$ & $3\,300{,}985$\\
precisão dupla (64 bits)   & $15\,000$ & $3\,300{,}000\,000\,000\,628\,5$\\
\bottomrule
\end{tabular}
\end{center}

Com $0{,}5$ o resultado é exato; com $0{,}11$, não. A explicação está neste capítulo:
$0{,}5=(0{,}1)_2$ tem representação binária finita, mas $0{,}11$ não tem. Cada parcela já entra
com um pequeno erro, e cada soma acrescenta outro.
\end{exemplo}



Há três fontes de erro num cálculo numérico:
\begin{enumerate}[nosep]
  \item \textbf{incerteza dos dados}: os dados vêm de medidas com instrumentos de precisão finita;
  \item \textbf{erros de arredondamento}: vêm da forma finita de representar números na máquina;
  \item \textbf{erros de truncamento}: aparecem quando um processo infinito (uma série, um limite,
        uma integral) é substituído por um processo finito.
\end{enumerate}
Este capítulo trata do segundo tipo.

% ==================================================================
\section{Sistemas de numeração posicionais}

A \emph{base} de um sistema de numeração é a quantidade de símbolos usados para escrever
os números.

\begin{center}
\begin{tabular}{lll}
\toprule
Base & Nome & Símbolos\\
\midrule
2  & binário      & 0, 1\\
8  & octal        & 0, 1, 2, 3, 4, 5, 6, 7\\
10 & decimal      & 0, 1, \dots, 9\\
16 & hexadecimal  & 0, 1, \dots, 9, A, B, C, D, E, F\\
\bottomrule
\end{tabular}
\end{center}

A base 10 é a do nosso dia a dia. A base 2 é natural nos circuitos lógicos: 0 e 1 correspondem a
dois níveis de tensão (\emph{low} e \emph{high}), a falso e verdadeiro, a desligado e ligado.
As bases 8 e 16 servem para escrever números binários de forma compacta. Quando a base não é
10, ela vem indicada no índice: $(1010001011)_2$, $(453234)_8$, $(23A0D)_{16}$.

\begin{definicao}[notação posicional]
Num sistema posicional de base $\beta$, o número
$(d_{p-1}d_{p-2}\dots d_1d_0\,,\,d_{-1}\dots d_{-n})_\beta$, com $0\le d_i\le\beta-1$, vale
\[
D=\sum_{i=-n}^{p-1}d_i\,\beta^{\,i}
 =d_{p-1}\beta^{p-1}+\dots+d_1\beta+d_0+d_{-1}\beta^{-1}+\dots+d_{-n}\beta^{-n},
\]
em que $p$ é o número de dígitos à esquerda da vírgula e $n$ o número de dígitos à direita.
O valor de cada símbolo depende da sua posição.
\end{definicao}

Em base 2, cada dígito é um \emph{bit}. O bit mais à esquerda é o MSB
(\emph{most significant bit}) e o mais à direita é o LSB (\emph{least significant bit}).

\begin{center}
\small
\begin{tabular}{cccc|cccc}
\toprule
hex & dec & bin & & hex & dec & bin\\
\midrule
0 & 0 & 0000 && 8 & 8  & 1000\\
1 & 1 & 0001 && 9 & 9  & 1001\\
2 & 2 & 0010 && A & 10 & 1010\\
3 & 3 & 0011 && B & 11 & 1011\\
4 & 4 & 0100 && C & 12 & 1100\\
5 & 5 & 0101 && D & 13 & 1101\\
6 & 6 & 0110 && E & 14 & 1110\\
7 & 7 & 0111 && F & 15 & 1111\\
\bottomrule
\end{tabular}
\end{center}

% ==================================================================
\section{Da base \texorpdfstring{$\beta$}{beta} para a base 10}

Basta somar os dígitos multiplicados pelos pesos $\beta^{\,i}$.

\begin{exemplo}
\begin{align*}
(11011)_2  &=1\cdot2^4+1\cdot2^3+0\cdot2^2+1\cdot2^1+1\cdot2^0=16+8+2+1=27,\\
(372)_8    &=3\cdot8^2+7\cdot8+2=192+56+2=250,\\
(24{,}6)_8 &=2\cdot8+4+6\cdot8^{-1}=16+4+0{,}75=20{,}75,\\
(356)_{16} &=3\cdot16^2+5\cdot16+6=768+80+6=854,\\
(2AF)_{16} &=2\cdot16^2+10\cdot16+15=512+160+15=687,\\
(1000{,}011)_2&=8+0\cdot2^{-1}+1\cdot2^{-2}+1\cdot2^{-3}=8+0{,}25+0{,}125=8{,}375.
\end{align*}
\end{exemplo}

% ==================================================================
\section{Da base 10 para a base \texorpdfstring{$\beta$}{beta}}

O número é separado em parte inteira e parte fracionária, que são convertidas por processos
diferentes.

\subsection{Parte inteira: divisões sucessivas}

Seja $N$ inteiro com $(N)_{10}=(c_k\dots c_2c_1c_0)_\beta$. Então
\[
N=c_k\beta^k+\dots+c_2\beta^2+c_1\beta+c_0=\beta\,\underbrace{(c_k\beta^{k-1}+\dots+c_2\beta+c_1)}_{q_1}+c_0 .
\]
Logo $c_0$ é o \emph{resto} da divisão de $N$ por $\beta$ e $q_1$ é o quociente. Repetindo
o raciocínio com $q_1$, obtemos $c_1$, e assim por diante, até o quociente ser zero.

\begin{tcolorbox}[colback=azul!3,colframe=azul!60,title={Algoritmo: inteiro decimal $\to$ base $\beta$},fonttitle=\bfseries]
\begin{description}[nosep,font=\normalfont\bfseries]
  \item[Passo 0:] $k=0$, $N_0=N$.
  \item[Passo 1:] obtenha $q_k$ e $r_k$ tais que $N_k=\beta\,q_k+r_k$, $0\le r_k<\beta$; faça $c_k=r_k$.
  \item[Passo 2:] se $q_k=0$, pare. Caso contrário, faça $N_{k+1}=q_k$, $k=k+1$ e volte ao passo 1.
\end{description}
Os restos, lidos \textbf{de baixo para cima} (do último para o primeiro), formam o número na base $\beta$.
\end{tcolorbox}

\begin{exemplo}
\[
\begin{array}{rcl}
25&=&2\cdot12+\mathbf{1}\\
12&=&2\cdot6+\mathbf{0}\\
6&=&2\cdot3+\mathbf{0}\\
3&=&2\cdot1+\mathbf{1}\\
1&=&2\cdot0+\mathbf{1}
\end{array}
\ \Rightarrow\ 25=(11001)_2,
\qquad
\begin{array}{rcl}
2834&=&8\cdot354+\mathbf{2}\\
354&=&8\cdot44+\mathbf{2}\\
44&=&8\cdot5+\mathbf{4}\\
5&=&8\cdot0+\mathbf{5}
\end{array}
\ \Rightarrow\ 2834=(5422)_8 .
\]
Da mesma forma, $37=(100101)_2$ e $2834=16\cdot177+2$, $177=16\cdot11+1$, $11=16\cdot0+11$,
logo $2834=(B12)_{16}$.
\end{exemplo}

\subsection{Parte fracionária: multiplicações sucessivas}

Seja $0<r<1$ com $(r)_{10}=(0{,}d_1d_2d_3\dots)_\beta$. Multiplicando por $\beta$:
\[
\beta\,r=d_1+\underbrace{(0{,}d_2d_3\dots)_\beta}_{\text{parte fracionária}} .
\]
A parte inteira de $\beta r$ é $d_1$. Repetindo com a parte fracionária, obtemos $d_2$, $d_3$,\dots

\begin{tcolorbox}[colback=azul!3,colframe=azul!60,title={Algoritmo: fração decimal $\to$ base $\beta$},fonttitle=\bfseries]
\begin{description}[nosep,font=\normalfont\bfseries]
  \item[Passo 0:] $r_1=r$, $k=1$.
  \item[Passo 1:] calcule $\beta r_k$; faça $d_k=\lfloor\beta r_k\rfloor$ (parte inteira).
  \item[Passo 2:] faça $r_{k+1}=\beta r_k-d_k$. Se $r_{k+1}=0$, pare.
  \item[Passo 3:] faça $k=k+1$ e volte ao passo 1.
\end{description}
As partes inteiras, lidas \textbf{de cima para baixo}, formam $0{,}d_1d_2d_3\dots$ Se o processo não
terminar, paramos depois do número de dígitos desejado.
\end{tcolorbox}

\begin{exemplo}[{: $8{,}375$ e $4{,}8$ em binário}]
Parte inteira: $8=(1000)_2$ e $4=(100)_2$. Partes fracionárias:
\[
\begin{array}{rcll}
0{,}375\times2&=&\mathbf{0}{,}75\\
0{,}75\times2&=&\mathbf{1}{,}5\\
0{,}5\times2&=&\mathbf{1}{,}0&\text{(fim)}
\end{array}
\qquad\qquad
\begin{array}{rcll}
0{,}8\times2&=&\mathbf{1}{,}6\\
0{,}6\times2&=&\mathbf{1}{,}2\\
0{,}2\times2&=&\mathbf{0}{,}4\\
0{,}4\times2&=&\mathbf{0}{,}8&\text{(volta ao início)}
\end{array}
\]
Logo $8{,}375=(1000{,}011)_2$, uma expansão finita, e
$4{,}8=(100{,}\overline{1100})_2=(100{,}110011001100\dots)_2$, uma expansão infinita e periódica.
\end{exemplo}

\begin{exemplo}[{: $3{,}38$ e $57{,}3$ (Aula 01)}]
Para $0{,}38$, as multiplicações dão $0{,}76$; $1{,}52$; $1{,}04$; $0{,}08$; $0{,}16$; $0{,}32$; $0{,}64$; $1{,}28$;\dots,
e então $0{,}38\approx(0{,}01100001)_2=0{,}37890625$ com 8 dígitos. Portanto
$3{,}38\approx(11{,}01100001)_2$. Quanto mais dígitos, melhor a aproximação.

Para $57{,}3$: $57=(111001)_2$ e $0{,}3=(0{,}0\overline{1001})_2$, logo
$57{,}3=(111001{,}0100110011001\dots)_2$.
\end{exemplo}

\begin{observacao}
Uma fração $p/q$ irredutível tem expansão finita na base $\beta$ se, e somente se, todo fator primo
de $q$ divide $\beta$. Em base 2, só têm expansão finita os racionais com denominador potência de 2:
$1/8=(0{,}001)_2$, $1/16=(0{,}0001)_2$, $5{,}625=(101{,}101)_2$. Já $1/3$, $1/5$ e $1/10$ são periódicos:
\[
\tfrac13=(0{,}\overline{01})_2,\qquad \tfrac15=(0{,}\overline{0011})_2,\qquad
\tfrac1{10}=(0{,}0\overline{0011})_2 .
\]
Um número com expansão decimal finita, como $0{,}1$, pode ter expansão binária infinita.
\end{observacao}

\subsection{Entre as bases 2, 8 e 16}

Como $8=2^3$ e $16=2^4$, basta agrupar os bits de 3 em 3 (octal) ou de 4 em 4 (hexadecimal) a partir da
vírgula:
\[
(101101{,}101)_2=(\underbrace{101}_{5}\,\underbrace{101}_{5}{,}\underbrace{101}_{5})_2=(55{,}5)_8,
\qquad
(0010\,1101{,}1010)_2=(2D{,}A)_{16} .
\]
É por isso que as palavras de máquina são escritas em hexadecimal.

\begin{atividade}{mudança de base passo a passo}{\arqBase}
Digite um número decimal (até 9 casas decimais), escolha a base $\beta$ (de 2 a 16) e o número de
dígitos fracionários. O aplicativo mostra:
\begin{itemize}[nosep]
  \item à esquerda, as divisões sucessivas da parte inteira, com os restos em vermelho;
  \item no centro, as multiplicações sucessivas da parte fracionária;
  \item à direita, o resultado, a verificação pela notação posicional, o erro cometido ao parar e a
        conferência com o comando \texttt{ToBase} do GeoGebra;
  \item no quadro de baixo, a conversão de um inteiro de uma base qualquer para decimal
        (\texttt{FromBase}).
\end{itemize}
\textbf{Tarefas.} (a) Reproduza os exemplos $25$, $37$, $2834$ (bases 8 e 16), $8{,}375$, $4{,}8$,
$3{,}38$ e $57{,}3$.
(b) Converta $0{,}1$ para a base 2 com 4, 8, 16 e 24 dígitos e anote o erro em cada caso.
Quantos dígitos binários são necessários para que o erro fique abaixo de $10^{-6}$?
(c) Ache três números com expansão finita em base 2 e infinita em base 3.
\end{atividade}

% ==================================================================
\section{Representação em ponto flutuante}

Um número real $x\neq0$ pode ser escrito em \emph{ponto flutuante normalizado} na base $\beta$:
\[
x=\pm\,(0{,}d_1d_2\dots d_t)_\beta\times\beta^{\,e},\qquad d_1\neq0,\qquad -m\le e\le M,
\]
em que $0{,}d_1\dots d_t$ é a \emph{mantissa}, $t$ é o número de dígitos da mantissa e $e$ é o
\emph{expoente}, limitado por $-m$ e $M$. Esse conjunto de números é indicado por $F(\beta,t,m,M)$.
O zero é representado à parte, com mantissa nula.

Como a mantissa e o expoente têm tamanho finito:
\begin{itemize}[nosep]
  \item $F$ é um conjunto \emph{finito}, com $2(\beta-1)\beta^{t-1}(M+m+1)+1$ elementos;
  \item o menor positivo é $(0{,}10\dots0)_\beta\times\beta^{-m}=\beta^{-m-1}$, e o maior é
        $(0{,}(\beta{-}1)\dots(\beta{-}1))_\beta\times\beta^{M}=(1-\beta^{-t})\beta^{M}$;
  \item se $0<|x|<\beta^{-m-1}$, ocorre \textbf{underflow}; se $|x|>(1-\beta^{-t})\beta^{M}$,
        ocorre \textbf{overflow};
  \item os números de $F$ \textbf{não são igualmente espaçados}: entre $\beta^{e-1}$ e $\beta^{e}$ o
        espaçamento é $\beta^{e-t}$, que cresce com $e$.
\end{itemize}

\begin{exemplo}[{: $F(10,3,5,5)$}]
Com $t=3$ dígitos decimais e $-5\le e\le5$, o menor positivo é $0{,}100\times10^{-5}=10^{-6}$ e o maior é
$0{,}999\times10^{5}=99\,900$. Então
$x=235{,}89=0{,}23589\times10^{3}$ fica $0{,}236\times10^3$ (arredondado) ou $0{,}235\times10^3$
(truncado); $x=0{,}345\times10^{-7}$ dá underflow; $x=0{,}875\times10^{9}$ dá overflow.
\end{exemplo}

\begin{atividade}{o sistema $F(\beta,t,m,M)$ na reta}{\arqSisF}
Os controles deslizantes escolhem $\beta$, $t$, $m$ e $M$; os pontos azuis do eixo são \emph{todos}
os números do sistema. Arraste o ponto $P$: o ponto vermelho é $fl(x)$, por arredondamento ou,
desmarcando a caixa, por truncamento. O quadro mostra a quantidade de números, o maior, o menor
positivo, os erros absoluto e relativo e os avisos de underflow e overflow.

\textbf{Tarefas.} (a) Em $F(2,3,1,2)$, confira que há 33 números, que o maior é $3{,}5$ e o menor
positivo é $0{,}25$. (b) Amplie a região perto do zero (roda do mouse) e observe o ``buraco''
entre $0$ e $0{,}25$. (c) Compare o espaçamento entre $0{,}5$ e $1$ com o espaçamento entre $2$ e $3{,}5$.
(d) Mova $P$ por vários pontos e verifique que $|ER|$ nunca passa de $\tfrac12\beta^{1-t}$ com
arredondamento, enquanto $|EA|$ cresce com $|x|$.
\end{atividade}

% ==================================================================
\section{O padrão IEEE 754}

As mais diversas representações em ponto flutuante já foram usadas. O padrão IEEE~754, criado em 1985
e revisto em 2008 e 2019, é hoje implementado na grande maioria dos processadores. Em base 2, a
mantissa normalizada começa sempre com o bit 1, que por isso não precisa ser guardado (\emph{bit
escondido}). O padrão usa a forma $1{,}f$ e guarda três campos:
\[
\underbrace{\ s\ }_{\text{sinal}}\ \Big|\ \underbrace{\ \ E\ \ }_{\text{expoente com bias}}\ \Big|\ \underbrace{\qquad f\qquad}_{\text{fração}}
\qquad\qquad
x=(-1)^s\times(1{,}f)_2\times2^{E-\text{bias}} .
\]

\begin{itemize}[nosep]
  \item $s=0$ para positivo e $s=1$ para negativo;
  \item $E$ é um inteiro sem sinal; o expoente verdadeiro é $E-\text{bias}$, com
        $\text{bias}=2^{k-1}-1$ para um campo de $k$ bits (assim não é preciso guardar expoentes negativos);
  \item $f=(0{,}b_1b_2\dots b_p)_2$ é a fração; a mantissa é $1{,}f$, com $p+1$ bits de precisão.
\end{itemize}

\begin{center}
\begin{tabular}{lccccc}
\toprule
formato & total & expoente & fração & bias & faixa aproximada\\
\midrule
meia (\emph{binary16})            & 16  & 5  & 10  & 15    & $6{,}1\times10^{-5}$ a $65\,504$\\
simples (\texttt{float})          & 32  & 8  & 23  & 127   & $1{,}2\times10^{-38}$ a $3{,}4\times10^{38}$\\
dupla (\texttt{double})           & 64  & 11 & 52  & 1023  & $2{,}2\times10^{-308}$ a $1{,}8\times10^{308}$\\
estendida (Intel)                 & 80  & 15 & 63  & 16383 & $3{,}4\times10^{-4932}$ a $1{,}2\times10^{4932}$\\
quádrupla (\emph{binary128})      & 128 & 15 & 112 & 16383 & idem\\
\bottomrule
\end{tabular}
\end{center}

\subsection{Valores especiais e números subnormais}

Os valores extremos do campo $E$ são reservados:

\begin{center}
\begin{tabular}{lll}
\toprule
campo $E$ & fração $f$ & significado\\
\midrule
$00\dots0$ & $f=0$     & $\pm0$\\
$00\dots0$ & $f\neq0$  & subnormal: $x=(-1)^s\times(0{,}f)_2\times2^{\,1-\text{bias}}$\\
$1\le E\le 2^k-2$ & qualquer & normalizado: $x=(-1)^s\times(1{,}f)_2\times2^{E-\text{bias}}$\\
$11\dots1$ & $f=0$     & $\pm\infty$ (overflow, divisão por zero)\\
$11\dots1$ & $f\neq0$  & NaN, \emph{not a number} (por exemplo, $0/0$ ou $\sqrt{-1}$)\\
\bottomrule
\end{tabular}
\end{center}

Os \textbf{subnormais} (ou desnormalizados) preenchem o intervalo entre $0$ e o menor normalizado
$2^{1-\text{bias}}$ com números igualmente espaçados. O underflow fica ``gradual'', em vez de saltar
direto para zero.

\begin{exemplo}[{: $15\,213$ em precisão simples}]
$15\,213=(11101101101101)_2=(1{,}1101101101101)_2\times2^{13}$. Então
\begin{itemize}[nosep]
  \item $s=0$;
  \item $E=13+127=140=(10001100)_2$;
  \item $f=\bits{11011011011010000000000}$ (os 13 bits depois da vírgula, completados com zeros até 23).
\end{itemize}
A palavra é \bits{0\,10001100\,11011011011010000000000}, ou $466DB400_{16}$.
\end{exemplo}

\begin{exemplo}[{: $178{,}125$ em precisão dupla}]
$178{,}125=128+32+16+2+0{,}125=(10110010{,}001)_2=(1{,}0110010001)_2\times2^{7}$. Então
$s=0$, $E=7+1023=1030=(10000000110)_2$ e $f=\bits{0110010001}0\dots0$ (52 bits).
A palavra é $4066440000000000_{16}$.
\end{exemplo}

\begin{exemplo}[{: $0{,}2$ em todos os formatos (as planilhas das máquinas)}]
$0{,}2=(0{,}0\overline{0011})_2=(1{,}\overline{1001})_2\times2^{-3}$. O expoente é $-3$ e a fração
$1001\,1001\,1001\dots$ precisa ser cortada. Arredondando para o mais próximo:

\begin{center}
\small
\begin{tabular}{lccc}
\toprule
formato & palavra (hex) & valor guardado & erro relativo\\
\midrule
8 bits (bias 3)   & $0D$          & $0{,}203125$ & $-1{,}6\times10^{-2}$\\
16 bits           & $3266$        & $0{,}199951171875$ & $\phantom{-}2{,}4\times10^{-4}$\\
32 bits           & $3E4CCCCD$    & $0{,}20000000298023224$ & $-1{,}5\times10^{-8}$\\
64 bits           & $3FC999999999999A$ & $0{,}20000000000000001110$ & $-5{,}6\times10^{-17}$\\
\bottomrule
\end{tabular}
\end{center}

Na máquina de 8 bits, o expoente $-3$ é menor que o mínimo $1-3=-2$, e $0{,}2$ cai na faixa
subnormal: $0{,}2\approx(0{,}1101)_2\times2^{-2}=13/64$. Nas planilhas de Excel o campo $E=0$ era tratado
como normalizado. No padrão IEEE (e nos aplicativos) ele é reservado para o zero e os subnormais.
\end{exemplo}



% ==================================================================
\section{As máquinas de 8, 16, 32 e 64 bits}

As planilhas usadas em aula montavam um número bit a bit e calculavam
$(-1)^s(1+\text{fração})\,2^{E-\text{bias}}$. Os aplicativos a seguir fazem o mesmo e mais:
\begin{itemize}[nosep]
  \item cada bit é uma caixa de seleção (marcada $=1$), agrupada em sinal, expoente e fração, com o
        peso de cada bit indicado acima da caixa;
  \item a palavra aparece em binário e em hexadecimal;
  \item a \emph{decodificação} mostra $s$, $E$, $E-\text{bias}$, a fração, o tipo do número (zero,
        subnormal, normalizado, infinito ou NaN), o valor com 15 algarismos significativos e o
        espaçamento local (ulp, \emph{unit in the last place});
  \item na \emph{codificação}, você digita o número exato e escolhe \textbf{arredondamento} (para o mais
        próximo, empate para o par, como no padrão) ou \textbf{truncamento} (corte, em direção ao zero).
        O aplicativo marca os bits e calcula os erros absoluto e relativo;
  \item o quadro \emph{Faixa} mostra o maior número, o menor normalizado e o menor subnormal.
\end{itemize}

\subsection{A máquina de 8 bits}

O formato didático de 8 bits de W.~T.~Verts tem 1 bit de sinal, 3 de expoente
($\text{bias}=3$) e 4 de fração. Ele segue todas as regras do IEEE~754 e é pequeno o bastante para
listar os seus 256 padrões:
\begin{itemize}[nosep]
  \item $E=0$: zero e subnormais $(0{,}f)_2\times2^{-2}$, de $1/64=0{,}015625$ a $15/64$, com passo $1/64$;
  \item $1\le E\le6$: normalizados $(1{,}f)_2\times2^{E-3}$, de $0{,}25$ a $(1{,}1111)_2\times2^3=15{,}5$;
  \item $E=7$: $\pm\infty$ e NaN.
\end{itemize}
São $2\times(16+6\times16)=224$ padrões finitos (com $+0$ e $-0$). Entre $1$ e $2$ o passo é $1/16$;
entre $8$ e $15{,}5$ é $1/2$.

\begin{center}
\small
\begin{tabular}{cccl}
\toprule
$E$ & expoente & intervalo & passo\\
\midrule
0 (sub.) & $-2$ & $[0,\ 0{,}234375]$ & $1/64$\\
1 & $-2$ & $[0{,}25,\ 0{,}484375]$ & $1/64$\\
2 & $-1$ & $[0{,}5,\ 0{,}96875]$ & $1/32$\\
3 & $\phantom{-}0$ & $[1,\ 1{,}9375]$ & $1/16$\\
4 & $\phantom{-}1$ & $[2,\ 3{,}875]$ & $1/8$\\
5 & $\phantom{-}2$ & $[4,\ 7{,}75]$ & $1/4$\\
6 & $\phantom{-}3$ & $[8,\ 15{,}5]$ & $1/2$\\
\bottomrule
\end{tabular}
\end{center}

\begin{atividade}{máquina de 8 bits}{\arqMoito}
O aplicativo abre com $0{,}2$ arredondado: \bits{0 000 1101}, um subnormal que vale $0{,}203125$.

\textbf{Tarefas.} (a) Clique no bit de expoente de peso $2^0$ e veja o número passar a
$(1{,}1101)_2\times2^{-2}=0{,}453125$. (b) Monte o maior número ($15{,}5$), o menor normalizado ($0{,}25$),
o menor subnormal ($1/64$), $+\infty$ e um NaN. (c) Digite $x=2{,}7$ e codifique por arredondamento e por
truncamento: os resultados são $2{,}75$ e $2{,}625$. (d) Codifique $15{,}75$ por arredondamento (overflow,
$+\infty$) e por truncamento (fica o maior finito, $15{,}5$). (e) Codifique $0{,}007$, que é menor que
metade do menor subnormal: o resultado é zero (underflow).
\end{atividade}



\subsection{As máquinas de 16, 32 e 64 bits}

\begin{atividade}{máquina de 16 bits (meia precisão)}{\arqMdezesseis}
Formato $1+5+10$, $\text{bias}=15$. O aplicativo abre com $0{,}2$ arredondado
($3266_{16}$, $0{,}199951171875$).

\textbf{Tarefas.} (a) Monte $1$, $2$, $-1{,}5$ e o maior número finito ($65\,504$). (b) Codifique
$65\,519$ e $65\,520$ por arredondamento: o primeiro vira $65\,504$ e o segundo, $+\infty$.
(c) Codifique $2049$: o resultado é $2048$, porque com 11 bits de mantissa os inteiros acima de
$2^{11}$ não são todos representáveis. (d) Compare o erro relativo de $0{,}2$ com $\varepsilon_{\text{máq}}=2^{-10}$.
\end{atividade}

\begin{atividade}{máquina de 32 bits (precisão simples)}{\arqMtrintadois}
Formato $1+8+23$, $\text{bias}=127$, o \texttt{float} das linguagens de programação.

\textbf{Tarefas.} (a) Codifique $15\,213$ e confira a palavra $466DB400_{16}$. (b) Codifique $0{,}1$ e
$0{,}2$ e anote os erros relativos; compare com $2^{-24}\approx5{,}96\times10^{-8}$.
(c) Codifique $16\,777\,217=2^{24}+1$ e explique o resultado. (d) Codifique $10^{-40}$: o número é
subnormal e o erro relativo cresce, porque sobram menos bits significativos.
\end{atividade}

\begin{atividade}{máquina de 64 bits (precisão dupla)}{\arqMsessentaquatro}
Formato $1+11+52$, $\text{bias}=1023$, o \texttt{double}. Os 52 bits da fração estão em duas linhas
(use as barras de rolagem se a tela for pequena). Em hexadecimal, a palavra tem 16 dígitos e a fração 13.

\textbf{Tarefas.} (a) Codifique $178{,}125$ e confira $4066440000000000_{16}$. (b) Codifique $0{,}1$:
a fração é $999999999999A_{16}$; o último dígito ``A'' (em vez de 9) vem do arredondamento para cima.
(c) Compare a palavra de $0{,}1$ com a de $0{,}2$: só muda o expoente. Por quê?
(d) O GeoGebra calcula em precisão dupla, então o número digitado já entra como \texttt{double}: aqui o
erro aparece como $0$. Use o ulp e $\varepsilon_{\text{máq}}=2^{-52}$ para estimar o erro, ou faça a mesma conta na
máquina de 32 bits.
\end{atividade}

% ==================================================================
\section{Arredondamento e truncamento}

Considere um sistema que opera em aritmética de ponto flutuante de $t$ dígitos na base 10, e escreva
\[
x=f_x\times10^{e}+g_x\times10^{e-t},\qquad 0{,}1\le f_x<1,\quad 0\le g_x<1 .
\]
Por exemplo, com $t=4$ e $x=234{,}57$: $x=0{,}2345\times10^{3}+0{,}7\times10^{-1}$, logo $f_x=0{,}2345$ e $g_x=0{,}7$.
A parcela $g_x\times10^{e-t}$ não cabe na mantissa. Há dois critérios:

\begin{description}[nosep]
  \item[Truncamento:] $g_x$ é desprezado e $\bar x=f_x\times10^{e}$. Então
  \[
  |EA_x|=|g_x|\times10^{e-t}<10^{e-t},\qquad
  |ER_x|=\frac{|g_x|\times10^{e-t}}{|f_x|\times10^{e}}<\frac{10^{e-t}}{0{,}1\times10^{e}}=10^{-t+1} .
  \]
  \item[Arredondamento:] se $|g_x|<\tfrac12$, despreza-se $g_x$; caso contrário, soma-se 1 ao último dígito de
  $f_x$. Então
  \[
  |EA_x|\le\tfrac12\times10^{e-t},\qquad |ER_x|<\tfrac12\times10^{-t+1} .
  \]
\end{description}

O arredondamento comete erros menores, mas exige mais operações; por isso algumas máquinas antigas
truncavam. O IEEE~754 adota como padrão o \textbf{arredondamento para o mais próximo, com empate para o
par} (o último bit fica 0). Ele também oferece arredondamento em direção a $0$ (truncamento), a $+\infty$ e a
$-\infty$, úteis na análise por intervalos.

Na base 2 com $p$ bits de fração, os mesmos argumentos dão o \textbf{épsilon da máquina}
$\varepsilon_{\text{máq}}=2^{-p}$, a distância de $1$ ao número seguinte da máquina, e
\[
\frac{|x-fl(x)|}{|x|}\le\tfrac12\varepsilon_{\text{máq}}=2^{-(p+1)}\ \ \text{(arredondamento)},
\qquad
\frac{|x-fl(x)|}{|x|}<\varepsilon_{\text{máq}}\ \ \text{(truncamento)},
\]
para $x$ na faixa dos normalizados. Em precisão dupla, $\varepsilon_{\text{máq}}=2^{-52}\approx2{,}2\times10^{-16}$,
ou seja, cerca de 16 algarismos decimais significativos. Em precisão simples,
$\varepsilon_{\text{máq}}=2^{-23}\approx1{,}2\times10^{-7}$, cerca de 7 algarismos.

% ==================================================================
\section{Erros absoluto e relativo}

\begin{definicao}
Seja $\bar x$ uma aproximação de $x$. O \emph{erro absoluto} é $EA_x=x-\bar x$ e o \emph{erro relativo} é
\[
ER_x=\frac{EA_x}{\bar x}=\frac{x-\bar x}{\bar x}
\]
(na prática, divide-se pelo valor aproximado, que é o conhecido; quando $x$ é conhecido, como nos aplicativos,
pode-se dividir por $x$).
\end{definicao}

Em geral só $\bar x$ é conhecido, e o que se obtém é um limitante do erro. Se
$\pi\in(3{,}14;\,3{,}15)$ e tomamos $\bar\pi$ nesse intervalo, então $|EA_\pi|<0{,}01$.

\begin{exemplo}[{: o erro relativo compara melhor}]
Seja $\bar x=2112{,}9$ com $|EA_x|<0{,}1$, isto é, $x\in(2112{,}8;\,2113)$, e $\bar y=5{,}3$ com
$|EA_y|<0{,}1$, isto é, $y\in(5{,}2;\,5{,}4)$. Os erros absolutos têm o mesmo limitante, mas
\[
|ER_x|<\frac{0{,}1}{2112{,}9}\approx4{,}7\times10^{-5},\qquad |ER_y|<\frac{0{,}1}{5{,}3}\approx0{,}02 .
\]
O número $x$ está representado com precisão muito maior que $y$.
\end{exemplo}

% ==================================================================
\section{Operações em ponto flutuante e propagação de erros}

Mesmo quando as parcelas são exatas, o resultado de uma operação precisa ser normalizado e arredondado (ou
truncado) para $t$ dígitos. O resultado aproximado $\overline{OP}$ é o que fica na memória.

\begin{exemplo}[{: adição e multiplicação com $t=4$}]
Sejam $x=0{,}937\times10^{4}$ e $y=0{,}1272\times10^{2}$.

\textbf{Adição.} É preciso alinhar as vírgulas: a mantissa do número de menor expoente é deslocada para a
direita tantas casas quanto a diferença dos expoentes:
\[
x+y=(0{,}937+0{,}001272)\times10^{4}=0{,}938272\times10^{4} .
\]
Com $t=4$: $\overline{x+y}=0{,}9383\times10^{4}$ (arredondamento) ou $0{,}9382\times10^{4}$ (truncamento).

\textbf{Multiplicação.}
$xy=(0{,}937\times0{,}1272)\times10^{6}=0{,}1191864\times10^{6}$, logo
$\overline{xy}=0{,}1192\times10^{6}$ (arredondamento) ou $0{,}1191\times10^{6}$ (truncamento).
\end{exemplo}

Supondo as parcelas representadas exatamente, o erro relativo de uma operação é
$|ER_{OP}|<10^{-t+1}$ no truncamento e $|ER_{OP}|<\tfrac12\times10^{-t+1}$ no arredondamento.

\subsection{Fórmulas de propagação}

Sejam $x=\bar x+EA_x$ e $y=\bar y+EA_y$. Desprezando os produtos de erros, que são muito pequenos:

\begin{center}
\renewcommand{\arraystretch}{2}
\begin{tabular}{lll}
\toprule
operação & erro absoluto & erro relativo\\
\midrule
$x+y$ & $EA_{x+y}=EA_x+EA_y$ & $ER_{x+y}=ER_x\dfrac{\bar x}{\bar x+\bar y}+ER_y\dfrac{\bar y}{\bar x+\bar y}$\\
$x-y$ & $EA_{x-y}=EA_x-EA_y$ & $ER_{x-y}=ER_x\dfrac{\bar x}{\bar x-\bar y}-ER_y\dfrac{\bar y}{\bar x-\bar y}$\\
$xy$  & $EA_{xy}\approx\bar x\,EA_y+\bar y\,EA_x$ & $ER_{xy}\approx ER_x+ER_y$\\
$x/y$ & $EA_{x/y}\approx\dfrac{\bar y\,EA_x-\bar x\,EA_y}{\bar y^{\,2}}$ & $ER_{x/y}\approx ER_x-ER_y$\\
\bottomrule
\end{tabular}
\end{center}

Na divisão, usa-se $\dfrac1{1+EA_y/\bar y}=1-\dfrac{EA_y}{\bar y}+\Big(\dfrac{EA_y}{\bar y}\Big)^2-\cdots$ e
desprezam-se as potências maiores que 1. A cada operação soma-se ainda o erro de arredondamento do próprio
resultado, limitado por $\tfrac12\varepsilon_{\text{máq}}$.

\subsection{Cancelamento}

Na subtração, os fatores $\bar x/(\bar x-\bar y)$ e $\bar y/(\bar x-\bar y)$ ficam enormes quando
$\bar x\approx\bar y$. Os algarismos iniciais iguais se cancelam e sobram os finais, justamente os
contaminados por erro. É o \textbf{cancelamento}.

\begin{exemplo}[{: reescrever para evitar o cancelamento}]
Com $t=6$ dígitos decimais, calcule $\sqrt{x+1}-\sqrt{x}$ para $x=10^{8}$:
$\sqrt{100\,000\,001}=10\,000{,}000\,05\approx0{,}100000\times10^{5}$ e $\sqrt{10^8}=0{,}100000\times10^{5}$, cuja
diferença é $0$. Racionalizando,
\[
\sqrt{x+1}-\sqrt{x}=\frac{1}{\sqrt{x+1}+\sqrt{x}}\approx\frac1{0{,}200000\times10^{5}}=0{,}500000\times10^{-4},
\]
que é o valor correto ($4{,}9999999875\times10^{-5}$) com 6 algarismos.
\end{exemplo}

\begin{exemplo}[{: derivada numérica de $\sin x$ em $x=1$}]
A derivada é aproximada por $\dfrac{\sin(1+h)-\sin 1}{h}$; o valor exato é $\cos1\approx0{,}5403023$. Em precisão dupla:

\begin{center}
\small
\begin{tabular}{ccc}
\toprule
$h$ & aproximação & erro\\
\midrule
$10^{-1}$  & $0{,}4973637525$ & $4{,}3\times10^{-2}$\\
$10^{-4}$  & $0{,}5402602314$ & $4{,}2\times10^{-5}$\\
$10^{-8}$  & $0{,}5403023029$ & $3{,}0\times10^{-9}$\\
$10^{-12}$ & $0{,}5403455461$ & $4{,}3\times10^{-5}$\\
$10^{-15}$ & $0{,}5551115123$ & $1{,}5\times10^{-2}$\\
$10^{-16}$ & $0$ & $5{,}4\times10^{-1}$\\
\bottomrule
\end{tabular}
\end{center}

Até $h\approx10^{-8}\approx\sqrt{\varepsilon_{\text{máq}}}$, o erro cai junto com $h$ (erro de truncamento da
fórmula). Depois, $\sin(1+h)$ e $\sin1$ ficam tão próximos que o cancelamento domina. Com
$h=10^{-16}$, $1+h$ é arredondado para $1$ e a diferença é zero. A diferença central
$\frac{\sin(1+h)-\sin(1-h)}{2h}$ atinge erro de $10^{-11}$ com $h\approx10^{-5}$.
\end{exemplo}

\begin{observacao}
Outros cuidados que vêm da aritmética finita:
\begin{itemize}[nosep]
  \item a adição não é associativa: $(x+y)+z$ pode ser diferente de $x+(y+z)$;
  \item somar um número muito pequeno a um muito grande pode não mudar nada: em precisão dupla,
        $1+2^{-53}=1$;
  \item não teste igualdade de números reais (\texttt{if x == 0.3}); use uma tolerância,
        $|x-0{,}3|<\varepsilon$. Em Python, \texttt{0.1 + 0.2} mostra \texttt{0.30000000000000004}.
\end{itemize}
\end{observacao}



% ==================================================================
\section{Exercícios}

\begin{enumerate}[label=\arabic*.]
  \item Converta para a base 10: (a) $(101101)_2$; (b) $(0{,}1011)_2$; (c) $(157)_8$; (d) $(3A{,}8)_{16}$.
  \item Converta $45{,}625$ para as bases 2, 8 e 16.
  \item Mostre que $0{,}1=(0{,}0\overline{0011})_2$. Quantos bits seriam necessários para guardar
        $0{,}1$ exatamente?
  \item Represente $13{,}25$ em precisão simples (32 bits) e dê a palavra em hexadecimal.
  \item Decodifique a palavra de 32 bits \bits{0 10000001 01000000000000000000000}.
  \item Na máquina de 8 bits ($1+3+4$, bias 3), decodifique:
        (a) \bits{0 101 0110}; (b) \bits{1 001 1000}; (c) \bits{0 000 0011}; (d) \bits{0 111 0000}.
  \item Na máquina de 8 bits, codifique $x=2{,}7$ por arredondamento e por truncamento e calcule os erros
        absoluto e relativo.
  \item Em $F(2,3,1,2)$, quantos números existem? Qual o maior e qual o menor positivo?
  \item Com $t=4$ e arredondamento, calcule $(x+y)-x$ para $x=0{,}3491\times10^{4}$ e $y=0{,}2345\times10^{0}$.
        Compare com $y$.
  \item Explique por que, em precisão dupla, $1+2^{-53}=1$, mas $1+2^{-52}\neq1$.
  \item Calcule $\sqrt{x+1}-\sqrt{x}$ para $x=10^8$ com $t=6$ dígitos, diretamente e após racionalizar.
\end{enumerate}

\subsection*{Respostas}
\begin{enumerate}[label=\arabic*.,nosep]
  \item (a) $45$; (b) $0{,}6875$; (c) $111$; (d) $58{,}5$.
  \item $(101101{,}101)_2=(55{,}5)_8=(2D{,}A)_{16}$.
  \item As multiplicações dão a sequência $0{,}2\to0{,}4\to0{,}8\to0{,}6\to0{,}2$, que se repete. Nenhuma
        quantidade finita de bits basta.
  \item $13{,}25=(1101{,}01)_2=(1{,}10101)_2\times2^3$: $s=0$, $E=130=(10000010)_2$,
        $f=\bits{10101}0\dots0$; palavra $41540000_{16}$.
  \item $(1{,}01)_2\times2^{129-127}=1{,}25\times4=5$.
  \item (a) $(1{,}0110)_2\times2^{2}=5{,}5$; (b) $-(1{,}1)_2\times2^{-2}=-0{,}375$; (c) subnormal
        $3/64=0{,}046875$; (d) $+\infty$.
  \item $2{,}7=(1{,}0101\,1001\dots)_2\times2^1$. Arredondamento: \bits{0 100 0110}$=2{,}75$, $EA=-0{,}05$,
        $ER\approx-1{,}85\times10^{-2}$. Truncamento: \bits{0 100 0101}$=2{,}625$, $EA=0{,}075$, $ER\approx2{,}78\times10^{-2}$.
  \item $2\cdot1\cdot4\cdot4+1=33$ números; maior $(0{,}111)_2\times2^2=3{,}5$; menor positivo $(0{,}100)_2\times2^{-1}=0{,}25$.
  \item $x+y=3491{,}2345\to0{,}3491\times10^{4}$, logo $(x+y)-x=0$, embora $y\neq0$: o $y$ foi ``absorvido''.
  \item O número seguinte a $1$ em precisão dupla é $1+2^{-52}$. O valor $1+2^{-53}$ está exatamente no meio e,
        pelo empate para o par, é arredondado para $1$.
  \item Diretamente: $0$. Racionalizando: $0{,}500000\times10^{-4}$ (exato: $4{,}9999999875\times10^{-5}$).
\end{enumerate}

\subsection*{Referências}
\begin{itemize}[nosep]
  \item M.~A.~G.~Ruggiero e V.~L.~R.~Lopes, \emph{Cálculo Numérico: aspectos teóricos e computacionais}, 2.ª ed., Makron Books.
  \item M.~L.~Overton, \emph{Numerical Computing with IEEE Floating Point Arithmetic}, SIAM, 2001.
  \item IEEE Std 754-2019, \emph{IEEE Standard for Floating-Point Arithmetic}.
\end{itemize}

% Capítulo gerado a partir de zeros_de_funcoes.tex (arquivos usados neste capítulo)
\def\arqBisA{metodoBisseccao_2026_corrigido.ggb}
\def\arqBisB{metodoBisseccao_2026_correto_and.ggb}
\def\arqPosA{metodoposicaofalsa_2026_correto.ggb}
\def\arqPosB{metodoposicaofalsa_2026_correto_and.ggb}
\def\arqFixo{metodopontofixoJAVASCRIPT_2026_final_corrigido.ggb}
\def\arqFixoC{metodopontofixoJAVASCRIPT_2026_CORRIGIDO_ex02_CONVERGINDO.ggb}
\def\arqFixoD{metodopontofixoJAVASCRIPT_2026_CORRIGIDO_ex02_DIVERGINDO.ggb}
\def\arqNewton{metodoNewtonRaphson_2026_corrigido.ggb}
\def\arqSec{metododasecanteJAVASCRIPT_2026_CORRIGIDO.ggb}
\def\pdfLacos{material_lacos_fluxograma.pdf}
\chapter{Zeros de funções}\label{cap:zeros}

\begin{tcolorbox}[colback=gray!4,colframe=gray!50,title={Como usar os aplicativos e os textos deste capítulo},fonttitle=\bfseries]
Os quadros laranja \emph{Atividade no GeoGebra} trazem o link \textbf{[abrir o aplicativo]}.
São nove aplicativos:
\begin{itemize}[nosep]
  \item bisseção: \ggbabrir{\arqBisA}{versão 1, em $[1{,}6;\,1{,}8]$} e \ggbabrir{\arqBisB}{versão 2, em $[1{,}2;\,2]$};
  \item posição falsa: \ggbabrir{\arqPosA}{versão 1} e \ggbabrir{\arqPosB}{versão 2, com tolerâncias};
  \item ponto fixo: \ggbabrir{\arqFixo}{$\varphi(x)=\frac{2(1+x)}{x+2}$},
        \ggbabrir{\arqFixoC}{exemplo que converge} e \ggbabrir{\arqFixoD}{exemplo que diverge};
  \item \ggbabrir{\arqNewton}{Newton-Raphson} e \ggbabrir{\arqSec}{secante}.
\end{itemize}
Em todos eles há três botões: \textbf{INICIAR} (restaura os dados iniciais), \textbf{Xk}
(faz uma iteração por clique, mostrando no gráfico a construção do método) e
\textbf{JAVASCRIPT} (executa o laço completo e mostra a tabela das iterações numa janela).

Os quadros verdes \emph{Leitura complementar} abrem dois textos de apoio:
\pdfabrir{\pdfLacos}{laços de repetição e fluxogramas}.

Para que os links funcionem, mantenha este PDF e as pastas \texttt{ggb} e \texttt{complementos}
juntos e abra o PDF num leitor que permita abrir arquivos externos (Adobe Acrobat Reader, Okular,
SumatraPDF, Foxit). Os leitores de PDF dos navegadores costumam bloquear esses links; nesse caso,
abra os arquivos diretamente das pastas.
\end{tcolorbox}

% ==================================================================
\section{O problema}

Queremos encontrar $\bar x$ tal que $f(\bar x)=0$, isto é, um \emph{zero} (ou \emph{raiz}) de
$f$. Para polinômios de grau até 4 há fórmulas, mas para equações como
\[
x\ln x-1=0,\qquad x^5-x^3+x^2-3=0,\qquad \cos x=x
\]
não há expressão fechada para a solução. Os métodos numéricos produzem uma sequência de
aproximações $x_0,x_1,x_2,\dots$ que converge para $\bar x$, e param quando a aproximação é
boa o bastante.

O trabalho tem duas fases:
\begin{enumerate}[nosep]
  \item \textbf{isolamento}: encontrar um intervalo $[a,b]$ que contenha uma única raiz
        (ou uma boa aproximação inicial);
  \item \textbf{refinamento}: melhorar a aproximação até a precisão desejada.
\end{enumerate}

Ao longo do capítulo, usaremos a função $f(x)=x^2-2$, cujo zero positivo é
$\sqrt2\approx1{,}414214$, para comparar os métodos. Como sabemos a resposta, fica fácil
observar a velocidade de cada um.

% ==================================================================
\section{Isolamento das raízes}

\begin{teorema}[Bolzano]
Se $f$ é contínua em $[a,b]$ e $f(a)\,f(b)<0$, então existe pelo menos um
$\bar x\in(a,b)$ com $f(\bar x)=0$. Se, além disso, $f'$ existe e não muda de sinal em
$(a,b)$, a raiz é única.
\end{teorema}

\begin{exemplo}[{: tabela de sinais}]
Para $f(x)=x^3-9x+3$:
\begin{center}
\begin{tabular}{c|ccccccccc}
$x$ & $-4$ & $-3$ & $-2$ & $-1$ & $0$ & $1$ & $2$ & $3$ & $4$\\ \hline
$f(x)$ & $-25$ & $3$ & $13$ & $11$ & $3$ & $-5$ & $-7$ & $3$ & $31$\\
sinal & $-$ & $+$ & $+$ & $+$ & $+$ & $-$ & $-$ & $+$ & $+$
\end{tabular}
\end{center}
Há mudanças de sinal em $[-4,-3]$, $[0,1]$ e $[2,3]$. Como $f$ tem grau 3, essas são as três raízes.
\end{exemplo}

\begin{exemplo}[{: a função dos aplicativos de bisseção}]
Para $f(x)=x\ln x-1$ (com $x>0$), temos $f(1)=-1<0$ e $f(2)=2\ln2-1\approx0{,}386>0$.
Como $f'(x)=\ln x+1>0$ em $[1,2]$, há uma única raiz nesse intervalo:
$\bar x\approx1{,}763223$. Graficamente, ela é a abscissa do ponto em que $\ln x=\frac1x$.
\end{exemplo}

\begin{observacao}
Mudança de sinal é condição suficiente, não necessária: uma raiz de multiplicidade par não
troca o sinal de $f$. Por exemplo, $f(x)=x^3-5x^2+8x-4=(x-1)(x-2)^2$ (a função dos aplicativos
de posição falsa) tem uma raiz dupla em $x=2$, onde o gráfico só tangencia o eixo. Os métodos de
intervalo (bisseção e posição falsa) não conseguem ``cercar'' essa raiz.
\end{observacao}

% ==================================================================
\section{Critérios de parada e laços de repetição}

Todos os métodos deste capítulo são \emph{iterativos}: repetem o mesmo cálculo até que um
critério de parada seja satisfeito. Dada uma tolerância $\varepsilon$, os critérios mais usados são
\[
\text{(i) } |x_k-x_{k-1}|<\varepsilon,\qquad
\text{(ii) } |f(x_k)|<\varepsilon,\qquad
\text{(iii) } \frac{|x_k-x_{k-1}|}{|x_k|}<\varepsilon,
\]
\[
\text{(iv) } |b_k-a_k|<\varepsilon\ \text{(métodos de intervalo)},
\]
sempre acompanhados de um número máximo de iterações, para evitar um laço infinito quando o
método não converge.

Nenhum critério isolado é perfeito: $|f(x_k)|$ pode ser pequeno longe da raiz (função muito
``achatada'') e $|x_k-x_{k-1}|$ pode ser pequeno sem que se esteja perto dela (convergência muito
lenta). Por isso, é comum combinar dois critérios, e a escolha entre \emph{e} e \emph{ou} importa:
\begin{itemize}[nosep]
  \item \textbf{continuar enquanto} $|b-a|>\varepsilon_1$ \textbf{ou} $|f(x_k)|>\varepsilon_2$:
        só para quando \emph{os dois} critérios forem satisfeitos (mais exigente). É o que fazem os
        aplicativos de bisseção, com as tolerâncias $tol_1$ e $tol_2$;
  \item \textbf{continuar enquanto} $|x_k-x_{k-1}|>\varepsilon$ \textbf{e} $|f(x_k)|>\varepsilon$:
        para quando \emph{um} deles for satisfeito (menos exigente). É o que fazem os botões
        JAVASCRIPT dos aplicativos de posição falsa, secante, Newton e ponto fixo.
\end{itemize}

\begin{figure}[h]
\centering
\begin{tikzpicture}[node distance=7mm and 10mm,
  bloco/.style={draw,rounded corners,fill=blue!6,minimum width=46mm,minimum height=7mm,align=center,font=\small},
  decisao/.style={draw,diamond,aspect=2.4,fill=orange!10,align=center,font=\small,inner sep=1pt},
  seta/.style={-{Stealth[length=2.5mm]}}]
\node[bloco] (ini) {Início: dados $f$, $x_0$ (ou $[a,b]$), $\varepsilon$, $N$};
\node[bloco,below=of ini] (k0) {$k\leftarrow0$};
\node[decisao,below=of k0] (teste) {critério de parada\\satisfeito ou $k\ge N$?};
\node[bloco,below=9mm of teste] (passo) {calcula $x_{k+1}$ pela fórmula do método};
\node[bloco,below=of passo] (inc) {$k\leftarrow k+1$};
\node[bloco,right=16mm of teste] (fim) {Fim: $\bar x\approx x_k$};
\draw[seta] (ini)--(k0); \draw[seta] (k0)--(teste);
\draw[seta] (teste)--node[right,font=\small]{não}(passo);
\draw[seta] (passo)--(inc);
\draw[seta] (inc.west)--++(-12mm,0)|-(teste.west);
\draw[seta] (teste)--node[above,font=\small]{sim}(fim);
\end{tikzpicture}
\caption{Fluxograma comum a todos os métodos iterativos: um laço \texttt{while}.}
\label{fig:fluxo}
\end{figure}

Em linguagem de programação, esse fluxograma é um laço \texttt{while}. O botão JAVASCRIPT dos
aplicativos de bisseção, por exemplo, executa (resumidamente):
\begin{lstlisting}[language=Java]
var k = 0;  var xk = (a + b) / 2;
while (F(xk) !== 0 && (Math.abs(b - a) > tol1 || Math.abs(F(xk)) > tol2) && k < 200) {
    if (F(a) * F(xk) > 0) { a = xk; } else { b = xk; }   // escolhe o novo intervalo
    k++;
    xk = (a + b) / 2;                                    // novo ponto medio
}
\end{lstlisting}
A condição \texttt{k < 200} é o número máximo de iterações, que impede um laço infinito.

\begin{leitura}{laços de repetição e fluxogramas (4 páginas)}{\pdfLacos}
Introdução aos laços \texttt{for} e \texttt{while}, com exemplos em Python e fluxogramas:
variável de controle, condição lógica, contador e acumulador, laço infinito, \texttt{break} e
\texttt{continue}. É a base para ler e escrever o código dos botões JAVASCRIPT dos aplicativos.
\end{leitura}

% ==================================================================
\section{Método da bisseção}

Partindo de $[a_0,b_0]$ com $f(a_0)f(b_0)<0$, calculamos o ponto médio
\[
x_k=\frac{a_k+b_k}{2}
\]
e ficamos com a metade em que há troca de sinal: se $f(a_k)f(x_k)<0$, então
$[a_{k+1},b_{k+1}]=[a_k,x_k]$; caso contrário, $[a_{k+1},b_{k+1}]=[x_k,b_k]$.

\begin{proposicao}[Estimativa do erro]
Como o intervalo é dividido ao meio a cada passo,
$|\bar x-x_k|\le\dfrac{b_k-a_k}{2}=\dfrac{b_0-a_0}{2^{k+1}}$. Para garantir erro menor que
$\varepsilon$ com o critério $b_k-a_k<\varepsilon$, bastam
\[
k>\log_2\frac{b_0-a_0}{\varepsilon}
\]
iterações, número que se conhece \emph{antes} de começar.
\end{proposicao}

\begin{exemplo}[{: três iterações para $f(x)=x^2-2$ em $[1,2]$}]
$f(1)=-1<0$ e $f(2)=2>0$.
\begin{center}
\begin{tabular}{cccccc}
\toprule
$k$ & $a_k$ & $b_k$ & $x_k=\frac{a_k+b_k}{2}$ & $f(x_k)$ & novo intervalo\\
\midrule
1 & 1 & 2 & 1{,}5 & $\phantom{-}0{,}25$ & $[1;\,1{,}5]$\\
2 & 1 & 1{,}5 & 1{,}25 & $-0{,}4375$ & $[1{,}25;\,1{,}5]$\\
3 & 1{,}25 & 1{,}5 & 1{,}375 & $-0{,}109375$ & $[1{,}375;\,1{,}5]$\\
\bottomrule
\end{tabular}
\end{center}
Em cada linha, o novo intervalo é escolhido pelo sinal de $f(x_k)$. Para obter
$b_k-a_k<10^{-3}$, seriam necessárias $\log_2(1/10^{-3})\approx9{,}97$, isto é, 10 iterações.
\end{exemplo}

\begin{figure}[h]
\centering
\begin{tikzpicture}
\begin{axis}[width=11cm,height=6cm,axis lines=middle,xmin=0.8,xmax=2.2,ymin=-1.3,ymax=2.3,
  xtick={1,1.25,1.375,1.5,2},xticklabels={$1$,$1{,}25$,$1{,}375$,$1{,}5$,$2$},
  ytick={-1,1,2},grid=both,grid style={gray!20},xlabel={$x$},ylabel={$y$}]
\addplot[domain=0.85:2.1,samples=100,very thick,azul]{x^2-2};
\addplot[only marks,mark=*,vermelho] coordinates {(1.5,0)(1.25,0)(1.375,0)};
\addplot[only marks,mark=*,azul] coordinates {(1,-1)(2,2)(1.5,0.25)(1.25,-0.4375)(1.375,-0.109375)};
\draw[thick,laranja,|-|] (axis cs:1,-1.1)--(axis cs:2,-1.1) node[right,font=\scriptsize]{$[a_0,b_0]$};
\draw[thick,laranja,|-|] (axis cs:1,-0.85)--(axis cs:1.5,-0.85);
\draw[thick,laranja,|-|] (axis cs:1.25,-0.6)--(axis cs:1.5,-0.6);
\draw[thick,laranja,|-|] (axis cs:1.375,-0.35)--(axis cs:1.5,-0.35);
\end{axis}
\end{tikzpicture}
\caption{Bisseção para $f(x)=x^2-2$: a cada passo o intervalo (em laranja) cai pela metade.}
\end{figure}

\begin{atividade}{bisseção (versão 1)}{\arqBisA}
A função é $f(x)=x\ln x-1$. O botão INICIAR usa $[a,b]=[1{,}6;\,1{,}8]$ e
$tol_1=tol_2=0{,}001$, que podem ser alteradas nas caixas $tol_1$ e $tol_2$. O botão Xk só avança
enquanto $f(a)f(b)<0$ e ($|b-a|>tol_1$ ou $|f(x_k)|>tol_2$).
\begin{enumerate}[nosep]
  \item Clique em INICIAR e depois em Xk repetidas vezes, anotando $a_k$, $b_k$, $x_k$ e
        $f(x_k)$. Confira: $x_0=1{,}7$; $x_1=1{,}75$; $x_2=1{,}775$; $x_3=1{,}7625$.
  \item Em que iteração o botão deixa de avançar? (Esperado: $k=8$, com $[a,b]=[1{,}7625;\,1{,}763281]$,
        $|b-a|\approx0{,}00078$ e $|f(x_k)|\approx0{,}00052$.) Compare com a estimativa
        $\log_2(0{,}2/0{,}001)\approx7{,}6$.
  \item Use o botão JAVASCRIPT para ver todas as iterações de uma vez. Ele lê $a$, $b$, $tol_1$ e
        $tol_2$ do aplicativo; teste com $tol_1=tol_2=10^{-6}$.
  \item Escolha com os controles deslizantes um intervalo sem troca de sinal, por exemplo
        $[2;3]$. O que acontece?
\end{enumerate}
\end{atividade}

\begin{atividade}{bisseção (versão 2, intervalo maior)}{\arqBisB}
Mesma função, mas o botão INICIAR usa $[a,b]=[1{,}2;\,2]$, e o texto mostra também os dados da
iteração inicial ($a_0$, $b_0$, $x_0$). Os dois arquivos de bisseção trazem ainda uma tabela na
Planilha (menu Exibir $\to$ Planilha) que executa o algoritmo célula a célula, com as colunas
$a$, $b$, $x_k$, $f(a)$, $f(x_k)$, $f(a)f(x_k)$, $a_k$, $b_k$ e $|b_k-a_k|$.
\begin{enumerate}[nosep]
  \item Quantas iterações são necessárias agora? (Esperado: $k=10$.) Por que são duas a mais que na
        versão 1? (Sugestão: compare $b_0-a_0$ nos dois casos.)
  \item Na Planilha, a primeira linha da tabela usa as células $B22$ ($a$) e $C22$ ($b$). Digite
        nelas $1$ e $2$ e acompanhe linha a linha: como cada linha obtém $a_k$ e $b_k$ a partir da
        anterior? A coluna $|b_k-a_k|$ é sempre dividida por 2?
\end{enumerate}
\end{atividade}

% ==================================================================
\section{Método da posição falsa}

Em vez do ponto médio, usamos o ponto em que a reta que liga $(a,f(a))$ a $(b,f(b))$ corta o
eixo $x$:
\[
x_k=\frac{a_k\,f(b_k)-b_k\,f(a_k)}{f(b_k)-f(a_k)}.
\]
A escolha do novo intervalo é igual à da bisseção. Quando $|f(a)|$ e $|f(b)|$ são muito
diferentes, $x_k$ tende a ficar mais perto da raiz do que o ponto médio.

\begin{exemplo}[{: três iterações para $f(x)=x^2-2$ em $[1,2]$}]
\begin{center}
\begin{tabular}{cccccc}
\toprule
$k$ & $a_k$ & $b_k$ & $x_k$ & $f(x_k)$ & novo intervalo\\
\midrule
1 & 1 & 2 & $4/3\approx1{,}333333$ & $-0{,}222222$ & $[1{,}333333;\,2]$\\
2 & 1{,}333333 & 2 & 1{,}400000 & $-0{,}040000$ & $[1{,}4;\,2]$\\
3 & 1{,}4 & 2 & 1{,}411765 & $-0{,}006920$ & $[1{,}411765;\,2]$\\
\bottomrule
\end{tabular}
\end{center}
Por exemplo, $x_1=\frac{1\cdot2-2\cdot(-1)}{2-(-1)}=\frac43$. Repare que a extremidade $b=2$
\emph{nunca muda}: como $f$ é convexa, todas as aproximações ficam do mesmo lado da raiz. Assim,
$b_k-a_k$ não tende a zero, e o critério de parada deve usar $|f(x_k)|$ ou $|x_k-x_{k-1}|$.
\end{exemplo}

\begin{atividade}{posição falsa (versão 1)}{\arqPosA}
A função é $f(x)=x^3-5x^2+8x-4=(x-1)(x-2)^2$, com $[a,b]=[0{,}5;\,1{,}7]$ no botão INICIAR. O
gráfico mostra a reta secante $g$ que define cada $x_k$.
\begin{enumerate}[nosep]
  \item Confira $x_1\approx1{,}636364$, $x_2\approx1{,}557282$, $x_3\approx1{,}463714$.
  \item Qual extremidade fica fixa? Avance 10 iterações: $x_{10}\approx1{,}022$, ainda longe de
        $\bar x=1$. Por que a convergência é tão lenta aqui? (Observe o formato do gráfico perto de
        $x=1$ e o valor de $f(a)$.)
  \item Tente um intervalo que contenha $x=2$, como $[1{,}5;\,2{,}5]$. O aplicativo avisa
        ``$f(a)f(b)>0$''. Explique usando a raiz dupla.
  \item O botão JAVASCRIPT resolve outro exemplo, $f(x)=x^2-2$ em $[0{,}5;\,2]$ com $\varepsilon=10^{-4}$.
        Compare a velocidade com a da função do gráfico.
\end{enumerate}
\end{atividade}

\begin{atividade}{posição falsa (versão 2, com tolerâncias)}{\arqPosB}
Mesma função, agora com $[a,b]=[0{,}2;\,1{,}5]$ e as tolerâncias $tol_1=tol_2=0{,}001$; o texto
mostra $a_0$, $b_0$, $x_0$ e os erros a cada passo.
\begin{enumerate}[nosep]
  \item Faça 5 iterações e confira $x_5\approx1{,}199928$ (com $a=0{,}2$ fixo).
  \item Compare $|b_k-a_k|$ e $|f(x_k)|$ ao longo das iterações. Qual dos dois critérios seria
        atingido primeiro?
\end{enumerate}
\end{atividade}

% ==================================================================
\section{Método do ponto fixo}

Reescrevemos $f(x)=0$ numa forma equivalente $x=\varphi(x)$ e iteramos
\[
x_{k+1}=\varphi(x_k).
\]
Um ponto com $\bar x=\varphi(\bar x)$ é um \emph{ponto fixo} de $\varphi$ e, portanto, um zero
de $f$. Graficamente, é a interseção de $y=\varphi(x)$ com a reta $y=x$, e as iterações formam o
diagrama em ``teia'' ou ``escada'' (Figura~\ref{fig:teia}).

\begin{teorema}[Convergência do ponto fixo]
Se $\varphi$ e $\varphi'$ são contínuas num intervalo $I$ que contém $\bar x$, $\varphi(I)\subset I$
e $|\varphi'(x)|\le L<1$ para todo $x\in I$, então, para qualquer $x_0\in I$, a sequência
$x_{k+1}=\varphi(x_k)$ converge para o único ponto fixo $\bar x$ em $I$, e
$|x_{k+1}-\bar x|\le L\,|x_k-\bar x|$.
Se $|\varphi'(\bar x)|>1$, as iterações se afastam de $\bar x$.
\end{teorema}

Se $\varphi'>0$ perto da raiz, as aproximações se aproximam de um só lado (escada); se
$\varphi'<0$, elas alternam em torno da raiz (teia). Quanto menor $|\varphi'(\bar x)|$, mais rápida
a convergência.

\begin{exemplo}[{: $f(x)=x^2-2$ com $\varphi(x)=\frac{x+2}{x+1}$ e $x_0=1$}]
De $x^2=2$ vem $x^2+x=x+2$, isto é, $x=\frac{x+2}{x+1}$. Temos
$\varphi'(x)=-\frac{1}{(x+1)^2}$, e em $[1,2]$ vale $|\varphi'(x)|\le\frac14<1$; além disso,
$\varphi([1,2])=[\frac43,\frac32]\subset[1,2]$. Logo, há convergência.
\begin{center}
\begin{tabular}{ccccc}
\toprule
$k$ & $x_k$ & $x_{k+1}=\varphi(x_k)$ & $f(x_{k+1})$ & $|x_{k+1}-x_k|$\\
\midrule
0 & 1 & 1{,}5 & $\phantom{-}0{,}25$ & 0{,}5\\
1 & 1{,}5 & 1{,}4 & $-0{,}04$ & 0{,}1\\
2 & 1{,}4 & $17/12\approx1{,}416667$ & $\phantom{-}0{,}006944$ & 0{,}016667\\
\bottomrule
\end{tabular}
\end{center}
Como $\varphi'<0$, os iterados alternam em torno de $\sqrt2$, e o erro cai cerca de
$|\varphi'(\sqrt2)|\approx0{,}17$ a cada passo.

\textbf{Nem toda $\varphi$ serve.} A mesma equação pode ser escrita como $x=\frac2x$. Agora
$|\varphi'(\sqrt2)|=1$, e com $x_0=1$ obtemos $2,1,2,1,\dots$: a sequência nunca se aproxima de
$\sqrt2$.
\end{exemplo}

\begin{figure}[h]
\centering
\begin{tikzpicture}
\begin{axis}[width=8cm,height=7cm,axis lines=middle,xmin=0.9,xmax=1.6,ymin=0.9,ymax=1.6,
  grid=both,grid style={gray!20},xlabel={$x$},ylabel={$y$},
  legend style={at={(0.02,0.98)},anchor=north west,font=\small}]
\addplot[domain=0.9:1.6,thick,gray]{x};
\addplot[domain=0.9:1.6,very thick,azul,samples=60]{(x+2)/(x+1)};
\legend{$y=x$,$y=\varphi(x)$}
\addplot[thick,vermelho,-{Stealth[length=2mm]}] coordinates {(1,0.9)(1,1.5)(1.5,1.5)(1.5,1.4)(1.4,1.4)(1.4,1.416667)(1.416667,1.416667)};
\addplot[only marks,mark=*,mark size=1.5pt] coordinates {(1,0.9)(1.5,0.9)(1.4,0.9)};
\end{axis}
\end{tikzpicture}
\caption{Diagrama em teia para $\varphi(x)=\frac{x+2}{x+1}$ a partir de $x_0=1$:
os iterados alternam em torno de $\sqrt2$.}
\label{fig:teia}
\end{figure}

\begin{atividade}{ponto fixo com $\varphi(x)=\frac{2(1+x)}{x+2}$}{\arqFixo}
Aqui $f(x)=x^2-2$ e $\varphi(x)=\frac{2(1+x)}{x+2}$ (verifique que $x=\varphi(x)\iff x^2=2$). O
botão INICIAR usa $[a,b]=[1;3]$ e $x_0=b-0{,}3=2{,}7$, e o gráfico mostra a escada (vetores),
$\varphi$ e $|\varphi'|$. O texto indica ``Continue'' ou ``Pare'' conforme os critérios
$|x_{k+1}-x_k|>10^{-5}$ e $|f(x_{k+1})|>10^{-4}$.
\begin{enumerate}[nosep]
  \item Confira $x_1\approx1{,}574468$, $x_2\approx1{,}440476$, $x_3\approx1{,}418685$.
  \item Calcule $\varphi'(x)=\frac{2}{(x+2)^2}$. Por que a convergência é monótona (escada) e não
        alternada? Quanto vale $|\varphi'(\sqrt2)|$?
  \item Use os botões ``Apaga/Mostrar o Graf $f(x)$ e $\varphi'(x)$'' e verifique que
        $|\varphi'(x)|<1$ em todo o intervalo $[1,3]$.
  \item O botão JAVASCRIPT resolve outro exemplo: $x\ln x-1=0$ com $\varphi(x)=e^{1/x}$ e $x_0=1{,}3$.
        Mostre que $\varphi'(\bar x)=-1/\bar x\approx-0{,}57$ e explique por que essa convergência é
        alternada e bem mais lenta.
\end{enumerate}
\end{atividade}

\begin{atividade}{ponto fixo que converge e ponto fixo que diverge}{\arqFixoC}
Os dois aplicativos tratam $f(x)=x^2-x-2=(x-2)(x+1)$, com $x_0=2{,}7$, mas com funções de
iteração diferentes:
\begin{itemize}[nosep]
  \item \ggbabrir{\arqFixoC}{versão que converge}: $\varphi(x)=x-\dfrac{x^2-x-2}{2x-1}$;
  \item \ggbabrir{\arqFixoD}{versão que diverge}: $\varphi(x)=x^2-2$.
\end{itemize}
\begin{enumerate}[nosep]
  \item Na versão que converge, confira $x_1\approx2{,}111364$, $x_2\approx2{,}003848$,
        $x_3\approx2{,}000005$. Que método é esse? (Compare $\varphi$ com a fórmula de Newton.)
  \item Na versão que diverge, confira $x_1=5{,}29$ e $x_2\approx25{,}98$. Calcule $\varphi'(2)=4$
        e explique a divergência pelo teorema.
  \item Nas duas versões, ligue o gráfico de $|\varphi'|$ e compare-o com a reta $y=1$ perto de $x=2$.
\end{enumerate}
\small\textbf{Arquivo da versão que diverge:}\\ \ggbabrir{\arqFixoD}{\texttt{ggb/\arqtxt{\arqFixoD}}}
\end{atividade}

% ==================================================================
\section{Método de Newton-Raphson}

Aproximamos $f$ pela reta tangente em $x_k$ e tomamos o ponto em que ela corta o eixo $x$:
\[
x_{k+1}=x_k-\frac{f(x_k)}{f'(x_k)},\qquad f'(x_k)\neq0.
\]
É o método do ponto fixo com $\varphi(x)=x-\frac{f(x)}{f'(x)}$. Como $\varphi'(\bar x)=0$ quando
$f'(\bar x)\neq0$, a convergência perto de uma raiz simples é \emph{quadrática}: grosso modo, o
número de algarismos corretos dobra a cada iteração.

\begin{exemplo}[{: três iterações para $f(x)=x^2-2$ com $x_0=2$}]
$f'(x)=2x$, e então $x_{k+1}=x_k-\frac{x_k^2-2}{2x_k}$.
\begin{center}
\begin{tabular}{cccccc}
\toprule
$k$ & $x_k$ & $f(x_k)$ & $f'(x_k)$ & $x_{k+1}$ & $f(x_{k+1})$\\
\midrule
0 & 2 & 2 & 4 & 1{,}5 & 0{,}25\\
1 & 1{,}5 & 0{,}25 & 3 & 1{,}416667 & 0{,}006944\\
2 & 1{,}416667 & 0{,}006944 & 2{,}833333 & 1{,}414216 & $6\times10^{-6}$\\
\bottomrule
\end{tabular}
\end{center}
Em três iterações já temos $\sqrt2$ com 5 casas corretas.
\end{exemplo}

\begin{figure}[h]
\centering
\begin{tikzpicture}
\begin{axis}[width=10cm,height=6cm,axis lines=middle,xmin=1.2,xmax=2.15,ymin=-0.6,ymax=2.3,
  grid=both,grid style={gray!20},xlabel={$x$},ylabel={$y$},xtick={1.416667,1.5,2},
  xticklabels={$x_2$,$x_1$,$x_0$}]
\addplot[domain=1.2:2.1,samples=80,very thick,azul]{x^2-2};
\addplot[domain=1.45:2.1,thick,vermelho]{2+4*(x-2)};
\addplot[domain=1.38:1.75,thick,laranja]{0.25+3*(x-1.5)};
\addplot[dashed] coordinates {(2,0)(2,2)};
\addplot[dashed] coordinates {(1.5,0)(1.5,0.25)};
\addplot[only marks,mark=*] coordinates {(2,2)(1.5,0.25)(1.5,0)(1.416667,0)};
\end{axis}
\end{tikzpicture}
\caption{Newton para $f(x)=x^2-2$: cada nova aproximação é o zero da reta tangente.}
\end{figure}

\paragraph{Cuidados.} Newton exige $f'$ e um bom chute inicial. Se $f'(x_k)\approx0$, o passo
$f(x_k)/f'(x_k)$ fica enorme e a sequência pode ``saltar'' para longe, convergir para outra raiz
ou nem convergir.

\begin{atividade}{Newton-Raphson}{\arqNewton}
A função é $f(x)=x^5-x^3+x^2-3$, e o aplicativo calcula $f'$ e desenha a reta tangente $h$ em
cada passo. O botão INICIAR usa $x_0=1$; o chute também pode ser digitado na caixa $x_0$.
\begin{enumerate}[nosep]
  \item Confira $x_1=1{,}5$, $x_2\approx1{,}339130$, $x_3\approx1{,}286935$, $x_4\approx1{,}282090$,
        $x_5\approx1{,}282051$. Quantas casas corretas cada iteração acrescenta?
  \item Note que $x_1=1{,}5$ passou do outro lado da raiz. Por que, depois disso, todas as
        aproximações se aproximam pela direita? (Observe a concavidade de $f$.)
  \item Experimente $x_0=0$: a tangente é horizontal ($f'(0)=0$) e o método não pode ser aplicado.
        Com $x_0=0{,}3$, $f'(0{,}3)\approx0{,}34$ é pequeno e $x_1\approx8{,}22$; quantas iterações são
        necessárias para voltar perto da raiz?
  \item O botão JAVASCRIPT resolve $f(x)=x^3-x^2-x=0$ com $x_0=0{,}9$ e $\varepsilon=10^{-3}$.
        Como $f'(0{,}9)=-0{,}37$ é pequeno, $x_1\approx-1{,}751$, e o método converge para a raiz
        $\frac{1-\sqrt5}{2}\approx-0{,}618$, e não para a raiz mais próxima,
        $\frac{1+\sqrt5}{2}\approx1{,}618$. Explique o salto com a tangente.
\end{enumerate}
\end{atividade}

% ==================================================================
\section{Método da secante}

Para evitar o cálculo de $f'$, substituímos a tangente pela secante que passa pelos dois últimos
pontos:
\[
x_{k+1}=x_k-\frac{f(x_k)\,(x_k-x_{k-1})}{f(x_k)-f(x_{k-1})}.
\]
São necessários dois valores iniciais $x_0,x_1$. A fórmula lembra a da posição falsa, mas aqui
\emph{não} se mantém um intervalo com troca de sinal: usam-se sempre os dois últimos pontos. A
convergência é superlinear (ordem $\approx1{,}618$), mais rápida que a da posição falsa e um pouco
mais lenta que a de Newton, sem precisar de derivada.

\begin{exemplo}[{: três iterações para $f(x)=x^2-2$ com $x_0=1$, $x_1=2$}]
\begin{center}
\begin{tabular}{cccccc}
\toprule
iteração & $x_{k-1}$ & $x_k$ & $f(x_{k-1})$ & $f(x_k)$ & $x_{k+1}$\\
\midrule
1 & 1 & 2 & $-1$ & 2 & $x_2=4/3\approx1{,}333333$\\
2 & 2 & 1{,}333333 & 2 & $-0{,}222222$ & $x_3=1{,}4$\\
3 & 1{,}333333 & 1{,}4 & $-0{,}222222$ & $-0{,}04$ & $x_4\approx1{,}414634$\\
\bottomrule
\end{tabular}
\end{center}
Na terceira iteração, $f(x_3)$ e $f(x_4)$ têm o mesmo sinal, e mesmo assim o método continua:
não há intervalo a manter.
\end{exemplo}

\begin{atividade}{secante}{\arqSec}
A função é $f(x)=x^2-2$, com $x_0=a=0{,}5$ e $x_1=b=2$ no botão INICIAR. A cada clique em Xk,
o aplicativo traça a secante e substitui $(a,b)$ por $(b,x_{\text{novo}})$.
\begin{enumerate}[nosep]
  \item Confira $x_2=1{,}2$, $x_3=1{,}375$, $x_4\approx1{,}417476$, $x_5\approx1{,}414168$,
        $x_6\approx1{,}414214$.
  \item Em $x_2=1{,}2$ e $x_3=1{,}375$, os dois valores de $f$ são negativos. A posição falsa
        aceitaria esse par? O que isso mostra sobre a diferença entre os dois métodos?
  \item O botão JAVASCRIPT executa o mesmo exemplo com $\varepsilon=10^{-4}$. Compare o número de
        iterações com o da posição falsa no mesmo intervalo (botão JAVASCRIPT do aplicativo de
        posição falsa).
\end{enumerate}
\end{atividade}

% ==================================================================
\section{Comparação dos métodos}

Para $f(x)=x^2-2$, com o critério $|x_k-x_{k-1}|<10^{-6}$ (ou $b_k-a_k<10^{-6}$ na bisseção):

\begin{center}
\begin{tabular}{>{\raggedright\arraybackslash}p{2.6cm}>{\raggedright\arraybackslash}p{3.3cm}c>{\raggedright\arraybackslash}p{2.6cm}>{\raggedright\arraybackslash}p{3.4cm}}
\toprule
Método & Dados iniciais & Iterações & Ordem & Exige\\
\midrule
Bisseção & $[1,2]$ & 20 & linear (fator $\frac12$) & troca de sinal\\
Posição falsa & $[1,2]$ & 9 & linear & troca de sinal\\
Ponto fixo & $x_0=1$, $\varphi=\frac{x+2}{x+1}$ & 9 & linear (fator $|\varphi'(\bar x)|$) & $|\varphi'|<1$ perto da raiz\\
Newton & $x_0=2$ & 5 & quadrática & $f'$ e bom chute\\
Secante & $x_0=1$, $x_1=2$ & 6 & $\approx1{,}618$ & dois chutes\\
\bottomrule
\end{tabular}
\end{center}

\begin{itemize}[nosep]
  \item \textbf{Bisseção}: sempre converge (com troca de sinal) e o número de iterações é
        previsível, mas é lenta. Boa para obter um bom chute inicial.
  \item \textbf{Posição falsa}: em geral mais rápida que a bisseção, mas pode ficar lenta quando
        uma extremidade fica fixa (como no aplicativo com $(x-1)(x-2)^2$).
  \item \textbf{Ponto fixo}: simples, mas a convergência depende da escolha de $\varphi$.
  \item \textbf{Newton}: o mais rápido perto de raízes simples; exige a derivada e pode falhar com
        um chute ruim.
  \item \textbf{Secante}: quase tão rápida quanto Newton, sem derivada; também pode falhar longe da
        raiz.
\end{itemize}
Na prática, é comum combinar: alguns passos de bisseção para chegar perto da raiz e depois
Newton ou secante para refinar.



% ==================================================================
\section{Exercícios}

\begin{enumerate}[label=\textbf{\arabic*.},leftmargin=*]
\item Isole as raízes de $f(x)=x^3-9x+3$ por uma tabela de sinais e justifique que há exatamente
  uma raiz em cada intervalo encontrado.

\item Para $f(x)=e^{-x}-x$ em $[0,1]$:
  \begin{enumerate}[label=(\alph*),nosep]
  \item faça três iterações da bisseção, justificando a escolha de cada intervalo;
  \item quantas iterações garantem $b_k-a_k<10^{-4}$?
  \item faça três iterações da posição falsa e compare.
  \end{enumerate}

\item Aplique o método de Newton a $f(x)=\cos x-x$ com $x_0=1$ até que
  $|x_{k+1}-x_k|<10^{-6}$.

\item Aplique o método da secante a $f(x)=e^{-x}-x$ com $x_0=0$ e $x_1=1$ (quatro iterações).

\item A equação $x^3-x-1=0$ tem uma raiz em $[1,2]$. Considere
  $\varphi_1(x)=\sqrt[3]{x+1}$ e $\varphi_2(x)=x^3-1$.
  \begin{enumerate}[label=(\alph*),nosep]
  \item mostre que as duas são funções de iteração para essa equação;
  \item calcule $|\varphi_1'|$ e $|\varphi_2'|$ perto da raiz e decida qual converge;
  \item faça cinco iterações com a função adequada, a partir de $x_0=1$.
  \end{enumerate}

\item Escreva, em Python, um laço \texttt{while} para o método de Newton aplicado a
  $f(x)=x^2-2$, que pare quando $|x_{k+1}-x_k|<10^{-8}$ ou após 50 iterações. Desenhe o fluxograma
  correspondente.

\item \textbf{(GeoGebra)} No aplicativo de bisseção (versão 1), troque o critério de parada:
  use $tol_1=0{,}1$ e $tol_2=10^{-6}$. Em que iteração o método para? Qual dos dois critérios foi
  decisivo? Explique a diferença entre parar quando \emph{os dois} critérios são satisfeitos e quando
  \emph{um} deles é satisfeito.

\item \textbf{(GeoGebra)} Usando os aplicativos de ponto fixo, proponha uma terceira função de
  iteração para $x^2-x-2=0$, por exemplo $\varphi(x)=\sqrt{x+2}$, e preveja, calculando
  $\varphi'(2)$, se ela converge e se a convergência é monótona ou alternada.
\end{enumerate}

\subsection*{Respostas}
\begin{enumerate}[label=\textbf{\arabic*.},leftmargin=*,nosep]
\item Trocas de sinal em $[-4,-3]$, $[0,1]$ e $[2,3]$; como $f$ tem grau 3, cada intervalo tem
  exatamente uma raiz (também se verifica pelo sinal de $f'(x)=3x^2-9$ em cada um).
\item (a) $x_1=0{,}5$ ($f>0$, novo intervalo $[0{,}5;1]$); $x_2=0{,}75$ ($f<0$, $[0{,}5;0{,}75]$);
  $x_3=0{,}625$ ($f<0$, $[0{,}5;0{,}625]$). (b) $k>\log_2 10^4\approx13{,}3$: 14 iterações.
  (c) $x_1\approx0{,}612700$, $x_2\approx0{,}572181$, $x_3\approx0{,}567703$ (com $a=0$ fixo);
  a raiz é $\bar x\approx0{,}567143$.
\item $x_1\approx0{,}750364$, $x_2\approx0{,}739113$, $x_3\approx0{,}739085133$,
  $x_4\approx0{,}739085133$.
\item $x_2\approx0{,}612700$, $x_3\approx0{,}563838$, $x_4\approx0{,}567170$, $x_5\approx0{,}567143$.
\item (b) Na raiz $\bar x\approx1{,}3247$: $|\varphi_1'(\bar x)|\approx0{,}19<1$ (converge) e
  $|\varphi_2'(\bar x)|=3\bar x^2\approx5{,}3>1$ (diverge). (c) $1{,}259921$; $1{,}312294$;
  $1{,}322354$; $1{,}324269$; $1{,}324633$.
\item Uma solução:
\begin{lstlisting}[language=Python]
x, k = 2.0, 0
while k < 50:
    x_novo = x - (x**2 - 2) / (2*x)
    k = k + 1
    if abs(x_novo - x) < 1e-8:
        break
    x = x_novo
print(k, x_novo)
\end{lstlisting}
\item Para em $k=15$ ($x_k\approx1{,}7632233$, $|f(x_k)|\approx7\times10^{-7}$). O critério
  do intervalo ($|b-a|\le0{,}1$) é atingido já em $k=1$; o método só para quando também
  $|f(x_k)|\le10^{-6}$, portanto é o segundo critério que decide. Parar com \emph{um} critério
  economiza iterações, mas pode parar cedo demais.
\item $\varphi'(x)=\frac{1}{2\sqrt{x+2}}$, $\varphi'(2)=\frac14$: converge, de forma monótona
  (pois $\varphi'>0$).
\end{enumerate}

% Capítulo gerado a partir de sistemas_lineares.tex (arquivos usados neste capítulo)
\def\arqGauss{EliminacaoGauss_sistema4x5_compacto.ggb}
\def\arqLUtres{FatoracaoLU_3x3.ggb}
\def\arqLUquatro{FatoracaoLU_4x4.ggb}
\def\arqJacobi{GaussJacobi_sistema3x3_rolagem.ggb}
\def\arqSeidel{GaussSeidel_sistema3x3_rolagem.ggb}
\def\pdfCusto{CustoComputacional_SistemasLineares.pdf}
\def\pdfComp{ComparacaoMetodosDiretos.pdf}
\chapter{Sistemas lineares: métodos diretos e iterativos}\label{cap:sistemas}

\begin{tcolorbox}[colback=gray!4,colframe=gray!50,title={Como usar os aplicativos e os textos deste capítulo},fonttitle=\bfseries]
Os quadros laranja \emph{Atividade no GeoGebra} trazem o link \textbf{[abrir o aplicativo]},
que abre o arquivo correspondente no GeoGebra Classic:
\begin{itemize}[nosep]
  \item \ggbabrir{\arqGauss}{Eliminação de Gauss, sistema $4\times4$};
  \item \ggbabrir{\arqLUtres}{Fatoração $PA=LU$, matriz $3\times3$} e
        \ggbabrir{\arqLUquatro}{Fatoração $PA=LU$, matriz $4\times4$};
  \item \ggbabrir{\arqJacobi}{Método de Gauss-Jacobi, sistema $3\times3$};
  \item \ggbabrir{\arqSeidel}{Método de Gauss-Seidel, sistema $3\times3$}.
\end{itemize}
Os quadros verdes \emph{Leitura complementar} abrem dois textos de apoio sobre custo
computacional:
\pdfabrir{\pdfCusto}{Custo computacional na resolução de sistemas lineares} e
\pdfabrir{\pdfComp}{Comparação computacional dos métodos diretos}.

Para que os links funcionem, mantenha este PDF e as pastas \texttt{ggb} e \texttt{complementos}
juntos, na mesma pasta, e abra o PDF num leitor que permita abrir arquivos externos (Adobe
Acrobat Reader, Okular, SumatraPDF, Foxit). Os leitores de PDF dos navegadores costumam bloquear
esses links; nesse caso, abra os arquivos diretamente das pastas.

Em todos os aplicativos, o aluno digita os valores pedidos (multiplicadores ou componentes da
iteração) e recebe a correção imediata (\textcolor{verde}{\ding{51}}/\textcolor{vermelho}{\ding{55}});
cada etapa só é liberada quando a anterior está correta. Digite frações como \texttt{4/3}, sem arredondar.
\end{tcolorbox}

% ==================================================================
\section{Introdução}

Um \emph{sistema linear} de $n$ equações e $n$ incógnitas é
\[
\begin{cases}
a_{11}x_1+a_{12}x_2+\dots+a_{1n}x_n=b_1\\
a_{21}x_1+a_{22}x_2+\dots+a_{2n}x_n=b_2\\
\qquad\vdots\\
a_{n1}x_1+a_{n2}x_2+\dots+a_{nn}x_n=b_n
\end{cases}
\qquad\text{ou}\qquad A\x=\bb,
\]
com $A=(a_{ij})$ a \emph{matriz dos coeficientes}, $\x$ o vetor das incógnitas e $\bb$ o vetor
dos termos independentes. A matriz $[A\,|\,\bb]$ é a \emph{matriz aumentada}.

Sistemas lineares aparecem em quase todos os problemas de Cálculo Numérico: interpolação,
ajuste de curvas, discretização de equações diferenciais, circuitos, estruturas, economia.
Eles podem ser enormes (milhões de incógnitas), por isso importa não apenas \emph{resolver},
mas resolver com poucas operações e com controle dos erros de arredondamento.

\begin{definicao}
Um sistema linear é
\emph{possível e determinado} (SPD) se tem uma única solução,
\emph{possível e indeterminado} (SPI) se tem infinitas soluções e
\emph{impossível} (SI) se não tem solução. Para $A$ quadrada, o sistema é SPD se, e somente
se, $\det A\neq0$.
\end{definicao}

Estudaremos duas famílias de métodos:
\begin{itemize}[nosep]
  \item \textbf{métodos diretos}, que, em aritmética exata, fornecem a solução exata após um
  número finito de operações: eliminação de Gauss e fatoração $LU$;
  \item \textbf{métodos iterativos}, que geram uma sequência $\x^{(0)},\x^{(1)},\x^{(2)},\dots$
  que converge para a solução: Gauss-Jacobi e Gauss-Seidel.
\end{itemize}

% ==================================================================
\section{Sistemas triangulares}

Os métodos diretos reduzem o problema a sistemas triangulares, que são resolvidos facilmente.
Se $U$ é triangular superior com $u_{ii}\neq0$, o sistema $U\x=\y$ é resolvido de baixo para
cima (\emph{retrossubstituição}):
\[
x_n=\frac{y_n}{u_{nn}},\qquad
x_i=\frac{1}{u_{ii}}\Bigl(y_i-\sum_{j=i+1}^{n}u_{ij}x_j\Bigr),\quad i=n-1,\dots,1.
\]
Se $L$ é triangular inferior, $L\y=\bb$ é resolvido de cima para baixo (\emph{substituição
progressiva}):
\[
y_1=\frac{b_1}{l_{11}},\qquad
y_i=\frac{1}{l_{ii}}\Bigl(b_i-\sum_{j=1}^{i-1}l_{ij}y_j\Bigr),\quad i=2,\dots,n.
\]
Cada uma custa cerca de $n^2$ operações, muito menos que a eliminação ($\approx\frac23n^3$).

% ==================================================================
\section{Eliminação de Gauss}

A eliminação de Gauss transforma $[A\,|\,\bb]$ numa matriz triangular superior equivalente
usando operações elementares sobre as linhas, que não alteram a solução:
\begin{itemize}[nosep]
  \item trocar duas linhas: $L_i\leftrightarrow L_j$;
  \item substituir uma linha por ela menos um múltiplo de outra: $L_i\leftarrow L_i-m\,L_k$.
\end{itemize}

\paragraph{O algoritmo.} Na etapa $k$ ($k=1,\dots,n-1$), o elemento $a^{(k-1)}_{kk}$ é o
\emph{pivô}. Se ele for nulo, troca-se a linha $k$ com uma linha abaixo que tenha elemento não
nulo na coluna $k$. Depois, para cada linha $i>k$, calcula-se o \emph{multiplicador}
\[
m_{ik}=\frac{a^{(k-1)}_{ik}}{a^{(k-1)}_{kk}}
\qquad\text{e faz-se}\qquad L_i\leftarrow L_i-m_{ik}L_k,
\]
o que zera a coluna $k$ abaixo do pivô. Ao final, resolve-se o sistema triangular por
retrossubstituição.

\begin{exemplo}[: eliminação com trocas de linhas (exemplo inicial do aplicativo)]
Resolva
\[
[A\,|\,\bb]=\left(\begin{array}{rrrr|r}
0 & 1 & 2 & 1 & 4\\
1 & 1 & 1 & 0 & 2\\
2 & 1 & 0 & 1 & 2\\
1 & 0 & 1 & 1 & 4
\end{array}\right).
\]
\textbf{Etapa 1.} O pivô $a_{11}=0$; trocamos $L_1\leftrightarrow L_2$. Agora o pivô é $1$ e
$m_{21}=0$, $m_{31}=2$, $m_{41}=1$:
\[
B_1=\left(\begin{array}{rrrr|r}
1 & 1 & 1 & 0 & 2\\
0 & 1 & 2 & 1 & 4\\
0 & -1 & -2 & 1 & -2\\
0 & -1 & 0 & 1 & 2
\end{array}\right).
\]
\textbf{Etapa 2.} Pivô $1$; $m_{32}=-1$, $m_{42}=-1$:
\[
B_2=\left(\begin{array}{rrrr|r}
1 & 1 & 1 & 0 & 2\\
0 & 1 & 2 & 1 & 4\\
0 & 0 & 0 & 2 & 2\\
0 & 0 & 2 & 2 & 6
\end{array}\right).
\]
\textbf{Etapa 3.} O pivô $(3,3)$ é zero; trocamos $L_3\leftrightarrow L_4$, e então $m_{43}=0$:
\[
B_3=\left(\begin{array}{rrrr|r}
1 & 1 & 1 & 0 & 2\\
0 & 1 & 2 & 1 & 4\\
0 & 0 & 2 & 2 & 6\\
0 & 0 & 0 & 2 & 2
\end{array}\right).
\]
\textbf{Retrossubstituição:} $x_4=1$, $x_3=\frac{6-2}{2}=2$, $x_2=4-2\cdot2-1=-1$,
$x_1=2-(-1)-2=1$. Logo $\x=(1,-1,2,1)$, e o sistema é SPD.
\end{exemplo}

\paragraph{Classificação pela forma escalonada.} Se, ao final, aparece uma linha
$(0\ \cdots\ 0\,|\,c)$ com $c\neq0$, o sistema é \emph{impossível}; se aparece
$(0\ \cdots\ 0\,|\,0)$ e não há contradição, sobram variáveis livres e o sistema é
\emph{possível e indeterminado}.

\begin{atividade}{eliminação de Gauss $4\times4$}{\arqGauss}
O aplicativo começa com o sistema do exemplo acima. Em cada etapa, informe a linha do pivô
(trocando linhas quando o pivô for zero) e os multiplicadores; as matrizes $B_1$, $B_2$, $B_3$
mostram o pivô em azul e, em vermelho, os elementos que ainda deveriam ser zero.
\begin{enumerate}[nosep]
  \item Refaça o exemplo no aplicativo e confira a solução com ``Mostrar solução''.
  \item Clique em \emph{Exemplo: troca + SPI}. Em que etapa é necessária a troca? Qual linha
        da matriz final mostra que o sistema é indeterminado? Qual variável fica livre?
  \item Clique em \emph{Exemplo: sistema impossível (SI)}. Qual linha da matriz final contém a
        contradição?
  \item Gere sistemas com \emph{Novo sistema aleatório (SPD)} até conseguir resolver três deles
        sem erros.
\end{enumerate}
\end{atividade}

\paragraph{Pivoteamento parcial.} Na aritmética do computador, um pivô muito pequeno gera
multiplicadores enormes, que amplificam os erros de arredondamento. Por isso, na prática,
escolhe-se como pivô o elemento de \emph{maior módulo} da coluna (a partir da linha $k$) e
trocam-se as linhas. Assim $|m_{ik}|\le1$.

\begin{exemplo}[: por que pivotear]
Para $A=\begin{pmatrix}10^{-17}&1\\1&1\end{pmatrix}$ e $\bb=(1,2)$, a solução é
$\x\approx(1,1)$. Sem pivoteamento, $m_{21}=10^{17}$ e a segunda equação vira
$(1-10^{17})x_2=2-10^{17}$; em dupla precisão, $1-10^{17}$ e $2-10^{17}$ são arredondados
para $-10^{17}$, e o programa obtém $x_2=1$ e $x_1=0$, totalmente errado. Com pivoteamento
parcial, as linhas são trocadas, $m_{21}=10^{-17}$, e obtém-se $\x=(1,1)$.
\end{exemplo}

% ==================================================================
\section{Fatoração LU}

Na eliminação de Gauss, os multiplicadores são usados e depois descartados. A fatoração $LU$
os guarda numa matriz triangular inferior $L$; a matriz final escalonada é $U$.

\begin{teorema}[Fatoração LU]
Se a eliminação de Gauss pode ser feita sem trocas de linhas (todos os pivôs não nulos), então
\[
A=LU,\qquad
L=\begin{pmatrix}
1 & & & \\
m_{21} & 1 & & \\
\vdots & & \ddots & \\
m_{n1} & m_{n2} & \cdots & 1
\end{pmatrix},
\qquad U=\text{matriz escalonada final}.
\]
Em geral (com trocas), existe uma matriz de permutação $P$ tal que $PA=LU$.
\end{teorema}

Para resolver $A\x=\bb$, escrevemos $PA\x=P\bb$, isto é, $LU\x=P\bb$, e resolvemos dois
sistemas triangulares:
\[
L\y=P\bb\quad\text{(substituição progressiva)},\qquad U\x=\y\quad\text{(retrossubstituição)}.
\]
Além disso, $\det A=(-1)^{t}\,u_{11}u_{22}\cdots u_{nn}$, onde $t$ é o número de trocas.

\begin{exemplo}[: fatoração sem trocas ($P=I$)]
Para $A=\begin{pmatrix}2&1&1\\4&3&3\\8&7&9\end{pmatrix}$ e $\bb=(4,10,24)$:
etapa 1, $m_{21}=2$ e $m_{31}=4$; etapa 2, $m_{32}=3$. Então
\[
L=\begin{pmatrix}1&0&0\\2&1&0\\4&3&1\end{pmatrix},\qquad
U=\begin{pmatrix}2&1&1\\0&1&1\\0&0&2\end{pmatrix}.
\]
$L\y=\bb$: $y_1=4$, $y_2=10-2\cdot4=2$, $y_3=24-4\cdot4-3\cdot2=2$.
$U\x=\y$: $x_3=1$, $x_2=2-1=1$, $x_1=\frac{4-1-1}{2}=1$. Logo $\x=(1,1,1)$ e
$\det A=2\cdot1\cdot2=4$.
\end{exemplo}

\begin{exemplo}[: fatoração com troca ($PA=LU$)]
Para $A=\begin{pmatrix}0&1&2\\1&1&1\\2&1&3\end{pmatrix}$ e $\bb=(8,6,13)$, o pivô
$a_{11}=0$ obriga a trocar $L_1\leftrightarrow L_2$. Então
\[
P=\begin{pmatrix}0&1&0\\1&0&0\\0&0&1\end{pmatrix},\quad
L=\begin{pmatrix}1&0&0\\0&1&0\\2&-1&1\end{pmatrix},\quad
U=\begin{pmatrix}1&1&1\\0&1&2\\0&0&3\end{pmatrix},\quad
P\bb=\begin{pmatrix}6\\8\\13\end{pmatrix}.
\]
$L\y=P\bb$ dá $\y=(6,8,9)$ e $U\x=\y$ dá $\x=(1,2,3)$. Como houve uma troca,
$\det A=-(1\cdot1\cdot3)=-3$.
\end{exemplo}

\begin{atividade}{fatoração $PA=LU$, matriz $3\times3$}{\arqLUtres}
Em cada etapa, confira o pivô, troque linhas se ele for zero e informe os multiplicadores; eles
aparecem em verde na matriz $L$. Ao final, a caixa ``Resolver $A\x=\bb$ usando LU'' mostra
$P\bb$, $\y$ e $\x$.
\begin{enumerate}[nosep]
  \item Clique em \emph{Exemplo: sem troca} e refaça o primeiro exemplo acima.
  \item Clique em \emph{Exemplo: troca na etapa 1} e refaça o segundo exemplo. Observe como a
        troca de linhas também troca os multiplicadores já guardados em $L$.
  \item Clique em \emph{Exemplo: troca só na etapa 2}. Quem é $P$ neste caso?
  \item Clique em \emph{Exemplo: A singular}. Qual elemento de $U$ se anula? O que isso diz
        sobre $\det A$ e sobre o sistema?
\end{enumerate}
\end{atividade}

\begin{atividade}{fatoração $PA=LU$, matriz $4\times4$}{\arqLUquatro}
\begin{enumerate}[nosep]
  \item No \emph{Exemplo: sem troca}, obtenha
  $L=\left(\begin{smallmatrix}1&0&0&0\\2&1&0&0\\4&3&1&0\\3&4&1&1\end{smallmatrix}\right)$,
  $U=\left(\begin{smallmatrix}2&1&1&0\\0&1&1&1\\0&0&2&2\\0&0&0&2\end{smallmatrix}\right)$,
  $\y=(4,3,4,2)$ e $\x=(1,1,1,1)$. Calcule $\det A$.
  \item No \emph{Exemplo: trocas nas etapas 1 e 3} (é a matriz do exemplo de eliminação de
  Gauss), compare a matriz $U$ com a matriz $B_3$ obtida no aplicativo de Gauss.
  \item Usando a fatoração do item 1, resolva à mão $A\x=\bb$ com outro vetor, por exemplo
  $\bb=(3,8,24,24)$, sem refazer a eliminação. (Resposta: $\y=(3,2,6,1)$ e
  $\x=(\frac34,-1,\frac52,\frac12)$.)
\end{enumerate}
\end{atividade}

\paragraph{Por que fatorar?} Para um único sistema, a fatoração $LU$ mais as duas substituições
custa exatamente o mesmo que a eliminação de Gauss. A vantagem aparece quando a \emph{mesma}
matriz $A$ é usada com vários vetores $\bb$ (por exemplo, em processos iterativos ou em
problemas com várias cargas): a parte cara, de ordem $n^3$, é feita uma única vez, e cada novo
sistema custa só $2n^2$ operações.

% ==================================================================
\section{Custo computacional dos métodos diretos}

Contamos as operações de ponto flutuante: multiplicações e divisões (M/D) e adições e
subtrações (A/S). Para $n$ grande, só o termo de maior grau importa.

\begin{center}
\begin{tabular}{>{\raggedright\arraybackslash}p{5.2cm}>{\raggedright\arraybackslash}p{3.6cm}>{\raggedright\arraybackslash}p{3.6cm}c}
\toprule
Método & M/D & A/S & Ordem\\
\midrule
Eliminação de Gauss (com retrossubstituição) & $\frac{n^3}{3}+n^2-\frac{n}{3}$ & $\frac{n^3}{3}+\frac{n^2}{2}-\frac{5n}{6}$ & $\frac23n^3$\\
$LU$: fatoração & $\frac{n^3-n}{3}$ & $\frac{n^3}{3}-\frac{n^2}{2}+\frac{n}{6}$ & $\frac23n^3$\\
$LU$: resolução ($L\y=P\bb$, $U\x=\y$) & $n^2$ & $n^2-n$ & $2n^2$\\
Matriz inversa e produto $A^{-1}\bb$ & $n^3+n^2$ & $n^3-n^2$ & $2n^3$\\
Cramer, determinantes por eliminação & \multicolumn{2}{l}{$(n+1)$ determinantes de custo $\approx\frac23n^3$} & $\frac23n^4$\\
Cramer, determinantes por cofatores & \multicolumn{2}{l}{$M_n=nM_{n-1}+n$, $S_n=nS_{n-1}+n-1$} & $e\,(n+1)!$\\
\bottomrule
\end{tabular}
\end{center}

\begin{center}
\begin{tabular}{lrrrr}
\toprule
Método & $n=4$ & $n=10$ & $n=20$ & $n=100$\\
\midrule
Eliminação de Gauss / $LU$ & 62 & 805 & 5\,910 & 681\,550\\
Matriz inversa & 128 & 2\,000 & 16\,000 & 2\,000\,000\\
Cramer (eliminação) & 189 & 6\,874 & 108\,149 & $6{,}7\times10^{7}$\\
Cramer (cofatores) & 319 & $1{,}1\times10^{8}$ & $1{,}4\times10^{20}$ & $2{,}6\times10^{160}$\\
\bottomrule
\end{tabular}
\end{center}

Num computador de $10^{12}$ operações por segundo, com $n=20$, a eliminação de Gauss leva
alguns nanossegundos, enquanto a regra de Cramer por cofatores levaria cerca de quatro anos e
meio. No tamanho dos aplicativos ($n=4$), Cramer por cofatores já custa cerca de cinco vezes
a eliminação.

\paragraph{Conclusões.}
\begin{itemize}[nosep]
  \item Em ordem de custo, para um sistema: Gauss $\approx$ $LU$ $<$ inversa $<$ Cramer.
  \item Nunca se calcula $A^{-1}$ para resolver $A\x=\bb$: custa cerca de três vezes mais e é
        menos estável; mesmo com muitos vetores $\bb$, a $LU$ é melhor.
  \item A regra de Cramer tem grande valor teórico, mas não é usada no cálculo numérico.
  \item Aproveitar a estrutura da matriz muda tudo: para matrizes tridiagonais, o algoritmo de
        Thomas custa $\approx8n$ operações, em vez de $\frac23n^3$.
  \item Nenhum método direto corrige um problema \emph{mal condicionado}: o erro relativo é da
        ordem de $\varepsilon\,\kappa(A)$, com $\varepsilon\approx10^{-16}$ e $\kappa(A)$ o número
        de condição. Para a matriz de Hilbert $13\times13$, $\kappa\approx5\times10^{17}$ e a
        solução calculada não tem nenhum algarismo correto, qualquer que seja o método.
\end{itemize}

\begin{leitura}{custo computacional (3 páginas)}{\pdfCusto}
Dedução das contagens de operações da eliminação de Gauss, da fatoração $LU$, da matriz
inversa e da regra de Cramer (por cofatores e por eliminação), com a tabela comparativa e o
caso $n=4$ dos aplicativos.
\end{leitura}

\begin{leitura}{comparação computacional dos métodos diretos (10 páginas)}{\pdfComp}
Comparação ampliada (Gauss com e sem pivoteamento, Gauss-Jordan, $LU$, Cholesky, $LDL^T$,
$QR$, inversa, Cramer e Thomas): contagem exata de operações, memória, tempos medidos em
Python/NumPy e precisão numérica em matrizes bem e mal condicionadas (matriz de Hilbert).
\end{leitura}

% ==================================================================
\section{Métodos iterativos: a ideia}

Para sistemas muito grandes e esparsos (com a maioria dos coeficientes nulos), os métodos
diretos podem ser caros e destruir a esparsidade (a eliminação preenche posições que eram
zero). Os métodos iterativos partem de um \emph{chute inicial} $\x^{(0)}$ e geram uma
sequência $\x^{(1)},\x^{(2)},\dots$ que, sob certas condições, converge para a solução.

A ideia é reescrever $A\x=\bb$ na forma equivalente $\x=C\x+\mathbf g$ e iterar
\[
\x^{(k+1)}=C\x^{(k)}+\mathbf g,\qquad k=0,1,2,\dots
\]
Cada iteração custa da ordem de $2n^2$ operações para matriz cheia (bem menos para matriz
esparsa). Se forem necessárias $k$ iterações, o custo total $\approx2kn^2$ é menor que o de
Gauss ($\frac23n^3$) sempre que $k<n/3$.

\paragraph{Critério de parada.} Interrompemos as iterações quando
\[
d^{(k)}=\max_{1\le i\le n}\bigl|x_i^{(k)}-x_i^{(k-1)}\bigr|<\varepsilon
\]
(tolerância $\varepsilon$ dada) ou quando se atinge um número máximo $N$ de iterações. Também
é comum usar o erro relativo $d^{(k)}/\max_i|x_i^{(k)}|$. Os aplicativos usam $d^{(k)}<\varepsilon$
com $\varepsilon=0{,}001$.

% ==================================================================
\section{Método de Gauss-Jacobi}

Supondo $a_{ii}\neq0$, isolamos $x_i$ na $i$-ésima equação e usamos, no lado direito, os
valores da iteração anterior:
\[
x_i^{(k+1)}=\frac{1}{a_{ii}}\Bigl(b_i-\sum_{j\neq i}a_{ij}\,x_j^{(k)}\Bigr),\qquad i=1,\dots,n.
\]
Para $n=3$:
\[
\begin{cases}
x_1^{(k+1)}=\dfrac{b_1-a_{12}x_2^{(k)}-a_{13}x_3^{(k)}}{a_{11}}\\[8pt]
x_2^{(k+1)}=\dfrac{b_2-a_{21}x_1^{(k)}-a_{23}x_3^{(k)}}{a_{22}}\\[8pt]
x_3^{(k+1)}=\dfrac{b_3-a_{31}x_1^{(k)}-a_{32}x_2^{(k)}}{a_{33}}
\end{cases}
\]
Se algum $a_{ii}=0$, é preciso reordenar as equações.

\begin{teorema}[Critério das linhas]
Seja
\[
\alpha_i=\frac{1}{|a_{ii}|}\sum_{j\neq i}|a_{ij}|,\qquad i=1,\dots,n.
\]
Se $\displaystyle\max_i\alpha_i<1$ (isto é, $A$ é \emph{estritamente diagonalmente dominante}
por linhas), o método de Gauss-Jacobi converge para a solução, qualquer que seja $\x^{(0)}$.
\end{teorema}

O critério é apenas \emph{suficiente}: o método pode convergir mesmo quando ele falha.
Trocar a ordem das equações pode fazer o critério passar a valer.

\begin{exemplo}[: Gauss-Jacobi (exemplo inicial do aplicativo)]
Resolva com $\x^{(0)}=(0,0,0)$ e $\varepsilon=0{,}001$:
\[
\begin{cases}
5x_1+2x_2+x_3=4\\
-x_1-6x_2-3x_3=16\\
-3x_1+2x_2+7x_3=-12
\end{cases}
\qquad
\alpha_1=\tfrac{2+1}{5}=\tfrac35,\quad
\alpha_2=\tfrac{1+3}{6}=\tfrac23,\quad
\alpha_3=\tfrac{3+2}{7}=\tfrac57.
\]
Como $\max\alpha_i=\frac57<1$, há convergência. As iterações são
\[
x_1^{(k+1)}=\frac{4-2x_2^{(k)}-x_3^{(k)}}{5},\quad
x_2^{(k+1)}=\frac{16+x_1^{(k)}+3x_3^{(k)}}{-6},\quad
x_3^{(k+1)}=\frac{-12+3x_1^{(k)}-2x_2^{(k)}}{7}.
\]
\begin{center}
\begin{tabular}{rrrrr}
\toprule
$k$ & $x_1^{(k)}$ & $x_2^{(k)}$ & $x_3^{(k)}$ & $d^{(k)}$\\
\midrule
0 & 0 & 0 & 0 & --\\
1 & 0{,}8000 & $-2{,}6667$ & $-1{,}7143$ & 2{,}6667\\
2 & 2{,}2095 & $-1{,}9429$ & $-0{,}6095$ & 1{,}4095\\
3 & 1{,}6990 & $-2{,}7302$ & $-0{,}2122$ & 0{,}7873\\
4 & 1{,}9345 & $-2{,}8437$ & $-0{,}2061$ & 0{,}2355\\
5 & 1{,}9787 & $-2{,}8860$ & $-0{,}0727$ & 0{,}1334\\
6 & 1{,}9690 & $-2{,}9601$ & $-0{,}0417$ & 0{,}0740\\
7 & 1{,}9924 & $-2{,}9740$ & $-0{,}0247$ & 0{,}0234\\
8 & 1{,}9945 & $-2{,}9864$ & $-0{,}0107$ & 0{,}0140\\
9 & 1{,}9967 & $-2{,}9937$ & $-0{,}0062$ & 0{,}0074\\
10 & 1{,}9987 & $-2{,}9963$ & $-0{,}0032$ & 0{,}0030\\
11 & 1{,}9992 & $-2{,}9982$ & $-0{,}0016$ & 0{,}0019\\
\textbf{12} & \textbf{1{,}9996} & $\mathbf{-2{,}9991}$ & $\mathbf{-0{,}0009}$ & \textbf{0{,}0009}\\
\bottomrule
\end{tabular}
\end{center}
A parada ocorre em $k=12$. A solução exata é $\x=(2,-3,0)$.
\end{exemplo}

\begin{atividade}{método de Gauss-Jacobi}{\arqJacobi}
Nas etapas 1 e 2, calcule $\x^{(1)}$ e $\x^{(2)}$ (as fórmulas mostram em azul os valores da
iteração anterior). Depois veja o critério das linhas e marque ``Continuar as iterações no
GeoGebra'' para ver a tabela completa. Use as barras de rolagem para percorrer a tela.
\begin{enumerate}[nosep]
  \item Confira as duas primeiras linhas da tabela do exemplo.
  \item Com $N=10$, o aplicativo avisa que o erro ainda é $\ge\varepsilon$. Aumente $N$ para
        $12$ e confira a parada em $k=12$.
  \item Mude o chute inicial para $(1,1,1)$ e depois para $(10,-10,10)$. O número de iterações
        muda? A solução muda? Por quê?
  \item Troque a ordem das duas primeiras equações e observe o critério das linhas e a tabela.
  \item Diminua $\varepsilon$ para $10^{-6}$ e observe quantas iterações a mais são necessárias.
\end{enumerate}
\end{atividade}

% ==================================================================
\section{Método de Gauss-Seidel}

O método de Gauss-Seidel é uma modificação simples: ao calcular $x_i^{(k+1)}$, usamos os
valores \emph{já atualizados} $x_1^{(k+1)},\dots,x_{i-1}^{(k+1)}$:
\[
x_i^{(k+1)}=\frac{1}{a_{ii}}\Bigl(b_i-\sum_{j<i}a_{ij}\,x_j^{(k+1)}-\sum_{j>i}a_{ij}\,x_j^{(k)}\Bigr),
\qquad i=1,\dots,n.
\]
Para $n=3$:
\[
\begin{cases}
x_1^{(k+1)}=\dfrac{b_1-a_{12}x_2^{(k)}-a_{13}x_3^{(k)}}{a_{11}}\\[8pt]
x_2^{(k+1)}=\dfrac{b_2-a_{21}\textcolor{vermelho}{x_1^{(k+1)}}-a_{23}x_3^{(k)}}{a_{22}}\\[8pt]
x_3^{(k+1)}=\dfrac{b_3-a_{31}\textcolor{vermelho}{x_1^{(k+1)}}-a_{32}\textcolor{vermelho}{x_2^{(k+1)}}}{a_{33}}
\end{cases}
\]
Além de, em geral, convergir mais depressa, o método só precisa guardar um vetor.

\begin{teorema}[Critério de Sassenfeld]
Sejam
\[
\beta_1=\frac{1}{|a_{11}|}\sum_{j=2}^{n}|a_{1j}|,\qquad
\beta_i=\frac{1}{|a_{ii}|}\Bigl(\sum_{j<i}|a_{ij}|\,\beta_j+\sum_{j>i}|a_{ij}|\Bigr),\quad i=2,\dots,n.
\]
Se $\displaystyle\max_i\beta_i<1$, o método de Gauss-Seidel converge para a solução,
qualquer que seja $\x^{(0)}$. Se o critério das linhas vale, o de Sassenfeld também vale.
\end{teorema}

\begin{exemplo}[: Gauss-Seidel (exemplo inicial do aplicativo)]
Resolva com $\x^{(0)}=(0,0,0)$ e $\varepsilon=0{,}001$:
\[
\begin{cases}
6x_1-2x_2=14\\
3x_1+9x_2-2x_3=31\\
x_1+2x_2-6x_3=19
\end{cases}
\qquad
\beta_1=\tfrac26=\tfrac13,\quad
\beta_2=\tfrac{3\cdot\frac13+2}{9}=\tfrac13,\quad
\beta_3=\tfrac{1\cdot\frac13+2\cdot\frac13}{6}=\tfrac16.
\]
Como $\max\beta_i=\frac13<1$, há convergência.
\begin{center}
\begin{tabular}{rrrrr@{\qquad}r}
\toprule
$k$ & $x_1^{(k)}$ & $x_2^{(k)}$ & $x_3^{(k)}$ & $d^{(k)}$ & $d^{(k)}$ com Jacobi\\
\midrule
0 & 0 & 0 & 0 & -- & --\\
1 & 2{,}3333 & 2{,}6667 & $-1{,}8889$ & 2{,}6667 & 3{,}4444\\
2 & 3{,}2222 & 1{,}9506 & $-1{,}9794$ & 0{,}8889 & 1{,}5370\\
3 & 2{,}9835 & 2{,}0101 & $-1{,}9994$ & 0{,}2387 & 0{,}4938\\
4 & 3{,}0034 & 1{,}9990 & $-1{,}9998$ & 0{,}0198 & 0{,}0974\\
5 & 2{,}9997 & 2{,}0002 & $-2{,}0000$ & 0{,}0037 & 0{,}0325\\
\textbf{6} & \textbf{3{,}0001} & \textbf{2{,}0000} & $\mathbf{-2{,}0000}$ & \textbf{0{,}0004} & 0{,}0056\\
\bottomrule
\end{tabular}
\end{center}
Gauss-Seidel para em $k=6$; Gauss-Jacobi, no mesmo sistema, precisaria de $k=8$. A solução
exata é $\x=(3,2,-2)$.
\end{exemplo}

\begin{atividade}{método de Gauss-Seidel}{\arqSeidel}
Nas etapas 1 e 2, calcule $x_1$, $x_2$, $x_3$ \emph{nessa ordem}, usando sempre os valores já
atualizados (em vermelho nas fórmulas). Depois veja o critério de Sassenfeld, a tabela
completa e a comparação com Gauss-Jacobi.
\begin{enumerate}[nosep]
  \item Confira as duas primeiras linhas da tabela do exemplo.
  \item O aplicativo informa em que iteração Gauss-Jacobi pararia. Confira que é $k=8$.
  \item Digite no aplicativo de Gauss-Seidel o sistema do exemplo de Gauss-Jacobi. Em quantas
        iterações Gauss-Seidel para? (Resposta esperada: $k=8$, contra $k=12$ de Jacobi.)
  \item Gere sistemas com \emph{Novo sistema aleatório} e compare $\max\alpha_i$ (Jacobi) e
        $\max\beta_i$ (Sassenfeld). Qual costuma ser menor?
\end{enumerate}
\end{atividade}

% ==================================================================
\section{Métodos diretos ou iterativos?}

\begin{center}
\begin{tabular}{>{\raggedright\arraybackslash}p{4.2cm}>{\raggedright\arraybackslash}p{5.2cm}>{\raggedright\arraybackslash}p{5.2cm}}
\toprule
 & Diretos (Gauss, $LU$) & Iterativos (Jacobi, Seidel)\\
\midrule
Resultado & exato em aritmética exata (a menos de arredondamentos) & aproximação com tolerância $\varepsilon$\\
Custo & $\approx\frac23n^3$, previsível & $\approx2n^2$ por iteração (menos se esparsa)\\
Convergência & sempre termina (se $\det A\neq0$) & depende de $A$ (critérios das linhas e de Sassenfeld)\\
Memória & preenche a matriz (perde a esparsidade) & preserva a esparsidade\\
Erros de arredondamento & acumulam ao longo da eliminação & cada iteração ``corrige'' a anterior\\
Indicados para & sistemas pequenos e médios, matrizes cheias, vários $\bb$ ($LU$) & sistemas grandes e esparsos, diagonalmente dominantes\\
\bottomrule
\end{tabular}
\end{center}

% ==================================================================
\section{Exercícios}

\begin{enumerate}[label=\textbf{\arabic*.},leftmargin=*]
\item Resolva por eliminação de Gauss:
  $\begin{cases}x_1+x_2+x_3=6\\2x_1+x_2-x_3=1\\x_1-x_2+2x_3=5\end{cases}$.

\item No aplicativo de eliminação de Gauss, resolva os exemplos \emph{troca + SPI} e
  \emph{sistema impossível}. Para o SPI, escreva a solução geral.

\item Obtenha a fatoração $PA=LU$ de
  $A=\begin{pmatrix}1&2&1\\2&4&3\\1&3&2\end{pmatrix}$ e use-a para resolver $A\x=(4,9,6)$.
  Confira no aplicativo $3\times3$ (\emph{Exemplo: troca só na etapa 2}).

\item Use a fatoração do primeiro exemplo de $LU$ ($A=\left(\begin{smallmatrix}2&1&1\\4&3&3\\8&7&9\end{smallmatrix}\right)$)
  para resolver $A\x=(1,2,6)$ sem refazer a eliminação.

\item Calcule $\det A$ para as matrizes dos dois exemplos de $LU$ usando $U$ e o número de
  trocas, e confira pela regra de Sarrus.

\item Um computador faz $10^{9}$ operações por segundo. Estime o tempo para resolver um sistema
  com $n=1000$ por eliminação de Gauss e pela matriz inversa. Quantos sistemas com a mesma
  matriz podem ser resolvidos pela $LU$ no tempo de uma única inversão?

\item Para o sistema
  $\begin{cases}10x_1+2x_2+x_3=7\\x_1+5x_2+x_3=-8\\2x_1+3x_2+10x_3=6\end{cases}$,
  com $\x^{(0)}=(0,0,0)$ e $\varepsilon=0{,}001$:
  \begin{enumerate}[label=(\alph*),nosep]
  \item verifique os critérios das linhas e de Sassenfeld;
  \item faça duas iterações de Gauss-Jacobi e duas de Gauss-Seidel à mão;
  \item use os aplicativos para obter o número de iterações de cada método até a parada.
  \end{enumerate}

\item O sistema $\begin{cases}2x_1+10x_2=12\\8x_1+x_2=9\end{cases}$ não satisfaz o critério das
  linhas. Reordene as equações para que satisfaça e resolva por Gauss-Seidel com
  $\x^{(0)}=(0,0)$ e duas iterações.

\item \textbf{(GeoGebra)} No aplicativo de Gauss-Jacobi, construa um sistema $3\times3$ em que
  o critério das linhas falhe mas o método ainda convirja (sugestão: um único $\alpha_i$
  ligeiramente maior que 1). O que isso mostra sobre o critério?
\end{enumerate}

\subsection*{Respostas}
\begin{enumerate}[label=\textbf{\arabic*.},leftmargin=*,nosep]
\item $m_{21}=2$, $m_{31}=1$, $m_{32}=2$; $U$ com diagonal $1,-1,7$; $\x=(1,2,3)$.
\item SPI: após a troca $L_2\leftrightarrow L_3$ na etapa 2, a última linha fica
  $(0\ 0\ 0\ 0\,|\,0)$. Restam $x_1+2x_2-x_3+x_4=3$, $-x_2+2x_3=1$ e $3x_3-3x_4=0$; com
  $x_4=t$ livre, $\x=(5-4t,\;2t-1,\;t,\;t)$, $t\in\mathbb R$. SI: aparece a linha
  $(0\ 0\ 0\ 0\,|\,-2)$, que vem de $L_2-L_1$.
\item $P$ troca as linhas 2 e 3; $L=\left(\begin{smallmatrix}1&0&0\\1&1&0\\2&0&1\end{smallmatrix}\right)$,
  $U=\left(\begin{smallmatrix}1&2&1\\0&1&1\\0&0&1\end{smallmatrix}\right)$, $P\bb=(4,6,9)$,
  $\y=(4,2,1)$, $\x=(1,1,1)$.
\item $\y=(1,0,2)$ e $\x=(\tfrac12,-1,1)$.
\item $4$ e $-3$.
\item Gauss: $\approx\frac23\cdot10^9$ operações, cerca de $0{,}67$ s; inversa: $\approx2\cdot10^9$, cerca de $2$ s.
  Com a $LU$ já feita, cada sistema custa $2\cdot10^{6}$ operações ($2$ ms): no tempo da
  inversão, fatora-se e resolvem-se cerca de $660$ sistemas.
\item (a) $\alpha=(0{,}3;\,0{,}4;\,0{,}5)$; $\beta=(0{,}3;\,0{,}26;\,0{,}138)$. (b) Jacobi:
  $\x^{(1)}=(0{,}7;\,-1{,}6;\,0{,}6)$, $\x^{(2)}=(0{,}96;\,-1{,}86;\,0{,}94)$; Seidel:
  $\x^{(1)}=(0{,}7;\,-1{,}74;\,0{,}982)$, $\x^{(2)}=(0{,}9498;\,-1{,}9864;\,1{,}0059)$.
  (c) Jacobi para em $k=8$ e Seidel em $k=5$; a solução é $\x=(1,-2,1)$.
\item Trocando as equações: $8x_1+x_2=9$, $2x_1+10x_2=12$; $\alpha=(\tfrac18,\tfrac15)$.
  Seidel: $\x^{(1)}=(1{,}125;\,0{,}975)$, $\x^{(2)}=(1{,}0031;\,0{,}9994)$; solução $(1,1)$.
\item O critério das linhas é apenas suficiente: sua falha não implica divergência.
\end{enumerate}

% Capítulo gerado a partir de interpolacao.tex (arquivos usados neste capítulo)
\def\arqGauss{EliminacaoGauss_Interpolacao_4pontos_automatico.ggb}
\def\arqLagrange{Interpolacao_Lagrange_4pontos_automatico.ggb}
\def\arqNewton{Interpolacao_Newton_4pontos_automatico.ggb}
\chapter{Interpolação polinomial}\label{cap:interp}

\begin{tcolorbox}[colback=gray!4,colframe=gray!50,title={Como usar os aplicativos deste capítulo},fonttitle=\bfseries]
Ao longo do capítulo, os quadros laranja \emph{Atividade no GeoGebra} trazem um link
\textbf{[abrir o aplicativo]} que abre o arquivo correspondente no GeoGebra Classic:
\begin{itemize}[nosep]
  \item \ggbabrir{\arqGauss}{Interpolação via sistema linear e eliminação de Gauss};
  \item \ggbabrir{\arqLagrange}{Interpolação de Lagrange};
  \item \ggbabrir{\arqNewton}{Interpolação de Newton (diferenças divididas)}.
\end{itemize}
Para que os links funcionem, mantenha este PDF e a pasta \texttt{ggb} juntos, na mesma
pasta, e abra o PDF num leitor que permita abrir arquivos externos (Adobe Acrobat Reader,
Okular, SumatraPDF, Foxit). Os leitores de PDF embutidos nos navegadores costumam bloquear
esse tipo de link; nesse caso, abra os arquivos da pasta \texttt{ggb} diretamente no GeoGebra.

Os três aplicativos começam com os mesmos quatro pontos
$A(-4,2)$, $B(-2,-2)$, $C(4,3)$ e $D(8,0)$, que são os pontos do exemplo
desenvolvido no texto. Ao arrastar qualquer ponto, todas as contas são refeitas.
Nos aplicativos, cada janela da direita tem barras de rolagem (vertical e horizontal)
para percorrer as contas.
\end{tcolorbox}

% ==================================================================
\section{O problema da interpolação}

Em muitas situações conhecemos uma função apenas em alguns pontos: uma tabela de
temperaturas medidas de hora em hora, valores tabelados de uma função especial, resultados
de um experimento ou de uma simulação cara. Interpolar é construir uma função simples que
\emph{passe exatamente} pelos pontos conhecidos e usá-la para estimar valores intermediários,
derivar, integrar ou simplesmente desenhar a curva.

\begin{definicao}
Dados $n$ pontos $(x_1,y_1),\,(x_2,y_2),\,\dots,\,(x_n,y_n)$ com abscissas
$x_1,\dots,x_n$ \emph{distintas} (os \emph{nós}), um \emph{polinômio interpolador} desses
pontos é um polinômio $p$ de grau menor ou igual a $n-1$ tal que
\[
p(x_i)=y_i,\qquad i=1,2,\dots,n.
\]
Quando $y_i=f(x_i)$ para alguma função $f$, dizemos que $p$ \emph{interpola $f$ nos nós}
$x_1,\dots,x_n$.
\end{definicao}

Usamos polinômios porque são fáceis de avaliar (só somas e produtos), de derivar e de
integrar, e porque, como veremos, o problema sempre tem uma única solução. Neste capítulo
estudamos três maneiras de obter esse mesmo polinômio:
\begin{enumerate}[nosep]
  \item resolvendo um \textbf{sistema linear} (matriz de Vandermonde), por eliminação de Gauss;
  \item pela \textbf{forma de Lagrange};
  \item pela \textbf{forma de Newton}, com diferenças divididas.
\end{enumerate}

\begin{figure}[h]
\centering
\begin{tikzpicture}
\begin{axis}[width=11cm,height=6.5cm,axis lines=middle,xmin=-6,xmax=10,ymin=-4,ymax=6,
  xtick={-4,-2,...,8},ytick={-2,0,...,4},grid=both,grid style={gray!20},
  xlabel={$x$},ylabel={$y$}]
\addplot[domain=-5.5:9.4,samples=150,very thick,vermelho]
  {-9/5 + 97/120*x + 43/160*x^2 - 41/960*x^3};
\addplot[only marks,mark=*,mark size=2.5pt] coordinates {(-4,2)(-2,-2)(4,3)(8,0)};
\node[above left] at (axis cs:-4,2) {$A$};
\node[below] at (axis cs:-2,-2) {$B$};
\node[above] at (axis cs:4,3.1) {$C$};
\node[above right] at (axis cs:8,0) {$D$};
\end{axis}
\end{tikzpicture}
\caption{O polinômio de grau 3 que passa pelos pontos $A$, $B$, $C$ e $D$ do exemplo
do capítulo.}
\label{fig:exemplo}
\end{figure}

% ==================================================================
\section{Existência e unicidade}

Escrevendo $p(x)=a_0+a_1x+\dots+a_{n-1}x^{n-1}$, as condições $p(x_i)=y_i$ formam um
sistema linear de $n$ equações nas $n$ incógnitas $a_0,\dots,a_{n-1}$:
\[
\begin{cases}
a_0+a_1x_1+a_2x_1^2+\dots+a_{n-1}x_1^{\,n-1}=y_1\\
a_0+a_1x_2+a_2x_2^2+\dots+a_{n-1}x_2^{\,n-1}=y_2\\
\qquad\vdots\\
a_0+a_1x_n+a_2x_n^2+\dots+a_{n-1}x_n^{\,n-1}=y_n
\end{cases}
\]
ou, na forma matricial,
\[
\underbrace{\begin{pmatrix}
1 & x_1 & x_1^2 & \cdots & x_1^{\,n-1}\\
1 & x_2 & x_2^2 & \cdots & x_2^{\,n-1}\\
\vdots & \vdots & \vdots & & \vdots\\
1 & x_n & x_n^2 & \cdots & x_n^{\,n-1}
\end{pmatrix}}_{V}
\begin{pmatrix}a_0\\ a_1\\ \vdots\\ a_{n-1}\end{pmatrix}
=
\begin{pmatrix}y_1\\ y_2\\ \vdots\\ y_n\end{pmatrix}.
\]
A matriz $V$ é a \emph{matriz de Vandermonde} dos nós.

\begin{teorema}[Existência e unicidade]\label{teo:unicidade}
Se os nós $x_1,\dots,x_n$ são distintos, existe um único polinômio $p$ de grau
menor ou igual a $n-1$ tal que $p(x_i)=y_i$ para $i=1,\dots,n$.
\end{teorema}

\begin{proof}
\emph{Unicidade.} Se $p$ e $q$ são dois polinômios interpoladores, então $r=p-q$ tem grau
no máximo $n-1$ e se anula nos $n$ pontos distintos $x_1,\dots,x_n$. Um polinômio não nulo
de grau no máximo $n-1$ tem no máximo $n-1$ raízes; logo $r\equiv 0$ e $p=q$.

\emph{Existência.} A unicidade diz que o sistema homogêneo $V\mathbf a=\mathbf 0$ só tem a
solução nula (ela corresponde ao polinômio que interpola os valores $y_i=0$). Então $V$ é
invertível e $V\mathbf a=\mathbf y$ tem solução para qualquer $\mathbf y$.
\end{proof}

\begin{observacao}
Pode-se mostrar que $\det V=\displaystyle\prod_{1\le i<j\le n}(x_j-x_i)$, que é diferente
de zero exatamente quando os nós são distintos (Exercício~\ref{ex:vandermonde}).
Se dois nós coincidem e os valores $y$ são diferentes, nenhuma função $y=p(x)$ passa pelos
dois pontos. Nos três aplicativos, colocar dois pontos na mesma vertical faz aparecer um aviso.
\end{observacao}

Como o polinômio interpolador é único, os três métodos deste capítulo produzem
\emph{o mesmo polinômio}, apenas escrito de formas diferentes e calculado com custos
diferentes.

% ==================================================================
\section{Método 1: sistema linear e eliminação de Gauss}

A forma mais direta é montar o sistema de Vandermonde e resolvê-lo por eliminação de Gauss
seguida de retrossubstituição, como em qualquer sistema linear.

\begin{exemplo}[: os pontos $A$, $B$, $C$, $D$ pelo sistema linear]
Vamos interpolar $A(-4,2)$, $B(-2,-2)$, $C(4,3)$, $D(8,0)$ por
$p(x)=a_0+a_1x+a_2x^2+a_3x^3$. Cada linha da matriz aumentada é
$\left(\begin{array}{cccc|c}1 & x_i & x_i^2 & x_i^3 & y_i\end{array}\right)$:
\[
B=\left(\begin{array}{rrrr|r}
1 & -4 & 16 & -64 & 2\\
1 & -2 & 4 & -8 & -2\\
1 & 4 & 16 & 64 & 3\\
1 & 8 & 64 & 512 & 0
\end{array}\right).
\]
\textbf{Etapa 1} (pivô $a_{11}=1$): $m_{21}=m_{31}=m_{41}=1$, e
$L_i\leftarrow L_i-m_{i1}L_1$ para $i=2,3,4$:
\[
B_1=\left(\begin{array}{rrrr|r}
1 & -4 & 16 & -64 & 2\\
0 & 2 & -12 & 56 & -4\\
0 & 8 & 0 & 128 & 1\\
0 & 12 & 48 & 576 & -2
\end{array}\right).
\]
\textbf{Etapa 2} (pivô $2$): $m_{32}=\frac{8}{2}=4$, $m_{42}=\frac{12}{2}=6$:
\[
B_2=\left(\begin{array}{rrrr|r}
1 & -4 & 16 & -64 & 2\\
0 & 2 & -12 & 56 & -4\\
0 & 0 & 48 & -96 & 17\\
0 & 0 & 120 & 240 & 22
\end{array}\right).
\]
\textbf{Etapa 3} (pivô $48$): $m_{43}=\frac{120}{48}=\frac52$:
\[
B_3=\left(\begin{array}{rrrr|r}
1 & -4 & 16 & -64 & 2\\
0 & 2 & -12 & 56 & -4\\
0 & 0 & 48 & -96 & 17\\
0 & 0 & 0 & 480 & -\frac{41}{2}
\end{array}\right).
\]
\textbf{Retrossubstituição:}
\begin{align*}
a_3&=\frac{-41/2}{480}=-\frac{41}{960}, &
a_2&=\frac{17+96a_3}{48}=\frac{43}{160},\\
a_1&=\frac{-4+12a_2-56a_3}{2}=\frac{97}{120}, &
a_0&=2+4a_1-16a_2+64a_3=-\frac95 .
\end{align*}
Portanto
\[
\boxed{\,p(x)=-\frac95+\frac{97}{120}\,x+\frac{43}{160}\,x^2-\frac{41}{960}\,x^3\,}
\]
(Figura~\ref{fig:exemplo}). Confira, por exemplo, que $p(4)=3$.
\end{exemplo}

\begin{observacao}
Na eliminação da matriz de Vandermonde, a primeira coluna é toda de 1, por isso
$m_{21}=m_{31}=\dots=1$ e o primeiro pivô nunca é zero. Mais ainda, o $k$-ésimo pivô
é $\prod_{j<k}(x_k-x_j)$ (no exemplo: $1$, $x_2-x_1=2$, $(x_3-x_1)(x_3-x_2)=48$ e
$(x_4-x_1)(x_4-x_2)(x_4-x_3)=480$). Assim, com nós distintos, a eliminação nunca precisa
de troca de linhas.
\end{observacao}

\begin{atividade}{eliminação de Gauss automática}{\arqGauss}
Abra o aplicativo. À esquerda ficam os pontos $A$, $B$, $C$, $D$ e a cúbica; à direita, a
matriz aumentada, os multiplicadores de cada etapa (com a conta
$m_{ik}=\text{elemento}/\text{pivô}$), as matrizes $B_1$, $B_2$, $B_3$ e a retrossubstituição.
\begin{enumerate}[nosep]
  \item Confira, no aplicativo, todos os números do exemplo acima.
  \item Arraste os pontos. Os multiplicadores $m_{21}$, $m_{31}$ e $m_{41}$ mudam? Explique usando a observação anterior.
  \item Verifique, para outras posições, que o pivô da etapa 2 é $x_B-x_A$ e o da etapa 3 é $(x_C-x_A)(x_C-x_B)$.
  \item Alinhe três dos pontos numa reta (por exemplo, todos com $y=0$). O que acontece com $a_3$? E se os quatro pontos estiverem alinhados?
  \item Coloque $B$ na mesma vertical de $A$. Por que as contas desaparecem?
\end{enumerate}
\end{atividade}

\paragraph{Vantagens e limitações.} O método é conceitualmente simples e fornece
diretamente os coeficientes na base $1,x,x^2,\dots$. Por outro lado, resolver o sistema custa
da ordem de $\tfrac23 n^3$ operações e a matriz de Vandermonde fica muito
\emph{mal condicionada} quando $n$ cresce (as colunas $x^k$ ficam quase paralelas), o que
amplifica os erros de arredondamento. Os métodos seguintes evitam o sistema linear.

% ==================================================================
\section{Método 2: forma de Lagrange}

A ideia de Lagrange é escrever $p$ como combinação de polinômios que ``ligam'' um nó
de cada vez.

\begin{definicao}
Os \emph{polinômios de Lagrange} associados aos nós distintos $x_1,\dots,x_n$ são
\[
L_i(x)=\prod_{\substack{j=1\\ j\neq i}}^{n}\frac{x-x_j}{x_i-x_j}
=\frac{(x-x_1)\cdots(x-x_{i-1})(x-x_{i+1})\cdots(x-x_n)}
{(x_i-x_1)\cdots(x_i-x_{i-1})(x_i-x_{i+1})\cdots(x_i-x_n)},
\qquad i=1,\dots,n.
\]
\end{definicao}

Cada $L_i$ tem grau $n-1$ e satisfaz
\[
L_i(x_j)=\delta_{ij}=\begin{cases}1,& j=i,\\ 0,& j\neq i,\end{cases}
\]
pois o numerador se anula em todos os nós diferentes de $x_i$ e, em $x_i$, numerador e
denominador são iguais.

\begin{teorema}[Forma de Lagrange]
O polinômio interpolador dos pontos $(x_i,y_i)$, $i=1,\dots,n$, é
\[
p(x)=\sum_{i=1}^{n}y_i\,L_i(x)=y_1L_1(x)+y_2L_2(x)+\dots+y_nL_n(x).
\]
\end{teorema}

\begin{proof}
$p$ tem grau no máximo $n-1$ e $p(x_j)=\sum_i y_i\,\delta_{ij}=y_j$. Pelo
Teorema~\ref{teo:unicidade}, é o polinômio interpolador.
\end{proof}

\begin{exemplo}[: os mesmos pontos pela forma de Lagrange]
Com $x_1=-4$, $x_2=-2$, $x_3=4$, $x_4=8$:
\begin{align*}
{\color{cL1}L_1(x)}&=\frac{(x+2)(x-4)(x-8)}{(-4+2)(-4-4)(-4-8)}=\frac{(x+2)(x-4)(x-8)}{-192}
=-\tfrac{1}{192}x^3+\tfrac{5}{96}x^2-\tfrac{1}{24}x-\tfrac13,\\[2pt]
{\color{cL2}L_2(x)}&=\frac{(x+4)(x-4)(x-8)}{(-2+4)(-2-4)(-2-8)}=\frac{(x+4)(x-4)(x-8)}{120}
=\tfrac{1}{120}x^3-\tfrac{1}{15}x^2-\tfrac{2}{15}x+\tfrac{16}{15},\\[2pt]
{\color{cL3}L_3(x)}&=\frac{(x+4)(x+2)(x-8)}{(4+4)(4+2)(4-8)}=\frac{(x+4)(x+2)(x-8)}{-192}
=-\tfrac{1}{192}x^3+\tfrac{1}{96}x^2+\tfrac{5}{24}x+\tfrac13,\\[2pt]
{\color{cL4}L_4(x)}&=\frac{(x+4)(x+2)(x-4)}{(8+4)(8+2)(8-4)}=\frac{(x+4)(x+2)(x-4)}{480}
=\tfrac{1}{480}x^3+\tfrac{1}{240}x^2-\tfrac{1}{30}x-\tfrac{1}{15}.
\end{align*}
Então
\[
p(x)=2L_1(x)-2L_2(x)+3L_3(x)+0\,L_4(x)
=-\frac95+\frac{97}{120}x+\frac{43}{160}x^2-\frac{41}{960}x^3,
\]
exatamente o polinômio obtido pelo sistema linear. Repare que $L_4$ nem precisaria ser
calculado, pois $y_4=0$.
\end{exemplo}

\begin{figure}[h]
\centering
\begin{tikzpicture}
\begin{axis}[width=11cm,height=6.5cm,axis lines=middle,xmin=-6,xmax=10,ymin=-1.5,ymax=2,
  xtick={-4,-2,...,8},ytick={-1,0,1},grid=both,grid style={gray!20},
  legend style={at={(0.02,0.98)},anchor=north west,font=\small},legend columns=4]
\addplot[domain=-5.5:9.5,samples=120,thick,cL1,dashed] {(x+2)*(x-4)*(x-8)/(-192)};
\addplot[domain=-5.5:9.5,samples=120,thick,cL2,dashed] {(x+4)*(x-4)*(x-8)/120};
\addplot[domain=-5.5:9.5,samples=120,thick,cL3,dashed] {(x+4)*(x+2)*(x-8)/(-192)};
\addplot[domain=-5.5:9.5,samples=120,thick,cL4,dashed] {(x+4)*(x+2)*(x-4)/480};
\legend{$L_1$,$L_2$,$L_3$,$L_4$}
\addplot[only marks,mark=*,mark size=1.8pt] coordinates {(-4,1)(-2,1)(4,1)(8,1)};
\addplot[only marks,mark=o,mark size=1.8pt] coordinates {(-4,0)(-2,0)(4,0)(8,0)};
\end{axis}
\end{tikzpicture}
\caption{As bases de Lagrange do exemplo: cada $L_i$ vale 1 no seu nó (pontos cheios) e 0
nos demais nós (círculos).}
\end{figure}

\begin{atividade}{forma de Lagrange automática}{\arqLagrange}
O aplicativo mostra, para cada base, a fórmula, a substituição dos nós, o denominador e a
forma expandida, além do polinômio $p$. No gráfico, as bases aparecem tracejadas nas
mesmas cores dos quadros.
\begin{enumerate}[nosep]
  \item Com a caixa ``Mostrar as bases'' marcada, confira que cada $L_i$ vale 1 em $x_i$ e 0 nos outros nós.
  \item Mova o ponto $C$ \emph{apenas na vertical}. Quais bases mudam? Qual parcela de $p$ muda? Por quê?
  \item Agora mova $C$ na horizontal. Por que, desta vez, todas as bases mudam?
  \item Coloque os quatro pontos com a mesma ordenada $y=c$. Qual é o polinômio $p$? Conclua que $L_1+L_2+L_3+L_4\equiv 1$.
  \item Compare a cúbica e os coeficientes com os do aplicativo de eliminação de Gauss.
\end{enumerate}
\end{atividade}

\paragraph{Vantagens e limitações.} A forma de Lagrange dispensa a resolução de sistemas e
é muito útil em teoria (dela saem fórmulas de derivação e integração numérica). Quando os
nós são fixos e só os valores $y_i$ mudam, basta recombinar as mesmas bases. Por outro
lado, se acrescentarmos um novo nó, \emph{todas} as bases precisam ser refeitas.

% ==================================================================
\section{Método 3: forma de Newton e diferenças divididas}

A forma de Newton constrói o polinômio aos poucos, um nó de cada vez:
\[
p(x)=c_1+c_2(x-x_1)+c_3(x-x_1)(x-x_2)+\dots+c_n(x-x_1)(x-x_2)\cdots(x-x_{n-1}).
\]
Impondo $p(x_1)=y_1$, $p(x_2)=y_2$, \dots\ sucessivamente, obtemos os coeficientes $c_k$ um a
um: o sistema é triangular. Esses coeficientes são as \emph{diferenças divididas}.

\begin{definicao}
As \emph{diferenças divididas} de $f$ nos nós $x_1,\dots,x_n$ são definidas recursivamente por
\begin{align*}
f[x_i]&=f(x_i)=y_i,\\
f[x_i,x_{i+1},\dots,x_{i+k}]&=\frac{f[x_{i+1},\dots,x_{i+k}]-f[x_i,\dots,x_{i+k-1}]}{x_{i+k}-x_i}.
\end{align*}
\end{definicao}

\begin{teorema}[Forma de Newton]
O polinômio interpolador de $f$ nos nós distintos $x_1,\dots,x_n$ é
\[
p(x)=f[x_1]+f[x_1,x_2](x-x_1)+f[x_1,x_2,x_3](x-x_1)(x-x_2)+\dots
+f[x_1,\dots,x_n]\prod_{j=1}^{n-1}(x-x_j).
\]
Além disso, se $p_{k}$ denota o polinômio que interpola os $k+1$ primeiros nós, vale
\[
p_k(x)=p_{k-1}(x)+f[x_1,\dots,x_{k+1}]\,(x-x_1)(x-x_2)\cdots(x-x_k).
\]
\end{teorema}

A segunda fórmula é a grande vantagem do método: ao acrescentar um ponto, aproveitamos
tudo o que já foi calculado e somamos apenas \emph{um termo}. As diferenças divididas são
organizadas numa tabela triangular; os coeficientes da forma de Newton ficam na primeira
linha (ou na diagonal superior, conforme a disposição).

\begin{exemplo}[: os mesmos pontos pela forma de Newton]
Para $A(-4,2)$, $B(-2,-2)$, $C(4,3)$, $D(8,0)$:

\textbf{Ordem 1:}
\[
f[x_1,x_2]=\frac{-2-2}{-2-(-4)}=-2,\quad
f[x_2,x_3]=\frac{3-(-2)}{4-(-2)}=\frac56,\quad
f[x_3,x_4]=\frac{0-3}{8-4}=-\frac34.
\]
\textbf{Ordem 2:}
\[
f[x_1,x_2,x_3]=\frac{\frac56-(-2)}{4-(-4)}=\frac{17}{48},\qquad
f[x_2,x_3,x_4]=\frac{-\frac34-\frac56}{8-(-2)}=-\frac{19}{120}.
\]
\textbf{Ordem 3:}
\[
f[x_1,x_2,x_3,x_4]=\frac{-\frac{19}{120}-\frac{17}{48}}{8-(-4)}=-\frac{41}{960}.
\]
\begin{center}
\renewcommand{\arraystretch}{1.6}
\begin{tabular}{c|c|c|c|c}
$x_i$ & $f[x_i]$ & ordem 1 & ordem 2 & ordem 3\\ \hline
$-4$ & $\color{vermelho}\boxed{2}$ & $\color{vermelho}\boxed{-2}$ & $\color{vermelho}\boxed{\frac{17}{48}}$ & $\color{vermelho}\boxed{-\frac{41}{960}}$\\
$-2$ & $-2$ & $\frac56$ & $-\frac{19}{120}$ & \\
$4$ & $3$ & $-\frac34$ & & \\
$8$ & $0$ & & &
\end{tabular}
\end{center}
Logo
\[
p(x)=2-2(x+4)+\frac{17}{48}(x+4)(x+2)-\frac{41}{960}(x+4)(x+2)(x-4),
\]
que, expandido, é novamente $p(x)=-\frac95+\frac{97}{120}x+\frac{43}{160}x^2-\frac{41}{960}x^3$.
Repare que o coeficiente de $x^3$ é a diferença dividida de maior ordem.
\end{exemplo}

\begin{figure}[h]
\centering
\begin{tikzpicture}
\begin{axis}[width=11cm,height=6.5cm,axis lines=middle,xmin=-6,xmax=10,ymin=-4,ymax=8,
  xtick={-4,-2,...,8},ytick={-2,0,...,6},grid=both,grid style={gray!20},
  legend style={at={(0.02,0.98)},anchor=north west,font=\small},legend columns=4]
\addplot[domain=-5.5:9.5,thick,cL1,dashed] {2};
\addplot[domain=-5.5:9.5,thick,cL2,dashed] {2-2*(x+4)};
\addplot[domain=-5.5:9.5,samples=120,thick,cL3,dashed] {2-2*(x+4)+17/48*(x+4)*(x+2)};
\addplot[domain=-5.5:9.4,samples=150,very thick,vermelho]
  {-9/5 + 97/120*x + 43/160*x^2 - 41/960*x^3};
\legend{$p_0$,$p_1$,$p_2$,$p_3=p$}
\addplot[only marks,mark=*,mark size=2.2pt] coordinates {(-4,2)(-2,-2)(4,3)(8,0)};
\end{axis}
\end{tikzpicture}
\caption{Construção incremental de Newton: $p_0$ passa por $A$; $p_1$ por $A,B$;
$p_2$ por $A,B,C$; $p_3$ pelos quatro pontos.}
\end{figure}

\paragraph{Avaliação eficiente (forma aninhada).} A forma de Newton pode ser avaliada como
no algoritmo de Horner:
\[
p(x)=c_1+(x-x_1)\Bigl(c_2+(x-x_2)\bigl(c_3+(x-x_3)\,c_4\bigr)\Bigr),
\]
o que exige apenas $n-1$ multiplicações para cada valor de $x$.

\begin{atividade}{diferenças divididas automáticas}{\arqNewton}
O aplicativo mostra as diferenças de cada ordem com as contas, a tabela triangular (com os
coeficientes destacados em vermelho) e o polinômio nas formas de Newton e expandida. No
gráfico, a caixa de seleção mostra os polinômios parciais $p_0$, $p_1$ e $p_2$.
\begin{enumerate}[nosep]
  \item Confira a tabela do exemplo acima.
  \item Arraste apenas o ponto $D$. Quais entradas da tabela mudam? Os polinômios $p_0$, $p_1$, $p_2$ mudam? Relacione com a fórmula $p_3=p_2+f[x_1,\dots,x_4](x-x_1)(x-x_2)(x-x_3)$.
  \item Arraste agora o ponto $A$. Por que, desta vez, toda a primeira linha da tabela muda?
  \item Posicione $D$ sobre a parábola $p_2$. Quanto vale $f[x_1,x_2,x_3,x_4]$? Interprete.
  \item Compare com os aplicativos de Gauss e de Lagrange e verifique que os três chegam ao mesmo polinômio.
\end{enumerate}
\end{atividade}

% ==================================================================
\section{Comparação dos três métodos}

\begin{center}
\renewcommand{\arraystretch}{1.3}
\begin{tabular}{>{\raggedright\arraybackslash}p{3.4cm}>{\raggedright\arraybackslash}p{3.6cm}>{\raggedright\arraybackslash}p{3.6cm}>{\raggedright\arraybackslash}p{3.6cm}}
\toprule
 & Sistema linear (Vandermonde) & Lagrange & Newton\\
\midrule
Base usada & $1,x,\dots,x^{n-1}$ & $L_1,\dots,L_n$ & $1,(x-x_1),(x-x_1)(x-x_2),\dots$\\
Custo para obter $p$ & $\approx\frac23 n^3$ (eliminação) & nenhum cálculo prévio, além dos denominadores & $\approx\frac{n^2}{2}$ divisões (tabela)\\
Avaliar $p(x)$ & $n-1$ mult.\ (Horner) & $\approx n^2$ operações & $n-1$ mult.\ (forma aninhada)\\
Acrescentar um nó & refazer o sistema & refazer todas as bases & somar um termo\\
Estabilidade numérica & ruim para $n$ grande & boa & boa\\
\bottomrule
\end{tabular}
\end{center}

Em resumo: o sistema linear é o caminho mais intuitivo e mostra a ligação com a Álgebra
Linear; a forma de Lagrange é a mais elegante e a mais usada na teoria; a forma de Newton é,
em geral, a mais conveniente no computador.

% ==================================================================
\section{O erro da interpolação}

Quando os dados vêm de uma função $f$, queremos saber quão perto $p(x)$ está de $f(x)$ fora
dos nós.

\begin{teorema}[Erro da interpolação]\label{teo:erro}
Seja $f$ com derivadas contínuas até a ordem $n$ num intervalo $[a,b]$ que contém os nós
distintos $x_1,\dots,x_n$, e seja $p$ o polinômio que interpola $f$ nesses nós. Para cada
$x\in[a,b]$ existe $\xi=\xi(x)\in(a,b)$ tal que
\[
f(x)-p(x)=\frac{f^{(n)}(\xi)}{n!}\,(x-x_1)(x-x_2)\cdots(x-x_n).
\]
Em particular, se $|f^{(n)}(t)|\le M_n$ em $[a,b]$,
\[
|f(x)-p(x)|\le\frac{M_n}{n!}\,\bigl|(x-x_1)(x-x_2)\cdots(x-x_n)\bigr|.
\]
\end{teorema}

\begin{proof}[Ideia da demonstração]
Fixado $x$ diferente dos nós, considere
$g(t)=f(t)-p(t)-K\prod_{i}(t-x_i)$, com $K$ escolhido para que $g(x)=0$. Então $g$ tem
$n+1$ zeros em $[a,b]$ (os nós e $x$). Aplicando o Teorema de Rolle repetidamente,
$g^{(n)}$ tem um zero $\xi$. Como $p^{(n)}\equiv0$ e a derivada de ordem $n$ de
$\prod_i(t-x_i)$ é $n!$, obtemos $0=f^{(n)}(\xi)-K\,n!$, isto é, $K=f^{(n)}(\xi)/n!$.
\end{proof}

\begin{exemplo}[: estimativa de erro]
Interpolando $f(x)=e^x$ nos nós $0$; $0{,}5$; $1$ e avaliando em $x=0{,}25$, obtemos
$p(0{,}25)\approx1{,}27176$, enquanto $e^{0{,}25}\approx1{,}28403$. O erro real é
$\approx0{,}01227$. Pelo Teorema~\ref{teo:erro}, com $n=3$ e $M_3=\max_{[0,1]}e^x=e$,
\[
|f(0{,}25)-p(0{,}25)|\le\frac{e}{3!}\,|0{,}25\cdot(-0{,}25)\cdot(-0{,}75)|\approx0{,}02124,
\]
uma cota coerente com o erro observado.
\end{exemplo}

\paragraph{Mais pontos nem sempre é melhor: o fenômeno de Runge.} O produto
$\prod(x-x_i)$ pode ficar grande perto das extremidades do intervalo quando os nós são
igualmente espaçados. Para $f(x)=\dfrac{1}{1+25x^2}$ em $[-1,1]$ com 11 nós igualmente
espaçados, o polinômio interpolador de grau 10 oscila fortemente perto de $\pm1$, com erro
máximo próximo de $1{,}9$ (Figura~\ref{fig:runge}). Nesses casos usam-se nós de Chebyshev
(mais concentrados nas extremidades) ou interpolação por partes (\emph{splines}), assuntos
que aparecem mais adiante.

\begin{figure}[h]
\centering
\begin{tikzpicture}
\begin{axis}[width=11cm,height=6.5cm,axis lines=middle,xmin=-1.1,xmax=1.1,ymin=-0.6,ymax=2.1,
  grid=both,grid style={gray!20},legend style={at={(0.5,0.98)},anchor=north,font=\small},legend columns=2]
\addplot[domain=-1:1,samples=200,very thick,azul] {1/(1+25*x^2)};
\addplot[thick,vermelho] coordinates {(-1.0000,0.0385)(-0.9800,1.2303)(-0.9600,1.8044)(-0.9400,1.9589)(-0.9200,1.8458)(-0.9000,1.5787)(-0.8800,1.2402)(-0.8600,0.8881)(-0.8400,0.5604)(-0.8200,0.2802)(-0.8000,0.0588)(-0.7800,-0.1007)(-0.7600,-0.2013)(-0.7400,-0.2496)(-0.7200,-0.2546)(-0.7000,-0.2262)(-0.6800,-0.1743)(-0.6600,-0.1083)(-0.6400,-0.0363)(-0.6200,0.0349)(-0.6000,0.1000)(-0.5800,0.1554)(-0.5600,0.1987)(-0.5400,0.2292)(-0.5200,0.2472)(-0.5000,0.2538)(-0.4800,0.2510)(-0.4600,0.2414)(-0.4400,0.2278)(-0.4200,0.2131)(-0.4000,0.2000)(-0.3800,0.1912)(-0.3600,0.1888)(-0.3400,0.1946)(-0.3200,0.2099)(-0.3000,0.2354)(-0.2800,0.2710)(-0.2600,0.3165)(-0.2400,0.3708)(-0.2200,0.4326)(-0.2000,0.5000)(-0.1800,0.5709)(-0.1600,0.6432)(-0.1400,0.7142)(-0.1200,0.7817)(-0.1000,0.8434)(-0.0800,0.8971)(-0.0600,0.9409)(-0.0400,0.9734)(-0.0200,0.9933)(0.0000,1.0000)(0.0200,0.9933)(0.0400,0.9734)(0.0600,0.9409)(0.0800,0.8971)(0.1000,0.8434)(0.1200,0.7817)(0.1400,0.7142)(0.1600,0.6432)(0.1800,0.5709)(0.2000,0.5000)(0.2200,0.4326)(0.2400,0.3708)(0.2600,0.3165)(0.2800,0.2710)(0.3000,0.2354)(0.3200,0.2099)(0.3400,0.1946)(0.3600,0.1888)(0.3800,0.1912)(0.4000,0.2000)(0.4200,0.2131)(0.4400,0.2278)(0.4600,0.2414)(0.4800,0.2510)(0.5000,0.2538)(0.5200,0.2472)(0.5400,0.2292)(0.5600,0.1987)(0.5800,0.1554)(0.6000,0.1000)(0.6200,0.0349)(0.6400,-0.0363)(0.6600,-0.1083)(0.6800,-0.1743)(0.7000,-0.2262)(0.7200,-0.2546)(0.7400,-0.2496)(0.7600,-0.2013)(0.7800,-0.1007)(0.8000,0.0588)(0.8200,0.2802)(0.8400,0.5604)(0.8600,0.8881)(0.8800,1.2402)(0.9000,1.5787)(0.9200,1.8458)(0.9400,1.9589)(0.9600,1.8044)(0.9800,1.2303)(1.0000,0.0385)};
\legend{$f(x)=\frac{1}{1+25x^2}$,$p_{10}$ (11 nós igualmente espaçados)}
\addplot[only marks,mark=*,mark size=1.5pt] coordinates {(-1.0,0.0385)(-0.8,0.0588)(-0.6,0.1000)(-0.4,0.2000)(-0.2,0.5000)(0.0,1.0000)(0.2,0.5000)(0.4,0.2000)(0.6,0.1000)(0.8,0.0588)(1.0,0.0385)};
\end{axis}
\end{tikzpicture}
\caption{Fenômeno de Runge: aumentar o número de nós igualmente espaçados pode piorar a
aproximação perto das extremidades.}
\label{fig:runge}
\end{figure}

% ==================================================================
\section{Exercícios}

\begin{enumerate}[label=\textbf{\arabic*.},leftmargin=*]
\item Encontre o polinômio que interpola $(0,1)$, $(1,3)$, $(2,7)$ pelos três métodos:
  sistema linear, Lagrange e Newton. Confira que os resultados coincidem.

\item Considere os pontos $(-1,4)$, $(0,1)$, $(1,0)$, $(2,1)$.
  \begin{enumerate}[label=(\alph*),nosep]
  \item Monte a tabela de diferenças divididas e escreva $p$ na forma de Newton.
  \item Qual é o grau de $p$? Por que ele é menor que 3?
  \item Acrescente o ponto $(3,4)$. Quantos termos novos aparecem na forma de Newton?
  \end{enumerate}

\item A temperatura de um forno foi medida em três instantes:
  \begin{center}
  \begin{tabular}{c|ccc}
  $t$ (min) & 1 & 3 & 4\\ \hline
  $T$ ($^\circ$C) & 21{,}2 & 24{,}8 & 25{,}6
  \end{tabular}
  \end{center}
  Use interpolação quadrática para estimar a temperatura em $t=2$ min.

\item Interpole $f(x)=\operatorname{sen}x$ nos nós $0$, $\pi/4$ e $\pi/2$. Estime
  $\operatorname{sen}(\pi/8)$, calcule o erro real e compare com a cota do
  Teorema~\ref{teo:erro}.

\item Mostre que, para quaisquer nós distintos, $\displaystyle\sum_{i=1}^{n}L_i(x)=1$ para
  todo $x$. (Sugestão: qual polinômio interpola a função constante $f(x)=1$?)

\item\label{ex:vandermonde} Mostre que, para três nós,
  $\det\begin{pmatrix}1&x_1&x_1^2\\1&x_2&x_2^2\\1&x_3&x_3^2\end{pmatrix}
  =(x_2-x_1)(x_3-x_1)(x_3-x_2)$.

\item \textbf{(GeoGebra)} Em cada um dos três aplicativos, coloque os quatro pontos sobre a
  reta $y=2x+2$ (por exemplo, $(0,2)$, $(1,4)$, $(2,6)$, $(3,8)$). Descreva o que acontece
  com: (a) a última linha da matriz escalonada; (b) a combinação das bases de Lagrange;
  (c) as diferenças divididas de ordem 2 e 3.

\item \textbf{(GeoGebra)} No aplicativo de Newton, mova somente o ponto $D$ e anote quais
  números mudam na tabela. Faça o mesmo no aplicativo de Lagrange. Qual dos dois métodos
  reaproveita mais contas quando um único dado muda?
\end{enumerate}

\subsection*{Respostas}
\begin{enumerate}[label=\textbf{\arabic*.},leftmargin=*,nosep]
\item $p(x)=x^2+x+1$. Na forma de Newton: $p(x)=1+2x+x(x-1)$.
\item (a) $f[x_i]$: $4,1,0,1$; ordem 1: $-3,-1,1$; ordem 2: $1,1$; ordem 3: $0$. Logo
  $p(x)=4-3(x+1)+(x+1)x=x^2-2x+1=(x-1)^2$. (b) Grau 2, pois $f[x_1,\dots,x_4]=0$: os quatro
  pontos estão sobre uma mesma parábola. (c) Um termo, com coeficiente
  $f[x_1,\dots,x_5]=0$; o ponto $(3,4)$ também está sobre a parábola.
\item $T(t)=\frac{92}{5}+\frac{47}{15}t-\frac13t^2$, e $T(2)=\frac{70}{3}\approx23{,}3\,^\circ$C.
\item $p(\pi/8)\approx0{,}40533$, $\operatorname{sen}(\pi/8)\approx0{,}38268$, erro
  $\approx0{,}02265$; cota $\frac16\,|\tfrac{\pi}{8}(\tfrac{\pi}{8}-\tfrac{\pi}{4})(\tfrac{\pi}{8}-\tfrac{\pi}{2})|\approx0{,}03028$.
\item O polinômio que interpola $f\equiv1$ é $p\equiv1$; pela forma de Lagrange, $p=\sum 1\cdot L_i$.
\item Escalone (ou desenvolva o determinante) e fatore.
\item (a) na matriz escalonada, os termos independentes das linhas 3 e 4 ficam iguais a 0, logo $a_3=a_2=0$;
  (b) $p(x)=2x+2$; (c) todas nulas.
\item Em Newton mudam apenas as entradas que envolvem $x_4$ (a última ``diagonal'' da tabela,
  de $f[x_4]$ até $f[x_1,\dots,x_4]$) e só o último termo do polinômio; $p_0,p_1,p_2$ não mudam.
  Em Lagrange, como $x_4$ aparece em todos os denominadores e numeradores, todas as bases mudam.
  Newton reaproveita mais contas.
\end{enumerate}

% Capítulo gerado a partir de ajuste_de_curvas.tex (arquivos usados neste capítulo)
\def\arqRetaRol{metodo_dos_minimos_quadrados_reta_rolagem.ggb}
\def\arqParabRol{metodo_dos_minimos_quadrados_exemplo02_parabola_rolagem.ggb}
\def\arqAproxRol{aproximacao_funcao_mmq_rolagem.ggb}
\def\arqMMQSimples{01_metodo_dos_minimos_quadrados_linear_simples_TCC_2026.ggb}
\def\arqGDReta{02_RegressaoLinearUsandoGradienteDescendente_junho.ggb}
\def\arqMMQComb{03_metodo_dos_minimos_quadrados_linear_combinada_TCC_2026.ggb}
\def\arqGDPadr{04_aproximacaoGDescCombFuncoespadronizada_03_07.ggb}
\def\arqGDComb{05_aproximacaoGDescCombFuncao.ggb}
\def\arqGDEsc{06_aproximacaoGDescCombFuncao_Escalonada.ggb}
\chapter[Ajuste de curvas]{Ajuste de curvas: método dos mínimos quadrados e gradiente descendente}\label{cap:ajuste}

\begin{tcolorbox}[colback=gray!4,colframe=gray!50,title={Como usar os aplicativos deste capítulo},fonttitle=\bfseries]
Os quadros laranja \emph{Atividade no GeoGebra} trazem o link \textbf{[abrir o aplicativo]}.
São nove aplicativos, em dois grupos.

\textbf{Método direto (sistema normal).}
\begin{itemize}[nosep]
  \item \ggbabrir{\arqRetaRol}{reta com 4 pontos}, \ggbabrir{\arqParabRol}{parábola com 11 pontos} e
        \ggbabrir{\arqAproxRol}{parábola com os pontos $A,\dots,K$}: as contas aparecem passo a passo,
        com barras de rolagem, e o gráfico fica na segunda janela;
  \item \ggbabrir{\arqMMQSimples}{reta com 5 pontos} e
        \ggbabrir{\arqMMQComb}{combinação de $\operatorname{sen}x$, $\cos x$, $x^2$, $x$ e $1$}, montados na planilha.
\end{itemize}

\textbf{Método iterativo (gradiente descendente).}
\begin{itemize}[nosep]
  \item \ggbabrir{\arqGDReta}{reta com 5 pontos}, com as curvas de nível do erro;
  \item \ggbabrir{\arqGDComb}{combinação de funções sem padronizar},
        \ggbabrir{\arqGDEsc}{com as funções padronizadas} e
        \ggbabrir{\arqGDPadr}{padronizadas, com volta aos coeficientes originais}.
\end{itemize}
Nos aplicativos iterativos, o botão de iteração faz um passo do método e o botão \emph{Limpar}
recomeça; o \emph{passo} (taxa de aprendizagem) pode ser alterado na caixa de texto.

Para que os links funcionem, mantenha este PDF e a pasta \texttt{ggb} juntos e abra o PDF num leitor
que permita abrir arquivos externos (Adobe Acrobat Reader, Okular, SumatraPDF, Foxit).
\end{tcolorbox}

% ==================================================================
\section{Interpolar ou ajustar?}

No capítulo de interpolação, procurávamos um polinômio que passasse \emph{exatamente} pelos pontos
dados. Isso nem sempre é desejável:
\begin{itemize}[nosep]
  \item os dados experimentais vêm com erros de medida, e forçar a curva a passar por todos eles
        reproduz também esses erros;
  \item com muitos pontos, o polinômio interpolador tem grau alto e oscila entre os nós;
  \item muitas vezes já se sabe que tipo de função descreve o fenômeno (uma reta, uma parábola, uma
        exponencial) e o que se quer são os seus parâmetros.
\end{itemize}

No \textbf{ajuste de curvas}, escolhemos uma família de funções e procuramos, dentro dela, a que
fica ``mais perto'' dos pontos, sem exigir que passe por eles.

\begin{center}
\begin{tikzpicture}
\begin{axis}[width=0.62\textwidth,height=5.2cm,xmin=-0.3,xmax=2.3,ymin=0,ymax=2,
  axis lines=middle,xlabel={$x$},ylabel={$y$},legend pos=north west,legend style={font=\small}]
\addplot[only marks,mark=*,mark size=2pt] coordinates {(0,1) (0.5,0.5) (1,1.5) (1.5,1) (2,1.5)};
\addlegendentry{dados}
\addplot[domain=-0.2:2.2,thick,red] {0.8+0.3*x};
\addlegendentry{$\varphi(x)=0{,}8+0{,}3x$}
\addplot[orange,thick] coordinates {(0,1) (0,0.8)};
\addplot[orange,thick] coordinates {(0.5,0.5) (0.5,0.95)};
\addplot[orange,thick] coordinates {(1,1.5) (1,1.1)};
\addplot[orange,thick] coordinates {(1.5,1) (1.5,1.25)};
\addplot[orange,thick] coordinates {(2,1.5) (2,1.4)};
\end{axis}
\end{tikzpicture}

\small Os segmentos laranja são os resíduos $r_k=f(x_k)-\varphi(x_k)$.
\end{center}

% ==================================================================
\section{O método dos mínimos quadrados (caso discreto)}

São dados $m$ pontos $(x_1,f(x_1)),\dots,(x_m,f(x_m))$ e $n$ funções $g_1(x),\dots,g_n(x)$,
escolhidas conforme o problema. Procuramos
\[
\varphi(x)=\alpha_1g_1(x)+\alpha_2g_2(x)+\dots+\alpha_ng_n(x)
\]
de modo que os \emph{resíduos} $r_k=f(x_k)-\varphi(x_k)$ sejam pequenos. Como resíduos positivos e
negativos se cancelariam numa soma simples, minimizamos a soma dos seus quadrados:
\[
D(\alpha_1,\dots,\alpha_n)=\sum_{k=1}^{m}\big[f(x_k)-\varphi(x_k)\big]^2
=\sum_{k=1}^{m}\Big[f(x_k)-\sum_{j=1}^{n}\alpha_jg_j(x_k)\Big]^2 .
\]
Exemplos de famílias: a reta ($g_1=1$, $g_2=x$), a parábola ($g_1=1$, $g_2=x$, $g_3=x^2$) e a
combinação $\alpha_1\operatorname{sen}x+\alpha_2\cos x+\alpha_3x^2+\alpha_4x+\alpha_5$.

\subsection{O sistema normal}

$D$ é uma função quadrática dos $\alpha_j$, limitada inferiormente por zero. No ponto de mínimo, todas
as derivadas parciais se anulam:
\[
\frac{\partial D}{\partial\alpha_i}=-2\sum_{k=1}^{m}\Big[f(x_k)-\sum_{j=1}^{n}\alpha_jg_j(x_k)\Big]g_i(x_k)=0,
\qquad i=1,\dots,n,
\]
ou seja,
\[
\sum_{j=1}^{n}\Big(\sum_{k=1}^{m}g_i(x_k)g_j(x_k)\Big)\alpha_j=\sum_{k=1}^{m}f(x_k)g_i(x_k),\qquad i=1,\dots,n.
\]

\begin{definicao}[sistema normal]
Com $a_{ij}=\sum_{k=1}^{m}g_i(x_k)g_j(x_k)$ e $b_i=\sum_{k=1}^{m}f(x_k)g_i(x_k)$, os coeficientes
do ajuste são a solução do sistema linear $n\times n$
\[
\begin{bmatrix}a_{11}&\cdots&a_{1n}\\ \vdots&&\vdots\\ a_{n1}&\cdots&a_{nn}\end{bmatrix}
\begin{bmatrix}\alpha_1\\ \vdots\\ \alpha_n\end{bmatrix}=
\begin{bmatrix}b_1\\ \vdots\\ b_n\end{bmatrix}.
\]
A matriz é simétrica ($a_{ij}=a_{ji}$). Se os vetores $(g_j(x_1),\dots,g_j(x_m))$ forem linearmente
independentes, ela é definida positiva, o sistema tem solução única e essa solução é o ponto de mínimo de $D$.
\end{definicao}

Em notação matricial, com $G$ a matriz $m\times n$ de entradas $G_{kj}=g_j(x_k)$ e $\mathbf f=(f(x_1),\dots,f(x_m))^T$,
o sistema normal é $G^TG\,\boldsymbol\alpha=G^T\mathbf f$. É o \textbf{método direto}: monta-se e
resolve-se um sistema linear, por exemplo por eliminação de Gauss (Capítulo~\ref{cap:sistemas}).

\subsection{Ajuste por uma reta}

Para $\varphi(x)=\alpha_1+\alpha_2x$ ($g_1=1$, $g_2=x$), o sistema normal fica
\[
\begin{bmatrix} m & \sum x_k\\ \sum x_k & \sum x_k^2\end{bmatrix}
\begin{bmatrix}\alpha_1\\ \alpha_2\end{bmatrix}=
\begin{bmatrix}\sum f(x_k)\\ \sum x_kf(x_k)\end{bmatrix}.
\]

\begin{exemplo}[{: reta por 5 pontos}]
Os pontos são $(0;1)$, $(0{,}5;\,0{,}5)$, $(1;\,1{,}5)$, $(1{,}5;\,1)$ e $(2;\,1{,}5)$.

\begin{center}
\begin{tabular}{c|rr|rrr}
\toprule
$k$ & $x_k$ & $f(x_k)$ & $g_1g_2=x_k$ & $g_2g_2=x_k^2$ & $fg_2=x_kf(x_k)$\\
\midrule
1 & 0   & 1   & 0   & 0    & 0\\
2 & 0,5 & 0,5 & 0,5 & 0,25 & 0,25\\
3 & 1   & 1,5 & 1   & 1    & 1,5\\
4 & 1,5 & 1   & 1,5 & 2,25 & 1,5\\
5 & 2   & 1,5 & 2   & 4    & 3\\
\midrule
$\sum$ & 5 & 5,5 & 5 & 7,5 & 6,25\\
\bottomrule
\end{tabular}
\end{center}

Como $g_1g_1=1$, $a_{11}=5$. O sistema e a eliminação de Gauss são
\[
\left[\begin{array}{cc|c}5&5&5{,}5\\5&7{,}5&6{,}25\end{array}\right]
\xrightarrow{\;m_{21}=1\;}
\left[\begin{array}{cc|c}5&5&5{,}5\\0&2{,}5&0{,}75\end{array}\right],
\]
logo $\alpha_2=0{,}75/2{,}5=0{,}3$ e $\alpha_1=(5{,}5-5\cdot0{,}3)/5=0{,}8$. Então
$\varphi(x)=0{,}8+0{,}3x$, com $D=\sum r_k^2=0{,}475$.
\end{exemplo}

\begin{atividade}{reta pelo método dos mínimos quadrados}{\arqMMQSimples}
O aplicativo monta na planilha a tabela de $x_k$, $f(x_k)$, $g_1(x_k)=1$, $g_2(x_k)=x_k$, dos produtos
$g_ig_j$ e $fg_i$ e das somas, e resolve o sistema normal por eliminação de Gauss ($m_{21}$, $\alpha_2$,
$\alpha_1$). A reta $\varphi$ aparece no gráfico com os pontos $A_0,\dots,A_4$.

\textbf{Tarefas.} (a) Confira a tabela e a reta do exemplo. (b) Mova um ponto para longe dos demais
(um \emph{dado discrepante}) e observe quanto a reta se desloca. (c) Alinhe os cinco pontos numa reta:
qual é então o valor de $D$?
\end{atividade}

\begin{atividade}{reta com rolagem (4 pontos)}{\arqRetaRol}
Todas as etapas aparecem em textos na janela da esquerda, com barras de rolagem: tabela dos produtos,
sistema normal, eliminação de Gauss, substituição regressiva, tabela de resíduos e comparação com a
parábola. Na janela da direita estão os pontos (que podem ser arrastados), a reta, os resíduos e a
conferência com o comando \texttt{FitPoly} do GeoGebra. Com os pontos iniciais,
$\varphi(x)=1{,}0735+0{,}2941x$.

\textbf{Tarefas.} (a) Gere pontos aleatórios e compare $D$ para a reta e para a parábola.
(b) Explique por que o $D$ da parábola nunca é maior que o da reta.
\end{atividade}

\subsection{Ajuste por uma parábola}

Para $\varphi(x)=\alpha_1+\alpha_2x+\alpha_3x^2$ ($g_1=1$, $g_2=x$, $g_3=x^2$):
\[
\begin{bmatrix} m & \sum x_k & \sum x_k^2\\ \sum x_k & \sum x_k^2 & \sum x_k^3\\ \sum x_k^2 & \sum x_k^3 & \sum x_k^4\end{bmatrix}
\begin{bmatrix}\alpha_1\\ \alpha_2\\ \alpha_3\end{bmatrix}=
\begin{bmatrix}\sum f(x_k)\\ \sum x_kf(x_k)\\ \sum x_k^2f(x_k)\end{bmatrix}.
\]

\begin{atividade}{parábola com 11 pontos}{\arqParabRol}
Mesmo formato do aplicativo da reta, agora com $g_3(x)=x^2$, três multiplicadores ($m_{21}$, $m_{31}$,
$m_{32}$) e os pontos $A_0,\dots,A_{10}$, dispostos em forma de ``U''. O resultado inicial é
$\varphi(x)=0{,}2144+0{,}4555x+1{,}8492x^2$.

\textbf{Tarefas.} (a) Compare o $R^2$ da parábola e o da reta: qual modelo descreve melhor esses dados?
(b) Use o campo ``avaliar $\varphi$ em $x$'' para estimar o valor em $x=1{,}2$, fora do intervalo dos
dados (extrapolação). Que cuidado se deve ter?
\end{atividade}

\begin{atividade}{parábola com os pontos $A,\dots,K$}{\arqAproxRol}
Outro conjunto de 11 pontos, agora mais espalhados. O ajuste é
$\varphi(x)=0{,}7873-0{,}1454x+0{,}2068x^2$, com $D\approx2{,}06$ e $R^2\approx0{,}77$, contra
$D\approx7{,}86$ e $R^2\approx0{,}11$ da reta.

\textbf{Tarefa.} Mostre os resíduos e identifique o ponto que mais contribui para $D$.
\end{atividade}

\subsection{Combinação de outras funções}

O método não exige polinômios. Basta que $\varphi$ seja \emph{linear nos parâmetros}.

\begin{exemplo}[{: $\varphi(x)=\alpha_1\operatorname{sen}x+\alpha_2\cos x+\alpha_3x^2+\alpha_4x+\alpha_5$}]\label{ex:comb}
Os 8 pontos são
\[
(-2;2),\ (0;0),\ (4;2),\ (8;0),\ (10;5),\ (14;4),\ (16;8),\ (20;10).
\]
O sistema normal é $5\times5$, com $g_1=\operatorname{sen}x$, $g_2=\cos x$, $g_3=x^2$, $g_4=x$ e $g_5=1$:
\[
\begin{bmatrix}
4{,}572 & 1{,}969 & 478{,}80 & 28{,}787 & 0{,}395\\
1{,}969 & 3{,}428 & -160{,}47 & -16{,}584 & -1{,}467\\
478{,}80 & -160{,}47 & 278\,320 & 16\,408 & 1\,036\\
28{,}787 & -16{,}584 & 16\,408 & 1\,036 & 70\\
0{,}395 & -1{,}467 & 1\,036 & 70 & 8
\end{bmatrix}
\begin{bmatrix}\alpha_1\\ \alpha_2\\ \alpha_3\\ \alpha_4\\ \alpha_5\end{bmatrix}=
\begin{bmatrix}4{,}736\\ -9{,}368\\ 7\,372\\ 438\\ 31\end{bmatrix}.
\]
A solução é
\[
\varphi(x)\approx-2{,}2594\operatorname{sen}x+0{,}1701\cos x+0{,}02518x^2+0{,}09374x-0{,}06376,
\qquad D\approx0{,}5606 .
\]
Repare na diferença de escala entre as entradas: $4{,}6$ na diagonal para $\operatorname{sen}x$ e
$278\,320$ para $x^2$. Ela não atrapalha o método direto, mas, como veremos, prejudica muito o gradiente descendente.
\end{exemplo}

\begin{center}
\begin{tikzpicture}
\begin{axis}[width=0.75\textwidth,height=5.5cm,xmin=-3,xmax=21,ymin=-2,ymax=13,axis lines=middle,
  xlabel={$x$},ylabel={$y$},legend pos=south east,legend style={font=\small}]
\addplot[only marks,mark=*,mark size=2pt] coordinates {(-2,2) (0,0) (4,2) (8,0) (10,5) (14,4) (16,8) (20,10)};
\addlegendentry{dados}
\addplot[domain=-2.5:20.5,samples=300,thick,red] {-2.259395*sin(deg(x))+0.17013*cos(deg(x))+0.025184*x^2+0.093738*x-0.063763};
\addlegendentry{$\varphi(x)$}
\end{axis}
\end{tikzpicture}
\end{center}

\begin{atividade}{combinação de funções pelo método direto}{\arqMMQComb}
A planilha calcula os produtos $g_ig_j$ e $fg_i$ para as cinco funções, as somas que formam o sistema
normal $5\times5$, os multiplicadores da eliminação de Gauss e os coeficientes $\alpha_1,\dots,\alpha_5$
por substituição regressiva.

\textbf{Tarefas.} (a) Confira os valores do Exemplo~\ref{ex:comb}. (b) Retire a função $\operatorname{sen}x$
(isto é, ajuste só com $\cos x$, $x^2$, $x$ e $1$). O ajuste piora muito? Por quê?
\end{atividade}

\subsection{Casos não lineares: linearização}

Se $\varphi$ não é linear nos parâmetros, às vezes uma mudança de variáveis resolve. Para
$\varphi(x)=a\,e^{bx}$, tomando logaritmos:
\[
z=\ln\varphi(x)=\ln a+bx=\alpha_1+\alpha_2x,
\]
e ajustamos uma reta aos pontos $(x_k,\ln f(x_k))$. Com os dados
\begin{center}\small
\begin{tabular}{c|cccccccc}
$x_k$ & $-1$ & $-0{,}7$ & $-0{,}4$ & $-0{,}1$ & $0{,}2$ & $0{,}5$ & $0{,}8$ & $1$\\ \hline
$f(x_k)$ & 36,547 & 17,264 & 8,155 & 3,852 & 1,820 & 0,860 & 0,406 & 0,246
\end{tabular}
\end{center}
obtemos $\alpha_1\approx1{,}0986$ e $\alpha_2\approx-2{,}5002$, logo $a=e^{\alpha_1}\approx3{,}0$ e
$\varphi(x)\approx3{,}0\,e^{-2{,}5x}$.

\begin{observacao}
O ajuste linearizado minimiza os quadrados dos resíduos de $\ln f$, e não os de $f$. O resultado é bom,
mas não é exatamente o mínimo de $\sum[f(x_k)-a e^{bx_k}]^2$. Para esse mínimo, é preciso um método
iterativo, como o gradiente descendente da próxima seção.
\end{observacao}

% ==================================================================
\section{O método iterativo: gradiente descendente}

Em vez de igualar o gradiente a zero e resolver um sistema, podemos \emph{caminhar} em direção ao mínimo.
O gradiente de uma função aponta na direção de maior crescimento; andando no sentido oposto, a função
diminui. É a ideia do \textbf{gradiente descendente}, o método usado para treinar modelos de regressão
e redes neurais em aprendizado de máquina.

\subsection{A função custo e o seu gradiente}

Nos aplicativos, o erro é escrito com os resíduos no sentido do modelo, $e_k=\varphi(x_k)-f(x_k)$, e
dividido por $2m$, como é costume em aprendizado de máquina:
\[
J(\boldsymbol\alpha)=\frac{1}{2m}\sum_{k=1}^{m}e_k^2=\frac{1}{2m}\sum_{k=1}^{m}\Big[\sum_{j=1}^{n}\alpha_jg_j(x_k)-f(x_k)\Big]^2=\frac{D(\boldsymbol\alpha)}{2m}.
\]
Como $J$ é um múltiplo positivo de $D$, os dois têm o mesmo ponto de mínimo. As derivadas parciais são
\[
\frac{\partial J}{\partial\alpha_j}=\frac1m\sum_{k=1}^{m}e_k\,g_j(x_k),\qquad j=1,\dots,n .
\]

\begin{tcolorbox}[colback=azul!3,colframe=azul!60,title={Algoritmo: gradiente descendente},fonttitle=\bfseries]
\begin{description}[nosep,font=\normalfont\bfseries]
  \item[Passo 0:] escolha $\boldsymbol\alpha^{(0)}$ (por exemplo, zero), o passo $\eta>0$, a tolerância e o número máximo de iterações.
  \item[Passo 1:] calcule os resíduos $e_k=\varphi(x_k)-f(x_k)$ com os coeficientes atuais.
  \item[Passo 2:] calcule o gradiente $d_j=\frac1m\sum_k e_k\,g_j(x_k)$, $j=1,\dots,n$.
  \item[Passo 3:] atualize todos os coeficientes \emph{simultaneamente}: $\alpha_j^{(i+1)}=\alpha_j^{(i)}-\eta\,d_j$.
  \item[Passo 4:] se $\|\mathbf d\|$ (ou a variação de $J$) for menor que a tolerância, pare; caso contrário, volte ao passo 1.
\end{description}
\end{tcolorbox}

O passo $\eta$ é chamado \emph{taxa de aprendizagem}. Nos aplicativos, a atualização simultânea é feita
guardando os valores antigos (\texttt{a1\_antigo}, \texttt{d1\_antigo}, \dots) antes de alterar os coeficientes.

\subsection{Gradiente descendente para a reta}\label{sec:gdreta}

No aplicativo da reta, o modelo é $\varphi(x)=\alpha_1x+\alpha_2$ (aqui $\alpha_1$ é a inclinação e
$\alpha_2$ o termo independente) e
\[
\frac{\partial J}{\partial\alpha_1}=\frac1m\sum e_kx_k,\qquad \frac{\partial J}{\partial\alpha_2}=\frac1m\sum e_k .
\]

\begin{exemplo}[{: as primeiras iterações, com $\eta=0{,}1$ e os 5 pontos do exemplo anterior}]
Começando de $\alpha_1=\alpha_2=0$, os resíduos são $e=(-1;\,-0{,}5;\,-1{,}5;\,-1;\,-1{,}5)$, e então
\[
\frac{\partial J}{\partial\alpha_1}=\frac{0-0{,}25-1{,}5-1{,}5-3}{5}=-1{,}25,\qquad
\frac{\partial J}{\partial\alpha_2}=\frac{-5{,}5}{5}=-1{,}1,
\]
\[
\alpha_1^{(1)}=0-0{,}1(-1{,}25)=0{,}125,\qquad \alpha_2^{(1)}=0-0{,}1(-1{,}1)=0{,}11 .
\]
Continuando:

\begin{center}
\begin{tabular}{rccc}
\toprule
iteração & $\alpha_1$ (inclinação) & $\alpha_2$ (termo indep.) & $D=\sum e_k^2$\\
\midrule
0 & 0 & 0 & 6,75\\
1 & 0,125 & 0,11 & 4,2927\\
2 & 0,22025 & 0,1965 & 2,8251\\
3 & 0,29256 & 0,26483 & 1,9473\\
10 & 0,47673 & 0,48470 & 0,6491\\
50 & 0,39059 & 0,68397 & 0,4988\\
100 & 0,32991 & 0,76170 & 0,4776\\
500 & 0,30000 & 0,79999 & 0,4750\\
\bottomrule
\end{tabular}
\end{center}

O método converge para $\alpha_1=0{,}3$ e $\alpha_2=0{,}8$, a mesma reta $\varphi(x)=0{,}3x+0{,}8$ do método direto,
mas precisou de algumas centenas de iterações.
\end{exemplo}

\begin{center}
\begin{tikzpicture}
\begin{axis}[width=0.46\textwidth,height=5.2cm,xlabel={iteração},ylabel={$J$},
  title={\small $J$ ao longo das iterações ($\eta=0{,}1$)},xmin=0,xmax=300,ymin=0,ymax=0.7,
  yticklabel style={/pgf/number format/fixed,/pgf/number format/use comma}]
\addplot[thick,azul] coordinates {(0,0.675) (1,0.429269) (2,0.282505) (3,0.194728) (4,0.142113) (5,0.110462) (6,0.0913175) (7,0.0796366) (8,0.072415) (9,0.0678619) (10,0.0649096) (20,0.0565007) (30,0.0532653) (40,0.0512007) (50,0.0498754) (60,0.0490248) (70,0.0484788) (80,0.0481283) (90,0.0479033) (100,0.0477589) (120,0.0476067) (140,0.0475439) (160,0.0475181) (200,0.0475031) (250,0.0475003) (300,0.0475)};
\end{axis}
\end{tikzpicture}\hfill
\begin{tikzpicture}
\begin{axis}[width=0.46\textwidth,height=5.2cm,xlabel={$\alpha_1$},ylabel={$\alpha_2$},
  title={\small caminho no plano dos parâmetros},xmin=-0.05,xmax=0.6,ymin=-0.05,ymax=0.9]
\addplot[thick,red,mark=*,mark size=0.8pt] coordinates {(0.0000,0.0000) (0.1250,0.1100) (0.2203,0.1965) (0.2926,0.2648) (0.3472,0.3191) (0.3882,0.3625) (0.4187,0.3974) (0.4412,0.4258) (0.4574,0.4491) (0.4689,0.4684) (0.4767,0.4847) (0.4817,0.4986) (0.4846,0.5105) (0.4859,0.5210) (0.4859,0.5303) (0.4850,0.5387) (0.4812,0.5534) (0.4729,0.5718) (0.4633,0.5877) (0.4503,0.6064) (0.4379,0.6231) (0.4263,0.6381) (0.4156,0.6519) (0.4058,0.6645) (0.3968,0.6760) (0.3886,0.6865) (0.3811,0.6961) (0.3742,0.7050) (0.3466,0.7403) (0.3299,0.7617) (0.3192,0.7754) (0.3123,0.7842) (0.3079,0.7899) (0.3051,0.7935) (0.3033,0.7958) (0.3021,0.7973) (0.3013,0.7983) (0.3009,0.7989) (0.3004,0.7995)};
\addplot[only marks,mark=x,mark size=4pt,very thick] coordinates {(0.3,0.8)};
\end{axis}
\end{tikzpicture}
\end{center}

\begin{atividade}{regressão linear por gradiente descendente}{\arqGDReta}
Os pontos são $J$, $K$, $L$, $M$ e $T$, os mesmos do exemplo. O botão \textbf{Iteração} faz um passo do
método; \textbf{Limpar os Dados} volta a $\alpha_1=\alpha_2=0$. O aplicativo mostra:
\begin{itemize}[nosep]
  \item a reta $\varphi$, os resíduos e o erro $\text{SSE}=\sum e_k^2$;
  \item as tabelas de $x_k$, $y_k$ e $\hat y_k=\varphi(x_k)$;
  \item as curvas de nível do erro no plano $(\alpha_1,\alpha_2)$ (elipses), com o ponto $(\alpha_1,\alpha_2)$
        caminhando em direção ao centro.
\end{itemize}
\textbf{Tarefas.} (a) Faça 10 iterações e confira a tabela do exemplo ($\alpha_1\approx0{,}4767$,
$\alpha_2\approx0{,}4847$). (b) Observe que o caminho é sempre perpendicular às curvas de nível.
(c) Repita com $\eta=0{,}5$ e $\eta=0{,}85$. (d) Tente $\eta=0{,}9$: o que acontece?
\end{atividade}

\subsection{A escolha do passo e a convergência}

Como $J$ é quadrática, $J(\boldsymbol\alpha)=J(\boldsymbol\alpha^*)+\tfrac12(\boldsymbol\alpha-\boldsymbol\alpha^*)^TH(\boldsymbol\alpha-\boldsymbol\alpha^*)$,
com $H=\frac1mG^TG$ (a matriz do sistema normal dividida por $m$). A iteração fica
\[
\boldsymbol\alpha^{(i+1)}-\boldsymbol\alpha^*=(I-\eta H)\,(\boldsymbol\alpha^{(i)}-\boldsymbol\alpha^*).
\]
É uma iteração de ponto fixo, como as do Capítulo~\ref{cap:zeros} e os métodos iterativos do
Capítulo~\ref{cap:sistemas}. Se $\lambda_{\min}$ e $\lambda_{\max}$ são o menor e o maior autovalor de $H$:
\begin{itemize}[nosep]
  \item o método converge se, e somente se, $0<\eta<2/\lambda_{\max}$;
  \item mesmo com o melhor passo, o erro cai por um fator próximo de
        $\dfrac{\kappa-1}{\kappa+1}$ por iteração, onde $\kappa=\lambda_{\max}/\lambda_{\min}$ é o
        \emph{número de condição} de $H$. Quanto maior $\kappa$, mais lenta a convergência.
\end{itemize}

No exemplo da reta, os autovalores de $H$ são $0{,}219$ e $2{,}281$. O passo máximo é
$2/2{,}281\approx0{,}877$ (por isso $\eta=0{,}9$ diverge) e $\kappa\approx10{,}4$.

\subsection{Padronização das variáveis}

No Exemplo~\ref{ex:comb}, os autovalores de $H$ vão de $0{,}051$ a $34\,912$, e $\kappa\approx6{,}8\times10^{5}$.
O passo máximo é $2/34\,912\approx5{,}7\times10^{-5}$. Com um passo tão pequeno, os coeficientes de
$\operatorname{sen}x$, $\cos x$ e $1$ quase não se movem: a função $x^2$, com valores até 400, domina o gradiente.

A solução é \textbf{padronizar} cada função: subtrair a média e dividir pelo desvio padrão dos seus valores nos pontos,
\[
\tilde g_j(x)=\frac{g_j(x)-\mu_j}{\sigma_j},\qquad
\mu_j=\frac1m\sum_kg_j(x_k),\quad \sigma_j=\sqrt{\frac1m\sum_k\big(g_j(x_k)-\mu_j\big)^2},
\]
e ajustar $\varphi(x)=\beta_1\tilde g_1(x)+\dots+\beta_4\tilde g_4(x)+\beta_5$. Agora todas as funções
têm média 0 e desvio 1 nos pontos. O número de condição cai para $\kappa\approx102$, e o passo pode ser
$\eta=0{,}1$. Como as $\tilde g_j$ têm média zero, a derivada em relação a $\beta_5$ fica
$\frac1m\sum e_k=0$ quando $\beta_5=\bar f$: por isso os aplicativos começam com $\beta_5$ igual à média dos $y_k$.

Os coeficientes originais são recuperados por
\[
\alpha_j=\frac{\beta_j}{\sigma_j}\quad(j=1,\dots,4),\qquad
\alpha_5=\beta_5-\sum_{j=1}^{4}\frac{\beta_j\mu_j}{\sigma_j}.
\]

\begin{exemplo}[{: sem e com padronização}]
Para os dados do Exemplo~\ref{ex:comb}, partindo de coeficientes nulos ($\beta_5=\bar f=3{,}875$ no caso padronizado):

\begin{center}
\begin{tabular}{rcc}
\toprule
iteração & $J$ sem padronizar ($\eta=10^{-6}$) & $J$ padronizado ($\eta=0{,}1$)\\
\midrule
0    & 13,3125 & 5,8047\\
10   & 7,1039  & 0,8488\\
50   & 1,4570  & 0,1217\\
100  & 1,1176  & 0,0476\\
500  & 1,1069  & 0,0361\\
1000 & 1,1060  & 0,0351\\
5000 & 1,0990  & 0,03504\\
\bottomrule
\end{tabular}
\end{center}

O mínimo é $J=D/16\approx0{,}03504$. Sem padronizar, depois de 5000 iterações o erro ainda é 30 vezes maior;
com padronização, o método chega ao mínimo e os $\beta_j$ convertidos dão os mesmos $\alpha_j$ do método direto.
\end{exemplo}

\begin{center}
\begin{tikzpicture}
\begin{axis}[width=0.7\textwidth,height=5.5cm,xlabel={iteração},ylabel={$J$},xmode=log,ymode=log,
  legend style={font=\small,at={(0.5,1.03)},anchor=south},legend columns=2,xmin=1,xmax=5000]
\addplot[thick,red,mark=*,mark size=1pt] coordinates {(1,12.4752) (2,11.6953) (5,9.66234) (10,7.10386) (20,4.05357) (50,1.45697) (100,1.11755) (200,1.10738) (500,1.10685) (1000,1.10597) (2000,1.10422) (5000,1.09902)};
\addlegendentry{sem padronizar, $\eta=10^{-6}$}
\addplot[thick,azul,mark=*,mark size=1pt] coordinates {(1,4.11939) (2,3.04881) (5,1.56211) (10,0.84876) (20,0.446914) (50,0.121745) (100,0.0475599) (200,0.0396608) (500,0.0361489) (1000,0.0351439) (2000,0.0350413) (5000,0.0350404)};
\addlegendentry{padronizado, $\eta=0{,}1$}
\addplot[dashed,gray,domain=1:5000] {0.03504};
\addlegendentry{mínimo (método direto)}
\end{axis}
\end{tikzpicture}
\end{center}

\begin{atividade}{combinação de funções sem padronizar}{\arqGDComb}
Funções $\varphi_1=\operatorname{sen}x$, $\varphi_2=\cos x$, $\varphi_3=x^2$, $\varphi_4=x$ e $\varphi_5=1$,
os mesmos 8 pontos do Exemplo~\ref{ex:comb} e o passo $10^{-6}$. O botão de iteração aplica um passo; o campo
\texttt{iter} conta as iterações e o texto mostra o erro $J$ e os parâmetros.

\textbf{Tarefas.} (a) Faça 50 iterações e confira $J\approx1{,}457$. (b) Aumente o passo para $10^{-4}$ e
observe a divergência ($10^{-4}>5{,}7\times10^{-5}$). (c) Compare os parâmetros com os do método direto.
\end{atividade}

\begin{atividade}{combinação de funções padronizadas}{\arqGDEsc}
As funções são padronizadas pela média e pelo desvio padrão dos seus valores nos pontos
(\texttt{medsen}, \texttt{desvsen}, \dots), e $\beta_5$ começa igual à média dos $y_k$. O passo é $0{,}1$.

\textbf{Tarefas.} (a) Com 9 iterações, confira $J\approx0{,}931$. (b) Continue até $J$ estabilizar perto de
$0{,}035$. (c) Verifique que a derivada em relação a $\beta_5$ fica praticamente nula.
\end{atividade}

\begin{atividade}{padronizado, com volta aos coeficientes originais}{\arqGDPadr}
Mesmo ajuste, com passo $0{,}01$ e controles deslizantes para $\beta_1,\dots,\beta_4$ (também se pode
procurar o mínimo ``à mão''). O aplicativo calcula $\alpha_j=\beta_j/\sigma_j$ e
$\alpha_5=\beta_5-\sum\beta_j\mu_j/\sigma_j$.

\textbf{Tarefa.} Depois de convergir, compare os $\alpha_j$ com os do aplicativo do método direto
(\ggbabrir{\arqMMQComb}{combinação de funções}).
\end{atividade}

% ==================================================================
\section{Método direto \texorpdfstring{$\times$}{x} método iterativo}

\begin{center}
\small
\begin{tabular}{p{3.2cm}p{5.2cm}p{5.2cm}}
\toprule
 & método direto (sistema normal) & gradiente descendente\\
\midrule
ideia & igualar o gradiente a zero e resolver $G^TG\boldsymbol\alpha=G^T\mathbf f$ & caminhar contra o gradiente: $\boldsymbol\alpha\leftarrow\boldsymbol\alpha-\eta\nabla J$\\
resultado & exato (a menos de arredondamento) em número finito de operações & aproximação que melhora a cada iteração\\
custo & montar $G^TG$ e resolver um sistema $n\times n$ (cerca de $n^3/3$ operações) & cada iteração custa da ordem de $mn$ operações\\
parâmetros & nenhum & passo $\eta$, chute inicial, critério de parada\\
sensibilidade & mal condicionamento de $G^TG$ afeta a precisão & mal condicionamento deixa a convergência lenta; padronizar ajuda\\
quando usar & poucos parâmetros, modelo linear nos parâmetros & muitos dados ou parâmetros, modelos não lineares, redes neurais\\
\bottomrule
\end{tabular}
\end{center}

\begin{observacao}
Os dois métodos minimizam a mesma função. Para modelos lineares nos parâmetros e poucos coeficientes, o método
direto é o mais simples e preciso. O gradiente descendente mostra a sua força quando o sistema normal é grande
demais para ser montado, ou quando o modelo não é linear nos parâmetros (por exemplo, $a\,e^{bx}$ sem
linearizar, ou uma rede neural). Nesses casos não há sistema linear a resolver, mas o gradiente continua
disponível.
\end{observacao}

% ==================================================================
\section{Exercícios}

\begin{enumerate}[label=\arabic*.]
  \item Ajuste uma reta aos pontos $(1;\,2)$, $(2;\,3)$, $(3;\,5)$ e $(4;\,6)$ pelo sistema normal e calcule $D$.
  \item Para os mesmos pontos, faça duas iterações do gradiente descendente com $\varphi(x)=\alpha_1x+\alpha_2$,
        $\eta=0{,}05$, partindo de $\alpha_1=\alpha_2=0$.
  \item Mostre que, para a reta, a primeira equação do sistema normal equivale a $\bar f=\alpha_1+\alpha_2\bar x$:
        a reta de mínimos quadrados passa pelo ponto médio $(\bar x,\bar f)$.
  \item Escreva o sistema normal para $\varphi(x)=\alpha_1+\alpha_2\cos x$.
  \item Linearize $\varphi(x)=a\,x^{b}$ e $\varphi(x)=\dfrac{1}{a+bx}$.
  \item No exemplo da Seção~\ref{sec:gdreta}, os autovalores de $H$ são $0{,}219$ e $2{,}281$. Qual o maior passo
        que garante convergência? Que passo dá a convergência mais rápida?
  \item Explique por que, depois de padronizar as funções $g_1,\dots,g_4$, a derivada de $J$ em relação ao termo
        constante $\beta_5$ se anula quando $\beta_5=\bar f$, quaisquer que sejam $\beta_1,\dots,\beta_4$.
\end{enumerate}

\subsection*{Respostas}
\begin{enumerate}[label=\arabic*.,nosep]
  \item $\sum x=10$, $\sum x^2=30$, $\sum f=16$, $\sum xf=47$. O sistema é $4\alpha_1+10\alpha_2=16$,
        $10\alpha_1+30\alpha_2=47$, com solução $\alpha_1=0{,}5$, $\alpha_2=1{,}4$, logo $\varphi(x)=0{,}5+1{,}4x$ e $D=0{,}2$.
  \item Iteração 1: $e=(-2,-3,-5,-6)$, gradientes $-11{,}75$ e $-4$, logo $\alpha_1=0{,}5875$, $\alpha_2=0{,}2$.
        Iteração 2: $e=(-1{,}2125;\,-1{,}625;\,-3{,}0375;\,-3{,}45)$, gradientes $-6{,}84375$ e $-2{,}33125$, logo
        $\alpha_1\approx0{,}92969$, $\alpha_2\approx0{,}31656$.
  \item Dividindo $m\alpha_1+\alpha_2\sum x_k=\sum f(x_k)$ por $m$.
  \item $\begin{bmatrix} m & \sum\cos x_k\\ \sum\cos x_k & \sum\cos^2x_k\end{bmatrix}\begin{bmatrix}\alpha_1\\ \alpha_2\end{bmatrix}=\begin{bmatrix}\sum f(x_k)\\ \sum f(x_k)\cos x_k\end{bmatrix}$.
  \item $\ln\varphi=\ln a+b\ln x$ (reta nos pontos $(\ln x_k,\ln f(x_k))$); $1/\varphi=a+bx$ (reta nos pontos $(x_k,1/f(x_k))$).
  \item $\eta<2/2{,}281\approx0{,}877$. O passo ótimo é $2/(\lambda_{\min}+\lambda_{\max})=2/2{,}5=0{,}8$.
  \item $\partial J/\partial\beta_5=\frac1m\sum e_k=\frac1m\sum_k\big(\sum_j\beta_j\tilde g_j(x_k)+\beta_5-f(x_k)\big)
        =\sum_j\beta_j\overline{\tilde g_j}+\beta_5-\bar f=\beta_5-\bar f$, pois cada $\tilde g_j$ tem média zero.
\end{enumerate}

\subsection*{Referências}
\begin{itemize}[nosep]
  \item M.~A.~G.~Ruggiero e V.~L.~R.~Lopes, \emph{Cálculo Numérico: aspectos teóricos e computacionais}, 2.ª ed., Makron Books (cap.~6, ajuste de curvas).
  \item A.~Géron, \emph{Mãos à obra: aprendizado de máquina com Scikit-Learn, Keras e TensorFlow}, Alta Books (gradiente descendente e padronização).
\end{itemize}

% Capítulo de integração numérica (arquivos usados neste capítulo)
\def\arqTrap{RegraTrapezios_repetida.ggb}
\def\arqSimpA{RegraSimpson13_repetida.ggb}
\def\arqSimpB{RegraSimpson38_repetida.ggb}
\def\arqRomberg{MetodoRomberg.ggb}
\def\arqGauss{QuadraturaGaussLegendre.ggb}
\chapter{Integração numérica}\label{cap:integracao}

\begin{tcolorbox}[colback=gray!4,colframe=gray!50,title={Como usar os aplicativos deste capítulo},fonttitle=\bfseries]
Ao longo do capítulo, os quadros laranja \emph{Atividade no GeoGebra} trazem um link
\textbf{[abrir o aplicativo]} que abre o arquivo correspondente no GeoGebra Classic:
\begin{itemize}[nosep]
  \item \ggbabrir{\arqTrap}{Regra dos trapézios repetida};
  \item \ggbabrir{\arqSimpA}{Regra 1/3 de Simpson repetida};
  \item \ggbabrir{\arqSimpB}{Regra 3/8 de Simpson repetida};
  \item \ggbabrir{\arqRomberg}{Método de Romberg};
  \item \ggbabrir{\arqGauss}{Quadratura de Gauss-Legendre}.
\end{itemize}
Para que os links funcionem, mantenha este PDF e a pasta \texttt{ggb} juntos, na mesma
pasta, e abra o PDF num leitor que permita abrir arquivos externos (Adobe Acrobat Reader,
Okular, SumatraPDF, Foxit). Os leitores de PDF embutidos nos navegadores costumam bloquear
esse tipo de link; nesse caso, abra os arquivos da pasta \texttt{ggb} diretamente no GeoGebra.

Os cinco aplicativos começam com a integral $\int_0^1 e^{-x^2}\,dx$, que é o exemplo
desenvolvido no texto. Em todos eles é possível digitar outra função no campo
$f(x)=$, mudar o intervalo $[a,b]$ nos controles deslizantes e usar o botão
\emph{Enquadrar o gráfico} para ajustar a janela. A janela da direita mostra as contas e tem
barras de rolagem (vertical e horizontal). O valor de referência da integral é calculado
pelo próprio GeoGebra, com precisão bem maior que a dos métodos estudados.
\end{tcolorbox}

% ==================================================================
\section{O problema da integração numérica}

O Teorema Fundamental do Cálculo resolve a integral definida quando conhecemos uma
primitiva: $\int_a^b f(x)\,dx=F(b)-F(a)$, com $F'=f$. Na prática, esse caminho falha com
frequência:
\begin{itemize}[nosep]
  \item muitas funções simples não têm primitiva expressa por funções elementares, como
        $e^{-x^2}$, $\dfrac{\operatorname{sen}x}{x}$, $\sqrt{1+\cos^2x}$ (comprimento de arco de
        uma senoide) ou $\dfrac{1}{\ln x}$;
  \item às vezes a primitiva existe, mas é tão complicada que avaliá-la custa mais (e
        acumula mais erros de arredondamento) do que aproximar a integral diretamente;
  \item muitas vezes $f$ só é conhecida numa tabela de valores medidos (velocidades,
        vazões, temperaturas) e não há fórmula para integrar.
\end{itemize}

A ideia central deste capítulo vem do Capítulo~\ref{cap:interp}: \emph{trocamos $f$ por um
polinômio interpolador $p$ e integramos $p$}, o que é sempre fácil. Como $p$ é uma
combinação dos valores $f(x_i)$ nos nós, a integral aproximada também é:
\[
\int_a^b f(x)\,dx\;\approx\;\int_a^b p(x)\,dx\;=\;\sum_{i=0}^{n}w_i\,f(x_i).
\]
Uma expressão desse tipo chama-se \emph{fórmula de quadratura}; os $x_i$ são os \emph{nós}
e os $w_i$ são os \emph{pesos}. Os pesos dependem só dos nós, não de $f$.

\begin{definicao}
Uma fórmula de quadratura tem \emph{grau de exatidão} $m$ se é exata (erro zero) para
todos os polinômios de grau menor ou igual a $m$ e não é exata para algum polinômio de
grau $m+1$. Como a fórmula e a integral são lineares, basta testar
$1,x,x^2,\dots,x^{m+1}$.
\end{definicao}

Usaremos ao longo de todo o capítulo o mesmo exemplo,
\[
I=\int_0^1 e^{-x^2}\,dx = 0{,}746824132812\dots,
\]
uma integral sem primitiva elementar (ela aparece na distribuição normal da Estatística,
por meio da função erro $\operatorname{erf}$). O valor acima foi obtido com alta precisão e
serve para medir o erro de cada método.

Dividimos os métodos em dois grupos:
\begin{enumerate}[nosep]
  \item \textbf{Fórmulas de Newton-Cotes}: nós igualmente espaçados, $x_i=a+ih$. São as regras
        dos trapézios, $1/3$ de Simpson e $3/8$ de Simpson, quase sempre usadas na forma
        \emph{repetida} (composta). A extrapolação de Richardson e o método de Romberg
        melhoram os trapézios sem novos tipos de fórmula.
  \item \textbf{Quadratura de Gauss}: os nós também são escolhidos, de modo a obter o maior
        grau de exatidão possível com poucos pontos.
\end{enumerate}

% ==================================================================
\section{Regra dos trapézios}

\subsection{Regra simples}

Interpolando $f$ nos dois extremos, $x_0=a$ e $x_1=b$, obtemos a reta
$p_1(x)=f(x_0)+\dfrac{f(x_1)-f(x_0)}{h}(x-x_0)$, com $h=b-a$. A área sob essa reta é a área
de um trapézio:
\[
\int_{x_0}^{x_1}f(x)\,dx\approx T=\frac{h}{2}\bigl[f(x_0)+f(x_1)\bigr].
\]

\begin{teorema}[Erro da regra dos trapézios simples]\label{teo:errotrap}
Se $f''$ é contínua em $[x_0,x_1]$ e $h=x_1-x_0$, existe $\eta\in(x_0,x_1)$ tal que
\[
\int_{x_0}^{x_1}f(x)\,dx-\frac{h}{2}\bigl[f(x_0)+f(x_1)\bigr]=-\frac{h^3}{12}\,f''(\eta).
\]
\end{teorema}

\begin{proof}
Pelo Teorema~\ref{teo:erro} (erro da interpolação) com dois nós,
$f(x)-p_1(x)=\dfrac{f''(\xi(x))}{2}(x-x_0)(x-x_1)$. O produto $(x-x_0)(x-x_1)$ não muda de
sinal em $[x_0,x_1]$ (é $\le0$), então o Teorema do Valor Médio para integrais permite tirar
$f''$ de dentro da integral:
\[
\int_{x_0}^{x_1}\bigl(f-p_1\bigr)\,dx=\frac{f''(\eta)}{2}\int_{x_0}^{x_1}(x-x_0)(x-x_1)\,dx
=\frac{f''(\eta)}{2}\cdot\Bigl(-\frac{h^3}{6}\Bigr)=-\frac{h^3}{12}f''(\eta).
\qedhere
\]
\end{proof}

O erro depende de $f''$: a regra é exata para retas (grau de exatidão 1). Se $f$ é côncava
para cima ($f''>0$), o trapézio fica acima do gráfico e a regra superestima a integral; se
$f''<0$, subestima.

\subsection{Regra repetida}

Com um único trapézio, $h=b-a$ é grande e o erro também. A solução é dividir $[a,b]$ em $n$
subintervalos iguais, de comprimento $h=\dfrac{b-a}{n}$, com nós $x_i=a+ih$
($i=0,1,\dots,n$), e somar as áreas dos $n$ trapézios. Cada nó interno pertence a dois
trapézios, por isso aparece com peso 2:
\[
\boxed{\;T_n=\frac{h}{2}\Bigl[f(x_0)+2f(x_1)+2f(x_2)+\dots+2f(x_{n-1})+f(x_n)\Bigr]\;}
\]

\begin{teorema}[Erro da regra dos trapézios repetida]\label{teo:errotrapr}
Se $f''$ é contínua em $[a,b]$, existe $\eta\in(a,b)$ tal que
\[
\int_a^b f(x)\,dx-T_n=-\frac{(b-a)\,h^2}{12}\,f''(\eta).
\]
Em particular, se $|f''(x)|\le M_2$ em $[a,b]$, então
$\Bigl|\int_a^b f(x)\,dx-T_n\Bigr|\le\dfrac{(b-a)\,h^2}{12}\,M_2$.
\end{teorema}

\begin{proof}
Somando os erros dos $n$ trapézios (Teorema~\ref{teo:errotrap}), o erro total é
$-\dfrac{h^3}{12}\sum_{i=1}^{n}f''(\eta_i)$. A média
$\frac1n\sum f''(\eta_i)$ está entre o mínimo e o máximo de $f''$; pelo Teorema do Valor
Intermediário, ela é igual a $f''(\eta)$ para algum $\eta$. Assim, o erro é
$-\dfrac{h^3}{12}\,n\,f''(\eta)$, e $nh=b-a$.
\end{proof}

O erro é proporcional a $h^2$: dobrar $n$ (dividir $h$ por 2) divide o erro por
aproximadamente $4$.

\begin{exemplo}[: $\int_0^1 e^{-x^2}\,dx$ pela regra dos trapézios]
Com $n=4$, $h=0{,}25$ e os nós $0$; $0{,}25$; $0{,}5$; $0{,}75$; $1$:
\begin{center}
\renewcommand{\arraystretch}{1.3}
\begin{tabular}{c|ccccc}
$i$ & 0 & 1 & 2 & 3 & 4\\ \hline
$x_i$ & $0$ & $0{,}25$ & $0{,}5$ & $0{,}75$ & $1$\\
$f(x_i)$ & $1$ & $0{,}93941307$ & $0{,}77880078$ & $0{,}56978282$ & $0{,}36787944$\\
peso & 1 & 2 & 2 & 2 & 1
\end{tabular}
\end{center}
\begin{align*}
T_4&=\frac{0{,}25}{2}\bigl[1+2(0{,}93941307+0{,}77880078+0{,}56978282)+0{,}36787944\bigr]\\
&=\frac{0{,}25}{2}\cdot5{,}94387278=0{,}74298410.
\end{align*}
O erro real é $|I-T_4|\approx3{,}84\times10^{-3}$. Para a cota, $f''(x)=(4x^2-2)e^{-x^2}$
e $M_2=\max_{[0,1]}|f''(x)|=2$ (em $x=0$), logo
$\dfrac{(b-a)h^2}{12}M_2=\dfrac{0{,}0625}{12}\cdot2\approx1{,}04\times10^{-2}$, coerente com o
erro observado. Dobrando $n$ sucessivamente:
\begin{center}
\renewcommand{\arraystretch}{1.3}
\begin{tabular}{c|c|c|c}
$n$ & $T_n$ & $|I-T_n|$ & razão entre erros consecutivos\\ \hline
1 & $0{,}6839397206$ & $6{,}29\times10^{-2}$ & \\
2 & $0{,}7313702518$ & $1{,}55\times10^{-2}$ & $4{,}07$\\
4 & $0{,}7429840978$ & $3{,}84\times10^{-3}$ & $4{,}02$\\
8 & $0{,}7458656148$ & $9{,}59\times10^{-4}$ & $4{,}01$\\
16 & $0{,}7465845968$ & $2{,}40\times10^{-4}$ & $4{,}00$
\end{tabular}
\end{center}
A razão se aproxima de $4=2^2$, como previsto.
\end{exemplo}

\begin{figure}[h]
\centering
\begin{tikzpicture}
\begin{axis}[width=10cm,height=6cm,axis lines=middle,xmin=-0.1,xmax=1.2,ymin=0,ymax=1.15,
  xtick={0,0.25,0.5,0.75,1},xticklabels={$0$,$0{,}25$,$0{,}5$,$0{,}75$,$1$},
  ytick={0.5,1},yticklabels={$0{,}5$,$1$},grid=both,grid style={gray!20},xlabel={$x$},ylabel={$y$}]
\addplot[fill=laranja!25,draw=laranja,thick] coordinates
  {(0,0)(0,1)(0.25,0.9394)(0.5,0.7788)(0.75,0.5698)(1,0.3679)(1,0)} -- cycle;
\addplot[laranja,thick] coordinates {(0.25,0)(0.25,0.9394)};
\addplot[laranja,thick] coordinates {(0.5,0)(0.5,0.7788)};
\addplot[laranja,thick] coordinates {(0.75,0)(0.75,0.5698)};
\addplot[domain=-0.05:1.15,samples=100,very thick,azul] {exp(-x^2)};
\addplot[only marks,mark=*,mark size=2pt,vermelho] coordinates
  {(0,1)(0.25,0.9394)(0.5,0.7788)(0.75,0.5698)(1,0.3679)};
\end{axis}
\end{tikzpicture}
\caption{Regra dos trapézios repetida para $\int_0^1 e^{-x^2}\,dx$ com $n=4$.}
\label{fig:trap}
\end{figure}

\begin{atividade}{regra dos trapézios repetida}{\arqTrap}
Abra o aplicativo. À esquerda ficam o gráfico de $f$, os trapézios e os controles; à
direita, o passo $h$, a tabela com os nós, os valores $f(x_i)$, os pesos $c_i$ e os produtos
$c_i\,f(x_i)$, a soma $S$, o valor $T_n=\frac h2S$ e o quadro do erro.
\begin{enumerate}[nosep]
  \item Confira, com $n=4$, todos os números do exemplo acima.
  \item Faça $n=2,4,8,16$ e anote o erro $|I-T_n|$. Calcule as razões entre erros consecutivos.
  \item Marque \emph{Mostrar a área exata}. Em que partes do intervalo os trapézios ficam acima do
        gráfico? Relacione com o sinal de $f''(x)=(4x^2-2)e^{-x^2}$.
  \item Digite $f(x)=3x+1$ e depois $f(x)=x^2$. Para qual das duas a regra é exata? Por quê?
  \item Digite $f(x)=x^2$ em $[0,1]$ e confira que o erro real é \emph{igual} à cota. O que há de
        especial com $f''$ nesse caso?
\end{enumerate}
\end{atividade}

% ==================================================================
\section{Regra 1/3 de Simpson}

\subsection{Regra simples}

Agora interpolamos $f$ por uma parábola em três nós igualmente espaçados
$x_0$, $x_1=x_0+h$, $x_2=x_0+2h$. Pela forma de Lagrange,
$p_2(x)=f(x_0)L_0(x)+f(x_1)L_1(x)+f(x_2)L_2(x)$, e basta integrar as bases. Com a mudança de
variável $x=x_0+sh$ ($dx=h\,ds$, $s\in[0,2]$):
\begin{align*}
\int_{x_0}^{x_2}L_0(x)\,dx&=h\int_0^2\frac{(s-1)(s-2)}{2}\,ds=\frac{h}{3},&
\int_{x_0}^{x_2}L_1(x)\,dx&=h\int_0^2\frac{s(s-2)}{-1}\,ds=\frac{4h}{3},\\
\int_{x_0}^{x_2}L_2(x)\,dx&=h\int_0^2\frac{s(s-1)}{2}\,ds=\frac{h}{3}.
\end{align*}
Portanto
\[
\int_{x_0}^{x_2}f(x)\,dx\approx\frac{h}{3}\bigl[f(x_0)+4f(x_1)+f(x_2)\bigr].
\]
O nome vem do fator $h/3$. Os pesos somam $2h$, que é o comprimento do intervalo, como deve
ser para que a fórmula integre exatamente $f\equiv1$.

\begin{teorema}[Erro da regra 1/3 de Simpson simples]
Se $f^{(4)}$ é contínua em $[x_0,x_2]$, existe $\eta\in(x_0,x_2)$ tal que
\[
\int_{x_0}^{x_2}f(x)\,dx-\frac{h}{3}\bigl[f(x_0)+4f(x_1)+f(x_2)\bigr]=-\frac{h^5}{90}\,f^{(4)}(\eta).
\]
\end{teorema}

\begin{observacao}
Esperaríamos um erro com $f'''$, pois interpolamos por um polinômio de grau 2. Mas a
fórmula é exata também para $x^3$: por simetria, para $f(x)=(x-x_1)^3$ a integral em
$[x_0,x_2]$ é zero e a fórmula também dá $\frac h3[(-h)^3+0+h^3]=0$. Assim, a regra $1/3$
tem grau de exatidão 3, e o erro depende de $f^{(4)}$. Esse ``ganho grátis'' acontece em
todas as fórmulas de Newton-Cotes com número ímpar de nós.
\end{observacao}

\subsection{Regra repetida}

Dividimos $[a,b]$ em um número \textbf{par} $n$ de subintervalos, $h=\dfrac{b-a}{n}$, e
aplicamos a regra simples em cada par $[x_0,x_2]$, $[x_2,x_4]$, \dots, $[x_{n-2},x_n]$. Os
nós de índice ímpar são meios de pares (peso 4); os de índice par interno são extremos de
dois pares (peso 2):
\[
\boxed{\;S_n=\frac{h}{3}\Bigl[f(x_0)+4f(x_1)+2f(x_2)+4f(x_3)+\dots+2f(x_{n-2})+4f(x_{n-1})+f(x_n)\Bigr]\;}
\]

\begin{teorema}[Erro da regra 1/3 de Simpson repetida]
Se $f^{(4)}$ é contínua em $[a,b]$ e $n$ é par, existe $\eta\in(a,b)$ tal que
\[
\int_a^b f(x)\,dx-S_n=-\frac{(b-a)\,h^4}{180}\,f^{(4)}(\eta),
\qquad\text{logo}\qquad
\Bigl|\int_a^b f(x)\,dx-S_n\Bigr|\le\frac{(b-a)\,h^4}{180}\,M_4,
\]
com $M_4=\max_{[a,b]}|f^{(4)}(x)|$.
\end{teorema}

(A demonstração é a mesma do Teorema~\ref{teo:errotrapr}: são $n/2$ parcelas
$-\frac{h^5}{90}f^{(4)}(\eta_i)$ e $\frac n2h^5=\frac{(b-a)h^4}{2}$.)
Agora dobrar $n$ divide o erro por aproximadamente $2^4=16$.

\begin{exemplo}[: $\int_0^1 e^{-x^2}\,dx$ pela regra 1/3 de Simpson]
Com $n=4$ usamos os mesmos cinco valores do exemplo dos trapézios, com pesos
$1,4,2,4,1$:
\begin{align*}
S_4&=\frac{0{,}25}{3}\bigl[1+4(0{,}93941307)+2(0{,}77880078)+4(0{,}56978282)+0{,}36787944\bigr]\\
&=\frac{0{,}25}{3}\cdot8{,}96226456=0{,}74685538.
\end{align*}
O erro real é $|I-S_4|\approx3{,}12\times10^{-5}$, cerca de 120 vezes menor que o dos
trapézios com os \emph{mesmos} cinco valores de $f$. Para a cota,
$f^{(4)}(x)=(16x^4-48x^2+12)e^{-x^2}$ e $M_4=12$ (em $x=0$):
$\dfrac{1\cdot0{,}25^4}{180}\cdot12\approx2{,}60\times10^{-4}$.
\begin{center}
\renewcommand{\arraystretch}{1.3}
\begin{tabular}{c|c|c|c}
$n$ & $S_n$ & $|I-S_n|$ & razão\\ \hline
2 & $0{,}7471804289$ & $3{,}56\times10^{-4}$ & \\
4 & $0{,}7468553798$ & $3{,}12\times10^{-5}$ & $11{,}4$\\
8 & $0{,}7468261205$ & $1{,}99\times10^{-6}$ & $15{,}7$\\
16 & $0{,}7468242574$ & $1{,}25\times10^{-7}$ & $15{,}9$
\end{tabular}
\end{center}
\end{exemplo}

\begin{figure}[h]
\centering
\begin{tikzpicture}
\begin{axis}[width=10cm,height=6cm,axis lines=middle,xmin=-0.15,xmax=2.3,ymin=-0.12,ymax=1.15,
  xtick={0,1,2},xticklabels={$x_0$,$x_1$,$x_2$},ytick={0.5,1},yticklabels={$0{,}5$,$1$},
  grid=both,grid style={gray!20},legend style={at={(0.98,0.98)},anchor=north east,font=\small}]
\addplot[domain=0:2,samples=80,fill=laranja!25,draw=laranja,thick] {1-0.773400*x+0.141279*x^2} \closedcycle;
\addplot[domain=-0.1:2.25,samples=120,very thick,azul] {exp(-x^2)};
\legend{parábola $p_2$,$f(x)=e^{-x^2}$}
\addplot[only marks,mark=*,mark size=2pt,vermelho] coordinates {(0,1)(1,0.3679)(2,0.0183)};
\end{axis}
\end{tikzpicture}
\caption{Regra 1/3 de Simpson simples: a área sob a parábola que passa pelos três nós
substitui a área sob $f$. Aqui, $f(x)=e^{-x^2}$ em $[0,2]$ com $h=1$ (escolhido grande
para que a diferença entre as curvas fique visível).}
\label{fig:simpson}
\end{figure}

\begin{atividade}{regra 1/3 de Simpson repetida}{\arqSimpA}
O aplicativo desenha, em cada par de subintervalos, a região sob a parábola que passa pelos
três nós. À direita aparecem a tabela com os pesos $1,4,2,4,\dots,4,1$, a soma, $S_n$ e o
quadro do erro.
\begin{enumerate}[nosep]
  \item Confira os números do exemplo com $n=4$ e compare com a tabela dos trapézios (os
        valores $f(x_i)$ são os mesmos; só os pesos mudam).
  \item Faça $n=2,4,8,16$ e verifique que as razões entre os erros se aproximam de $16$.
  \item Digite $f(x)=x^3-2x+1$ em $[-1,2]$ com $n=2$. O erro é zero? Explique com o grau de
        exatidão. E com $f(x)=x^4$?
  \item Com $f(x)=e^{-x^2}$ em $[0,2]$ e $n=2$, compare o gráfico com a
        Figura~\ref{fig:simpson}. Aumente $n$ e observe as parábolas se confundirem com $f$.
\end{enumerate}
\end{atividade}

% ==================================================================
\section{Regra 3/8 de Simpson}

Interpolando $f$ por uma cúbica em quatro nós igualmente espaçados $x_0,\dots,x_3$ e
integrando as bases de Lagrange (com $x=x_0+sh$, $s\in[0,3]$), obtemos
\[
\int_{x_0}^{x_3}f(x)\,dx\approx\frac{3h}{8}\bigl[f(x_0)+3f(x_1)+3f(x_2)+f(x_3)\bigr],
\qquad\text{erro}=-\frac{3h^5}{80}\,f^{(4)}(\eta).
\]
O grau de exatidão também é 3: o quarto nó não aumenta o grau, apenas melhora a constante
do erro por trecho. Na forma repetida, com $n$ \textbf{múltiplo de 3}, os nós que separam dois
grupos têm peso 2 e os demais nós internos têm peso 3:
\[
\boxed{\;S^{3/8}_n=\frac{3h}{8}\Bigl[f(x_0)+3f(x_1)+3f(x_2)+2f(x_3)+3f(x_4)+\dots+3f(x_{n-1})+f(x_n)\Bigr]\;}
\]
com erro $-\dfrac{(b-a)\,h^4}{80}\,f^{(4)}(\eta)$.

Comparando as cotas, $\frac{1}{80}>\frac{1}{180}$: com o mesmo $h$, a regra $1/3$ costuma ser mais
precisa. A regra $3/8$ é útil quando o número de subintervalos é ímpar (por exemplo, numa
tabela de dados com $n=7$): aplica-se a $3/8$ em três subintervalos e a $1/3$ nos quatro
restantes.

\begin{exemplo}[: $\int_0^1 e^{-x^2}\,dx$ pela regra 3/8 de Simpson]
Com $n=3$, $h=\frac13$ e $f(0)=1$, $f(\frac13)=0{,}89483932$, $f(\frac23)=0{,}64118039$,
$f(1)=0{,}36787944$:
\begin{align*}
S^{3/8}_3&=\frac{3\cdot\frac13}{8}\bigl[1+3(0{,}89483932)+3(0{,}64118039)+0{,}36787944\bigr]\\
&=\frac18\cdot5{,}97593856=0{,}74699232,
\end{align*}
com erro $|I-S^{3/8}_3|\approx1{,}68\times10^{-4}$. Com $n=6$, $S^{3/8}_6=0{,}7468380575$
e o erro cai para $1{,}39\times10^{-5}$ (razão $\approx12$). Com o mesmo $n=6$, a regra $1/3$
dá erro $\approx 6{,}3\times10^{-6}$ (Exercício~\ref{ex:simpson6}).
\end{exemplo}

\begin{atividade}{regra 3/8 de Simpson repetida}{\arqSimpB}
Aqui cada grupo de três subintervalos é substituído pela cúbica que passa pelos quatro nós.
\begin{enumerate}[nosep]
  \item Confira o exemplo com $n=3$ e depois com $n=6$.
  \item Abra também o aplicativo da regra $1/3$ e use $n=6$ e $n=12$ nos dois. Qual regra
        tem o menor erro? A razão entre os erros é próxima de $\frac{180}{80}=2{,}25$?
  \item Teste $f(x)=x^3$ e $f(x)=x^4$ em $[0,3]$ com $n=3$. Confira que o grau de exatidão é 3.
\end{enumerate}
\end{atividade}

% ==================================================================
\section{Fórmulas de Newton-Cotes: resumo}

As três regras são casos da mesma construção: nós igualmente espaçados e integração do
polinômio interpolador. Na tabela, $h$ é o espaçamento e o erro é o da fórmula simples.

\begin{center}
\renewcommand{\arraystretch}{1.55}
\begin{tabular}{l|c|c|c}
fórmula (simples) & nós & grau de exatidão & erro\\ \hline
trapézios: $\frac h2(f_0+f_1)$ & 2 & 1 & $-\frac{h^3}{12}f''(\eta)$\\
$1/3$ de Simpson: $\frac h3(f_0+4f_1+f_2)$ & 3 & 3 & $-\frac{h^5}{90}f^{(4)}(\eta)$\\
$3/8$ de Simpson: $\frac{3h}{8}(f_0+3f_1+3f_2+f_3)$ & 4 & 3 & $-\frac{3h^5}{80}f^{(4)}(\eta)$\\
Boole: $\frac{2h}{45}(7f_0+32f_1+12f_2+32f_3+7f_4)$ & 5 & 5 & $-\frac{8h^7}{945}f^{(6)}(\eta)$
\end{tabular}
\end{center}

\begin{observacao}
Por que não usar um único polinômio de grau alto em todo o intervalo? Pelo mesmo motivo do
fenômeno de Runge (Figura~\ref{fig:runge}): com muitos nós igualmente espaçados o polinômio
interpolador pode oscilar, e, a partir de 9 nós, as fórmulas de Newton-Cotes passam a ter
pesos negativos, o que amplifica erros de arredondamento. Na prática, usamos fórmulas de
grau baixo \emph{repetidas} em muitos subintervalos.
\end{observacao}

% ==================================================================
\section{Quantos subintervalos usar?}

As cotas de erro permitem escolher $n$ \emph{antes} de calcular, desde que saibamos
limitar a derivada que aparece na fórmula.

\begin{exemplo}[: escolha de $n$ para $\int_1^2\frac{dx}{x}=\ln2$ com erro menor que $10^{-4}$]
Para $f(x)=\frac1x$ temos $f''(x)=\frac{2}{x^3}$ e $f^{(4)}(x)=\frac{24}{x^5}$, ambas
decrescentes em $[1,2]$; logo $M_2=2$ e $M_4=24$.

\textbf{Trapézios:}
$\dfrac{(2-1)h^2}{12}\cdot2\le10^{-4}\iff h^2\le6\times10^{-4}\iff h\le0{,}02449$. Como
$h=\frac1n$, precisamos de $n\ge40{,}8$, isto é, $n=41$ subintervalos (42 valores de $f$).

\textbf{Simpson 1/3:}
$\dfrac{(2-1)h^4}{180}\cdot24\le10^{-4}\iff h^4\le7{,}5\times10^{-4}\iff h\le0{,}1655$, isto é,
$n\ge6{,}04$; como $n$ deve ser par, $n=8$ (9 valores de $f$).

Os erros reais são $|\ln2-T_{41}|\approx3{,}7\times10^{-5}$ e
$|\ln2-S_8|\approx7{,}4\times10^{-6}$. As cotas garantem a tolerância, mas são pessimistas:
com $S_6$ o erro já seria $2{,}3\times10^{-5}$.
\end{exemplo}

Quando não conseguimos limitar a derivada (ou $f$ vem de uma tabela), uma alternativa é
comparar duas aproximações: se $T_n$ e $T_{2n}$ concordam em $k$ casas, temos uma boa ideia
da precisão. A próxima seção transforma essa comparação num método muito mais eficiente.

% ==================================================================
\section{Extrapolação de Richardson e método de Romberg}

\subsection{A ideia de Richardson}

Para $f$ suficientemente derivável, a fórmula de Euler-Maclaurin garante que o erro dos
trapézios tem uma expansão em potências \emph{pares} de $h$:
\[
T(h)=I+c_1h^2+c_2h^4+c_3h^6+\cdots,
\]
com $c_1,c_2,\dots$ que não dependem de $h$. Escrevendo a mesma expressão para $h/2$,
\[
T(h/2)=I+c_1\frac{h^2}{4}+c_2\frac{h^4}{16}+\cdots,
\]
e combinando as duas para eliminar o termo em $h^2$:
\[
\frac{4\,T(h/2)-T(h)}{3}=I-\frac{c_2}{4}h^4+\cdots=T(h/2)+\frac{T(h/2)-T(h)}{3}.
\]
O resultado tem erro $O(h^4)$, muito menor, sem nenhuma nova avaliação de $f$. Isso é a
\emph{extrapolação de Richardson}.

\begin{proposicao}\label{prop:richardson}
Para todo $n$, $\dfrac{4\,T_{2n}-T_n}{3}=S_{2n}$: um passo de Richardson aplicado aos
trapézios reproduz a regra $1/3$ de Simpson.
\end{proposicao}

\begin{proof}
Basta verificar num par de subintervalos $[x_0,x_2]$ de comprimento $2h$:
$T_1=h\,(f_0+f_2)$ e $T_2=\frac h2(f_0+2f_1+f_2)$. Então
$\frac{4T_2-T_1}{3}=\frac{h}{3}\bigl(2f_0+4f_1+2f_2-f_0-f_2\bigr)=\frac h3(f_0+4f_1+f_2)$.
Somando sobre os pares, obtemos $S_{2n}$.
\end{proof}

\subsection{O método de Romberg}

Repetindo a ideia, cada nova coluna elimina mais um termo da expansão ($h^4$, depois $h^6$,
\dots). Organizamos tudo numa tabela triangular.

\begin{tcolorbox}[colback=azul!3,colframe=azul!60,title={Algoritmo: método de Romberg},fonttitle=\bfseries]
\begin{enumerate}[nosep]
  \item Para $k=0,1,\dots,K$: calcule $R_{k,0}=T_{2^k}$ (trapézios com $2^k$ subintervalos).
        Cada $T_{2^k}$ reaproveita todos os valores de $f$ de $T_{2^{k-1}}$:
        $T_{2n}=\frac12T_n+h_{2n}\sum f(\text{novos pontos médios})$.
  \item Para $j=1,\dots,k$:
        \[
        R_{k,j}=R_{k,j-1}+\frac{R_{k,j-1}-R_{k-1,j-1}}{4^j-1}.
        \]
  \item Pare quando $|R_{k,k}-R_{k-1,k-1}|$ for menor que a tolerância. A melhor aproximação
        é $R_{K,K}$.
\end{enumerate}
\end{tcolorbox}

Pela Proposição~\ref{prop:richardson}, a coluna $R_{k,1}$ é a regra $1/3$ de Simpson com $2^k$
subintervalos; a coluna $R_{k,2}$ coincide com a regra de Boole.

\begin{exemplo}[: Romberg para $\int_0^1 e^{-x^2}\,dx$ com $K=4$]
\begin{center}
\renewcommand{\arraystretch}{1.3}
\small
\begin{tabular}{c|c|lllll}
$k$ & $n$ & $R_{k,0}=T_n$ & $R_{k,1}$ & $R_{k,2}$ & $R_{k,3}$ & $R_{k,4}$\\ \hline
0 & 1 & $0{,}6839397206$ & & & & \\
1 & 2 & $0{,}7313702518$ & $0{,}7471804289$ & & & \\
2 & 4 & $0{,}7429840978$ & $0{,}7468553798$ & $0{,}7468337098$ & & \\
3 & 8 & $0{,}7458656148$ & $0{,}7468261205$ & $0{,}7468241699$ & $0{,}7468240185$ & \\
4 & 16 & $0{,}7465845968$ & $0{,}7468242574$ & $0{,}7468241332$ & $0{,}7468241326$ & $0{,}7468241331$
\end{tabular}
\end{center}
Por exemplo, $R_{1,1}=0{,}7313702518+\frac{0{,}7313702518-0{,}6839397206}{3}=0{,}7471804289$.
Com os mesmos 17 valores de $f$ usados por $T_{16}$ (erro $2{,}4\times10^{-4}$), o valor
$R_{4,4}$ tem erro de apenas $2{,}8\times10^{-10}$.
\end{exemplo}

\begin{observacao}
Na linha $k=3$, $R_{3,3}$ (erro $1{,}1\times10^{-7}$) é um pouco pior que $R_{3,2}$ (erro
$3{,}7\times10^{-8}$). A expansão $c_1h^2+c_2h^4+\cdots$ só descreve bem o erro quando $h$ já
é pequeno; com $h=\frac18$ o termo seguinte ainda pesa. Por isso o critério de parada compara
linhas sucessivas da diagonal. Se $f$ não for suave (por exemplo, $\sqrt x$ perto de $0$), a
expansão nem vale, e o ganho do método desaparece (Exercício~\ref{ex:raiz}).
\end{observacao}

\begin{atividade}{método de Romberg}{\arqRomberg}
O controle $K$ escolhe quantas linhas da tabela aparecem; o controle $m$ escolhe quais
trapézios ($2^m$ subintervalos) são desenhados.
\begin{enumerate}[nosep]
  \item Com $K=4$, confira a tabela do exemplo e o erro de $R_{4,4}$.
  \item Compare a coluna $R_{k,1}$ com os valores $S_n$ obtidos no aplicativo da regra $1/3$.
  \item Com $K=5$, quantas avaliações de $f$ foram usadas? Qual é o erro de $R_{5,5}$?
  \item Digite $f(x)=\sqrt{x}$ em $[0,1]$ e compare os erros dos trapézios e de Romberg com os
        de $e^{-x^2}$. O que mudou?
\end{enumerate}
\end{atividade}

% ==================================================================
\section{Quadratura de Gauss-Legendre}

\subsection{A ideia}

Nas fórmulas de Newton-Cotes os nós são fixados de antemão e só os pesos são escolhidos;
com $n$ nós, o grau de exatidão é $n-1$ ou $n$. Se escolhermos também os nós, teremos $2n$
parâmetros livres ($x_1,\dots,x_n$ e $w_1,\dots,w_n$) e poderemos exigir exatidão para
$1,x,\dots,x^{2n-1}$, ou seja, grau de exatidão $2n-1$.

Trabalhamos no intervalo padrão $[-1,1]$. Para $n=2$, queremos
$\int_{-1}^{1}g(t)\,dt\approx w_1g(t_1)+w_2g(t_2)$ exata para $g=1,t,t^2,t^3$:
\[
\begin{cases}
w_1+w_2=2\\
w_1t_1+w_2t_2=0\\
w_1t_1^2+w_2t_2^2=\frac23\\
w_1t_1^3+w_2t_2^3=0
\end{cases}
\]
Por simetria, procuramos $t_2=-t_1$ e $w_1=w_2$. A primeira equação dá $w_1=w_2=1$ e a
terceira, $2t_1^2=\frac23$, ou seja, $t_{1,2}=\mp\frac{1}{\sqrt3}$. Assim,
\[
\int_{-1}^{1}g(t)\,dt\approx g\Bigl(-\tfrac{1}{\sqrt3}\Bigr)+g\Bigl(\tfrac{1}{\sqrt3}\Bigr),
\]
exata para cúbicas com apenas \emph{duas} avaliações (a regra $1/3$ de Simpson precisa de três).

Para $n$ qualquer, os nós são as raízes do \emph{polinômio de Legendre} $P_n(t)$ (por isso
o nome Gauss-Legendre) e os pesos são positivos. A tabela traz os primeiros casos.
\begin{center}
\renewcommand{\arraystretch}{1.4}
\begin{tabular}{c|c|c}
$n$ & nós $t_i$ & pesos $w_i$\\ \hline
1 & $0$ & $2$\\
2 & $\pm\frac{1}{\sqrt3}\approx\pm0{,}5773502692$ & $1$\\
3 & $0$;\quad $\pm\sqrt{\frac35}\approx\pm0{,}7745966692$ & $\frac89$;\quad $\frac59$\\
4 & $\pm0{,}3399810436$;\quad $\pm0{,}8611363116$ & $0{,}6521451549$;\quad $0{,}3478548451$\\
5 & $0$;\quad $\pm0{,}5384693101$;\quad $\pm0{,}9061798459$ & $0{,}5688888889$;\quad $0{,}4786286705$;\quad $0{,}2369268851$
\end{tabular}
\end{center}

\subsection{Mudança de intervalo}

Para integrar em $[a,b]$, usamos a mudança de variável afim
$x=\dfrac{b-a}{2}\,t+\dfrac{a+b}{2}$, que leva $[-1,1]$ em $[a,b]$, com
$dx=\frac{b-a}{2}\,dt$:
\[
\boxed{\;\int_a^b f(x)\,dx\approx G_n=\frac{b-a}{2}\sum_{i=1}^{n}w_i\,f(x_i),\qquad
x_i=\frac{b-a}{2}\,t_i+\frac{a+b}{2}.\;}
\]

\begin{teorema}[Erro da quadratura de Gauss-Legendre]
Se $f^{(2n)}$ é contínua em $[a,b]$, existe $\eta\in(a,b)$ tal que
\[
\int_a^b f(x)\,dx-G_n=\frac{(b-a)^{2n+1}\,(n!)^4}{(2n+1)\,\bigl[(2n)!\bigr]^3}\,f^{(2n)}(\eta).
\]
\end{teorema}

\begin{exemplo}[: $\int_0^1 e^{-x^2}\,dx$ por Gauss-Legendre]
Aqui $\frac{b-a}{2}=\frac{a+b}{2}=\frac12$, logo $x_i=\frac{t_i+1}{2}$.

\textbf{Dois pontos:} $x_{1,2}=\frac12\mp\frac{1}{2\sqrt3}\approx0{,}21132487$ e $0{,}78867513$;
\[
G_2=\frac12\bigl[e^{-0{,}21132487^2}+e^{-0{,}78867513^2}\bigr]
=\frac12\bigl[0{,}95632430+0{,}53686508\bigr]=0{,}74659469,
\]
com erro $2{,}29\times10^{-4}$. Com o mesmo número de avaliações, o trapézio simples tem erro
$6{,}3\times10^{-2}$.

\textbf{Três pontos:} $x_2=\frac12$ e $x_{1,3}=\frac12\mp\frac12\sqrt{0{,}6}\approx0{,}11270167$ e $0{,}88729833$;
\[
G_3=\frac12\Bigl[\tfrac59\,f(0{,}11270167)+\tfrac89\,f(0{,}5)+\tfrac59\,f(0{,}88729833)\Bigr]=0{,}74681458,
\]
com erro $9{,}5\times10^{-6}$, cerca de 37 vezes menor que o de Simpson simples
($S_2$, também com três avaliações). Com cinco pontos, $G_5=0{,}7468241268$, com erro
$6\times10^{-9}$.
\end{exemplo}

\begin{atividade}{quadratura de Gauss-Legendre}{\arqGauss}
O aplicativo mostra os nós $t_i$ e os pesos $w_i$, os nós $x_i$ já transportados para $[a,b]$,
o valor $G_{np}$ e uma comparação com os trapézios e com Simpson. No gráfico, a curva
tracejada é o polinômio que passa pelos pontos $(x_i,f(x_i))$, e a área laranja é a sua
integral.
\begin{enumerate}[nosep]
  \item Confira os valores do exemplo para $np=2$ e $np=3$.
  \item Verifique que a área laranja é igual a $G_{np}$. Por que isso acontece? (Pense na
        fórmula como integral de um polinômio interpolador.)
  \item Com $np=2$, digite $f(x)=x^3$ e depois $f(x)=x^4$ em $[0,1]$. Em qual caso o erro é
        zero? Confira o erro do segundo caso com o teorema acima.
  \item Digite $f(x)=1/(1+25x^2)$ em $[-1,1]$ e aumente $np$ de 1 a 5. O polinômio tracejado
        oscila, mas o erro diminui (com oscilações). Para a quadratura de Gauss, sabe-se que
        $G_n\to I$ para toda $f$ contínua, pois os pesos são positivos; com 20 pontos o erro
        já seria $3{,}6\times10^{-4}$.
\end{enumerate}
\end{atividade}

\begin{observacao}
A quadratura de Gauss não usa os extremos $a$ e $b$. Isso permite integrar funções que não
estão definidas nos extremos, como $\int_0^1\frac{\operatorname{sen}x}{x}\,dx$ ou
$\int_0^1\frac{dx}{\sqrt x}$ (esta última converge devagar, mas a fórmula pode ser aplicada).
Por outro lado, os nós são números irracionais, o que torna o método inadequado para dados
tabelados em pontos igualmente espaçados.
\end{observacao}

% ==================================================================
\section{Comparação dos métodos}

A tabela reúne os resultados do exemplo $\int_0^1 e^{-x^2}\,dx$. O custo de um método de
integração numérica se mede, principalmente, pelo número de avaliações de $f$.

\begin{center}
\renewcommand{\arraystretch}{1.35}
\begin{tabular}{l|c|c|c}
método & avaliações de $f$ & aproximação & erro\\ \hline
trapézios, $n=4$ & 5 & $0{,}7429840978$ & $3{,}8\times10^{-3}$\\
trapézios, $n=16$ & 17 & $0{,}7465845968$ & $2{,}4\times10^{-4}$\\
Simpson $1/3$, $n=4$ & 5 & $0{,}7468553798$ & $3{,}1\times10^{-5}$\\
Simpson $3/8$, $n=6$ & 7 & $0{,}7468380575$ & $1{,}4\times10^{-5}$\\
Simpson $1/3$, $n=16$ & 17 & $0{,}7468242574$ & $1{,}2\times10^{-7}$\\
Romberg, $R_{4,4}$ & 17 & $0{,}7468241331$ & $2{,}8\times10^{-10}$\\
Gauss, 3 pontos & 3 & $0{,}7468145842$ & $9{,}5\times10^{-6}$\\
Gauss, 5 pontos & 5 & $0{,}7468241268$ & $6{,}0\times10^{-9}$
\end{tabular}
\end{center}

\begin{figure}[h]
\centering
\begin{tikzpicture}
\begin{loglogaxis}[width=11cm,height=7cm,xlabel={avaliações de $f$},ylabel={erro $|I-\text{aproximação}|$},
  grid=both,grid style={gray!20},xmin=0.8,xmax=80,ymin=1e-9,ymax=0.2,
  legend style={at={(0.98,0.98)},anchor=north east,font=\small}]
\addplot[thick,cL1,mark=*] coordinates {(2,6.2884e-02)(3,1.5454e-02)(5,3.8400e-03)(9,9.5852e-04)(17,2.3954e-04)(33,5.9878e-05)(65,1.4969e-05)};
\addplot[thick,cL2,mark=square*] coordinates {(3,3.5630e-04)(5,3.1247e-05)(9,1.9877e-06)(17,1.2462e-07)(33,7.7946e-09)};
\addplot[thick,cL3,mark=triangle*] coordinates {(1,3.1977e-02)(2,2.2944e-04)(3,9.5486e-06)(4,3.3532e-07)(5,6.0462e-09)};
\legend{trapézios repetida,Simpson $1/3$ repetida,Gauss-Legendre}
\end{loglogaxis}
\end{tikzpicture}
\caption{Erro em função do número de avaliações de $f$ para $\int_0^1 e^{-x^2}\,dx$, em
escala logarítmica. As retas dos trapézios e de Simpson têm inclinações $-2$ e $-4$
(erros proporcionais a $h^2$ e $h^4$).}
\label{fig:compara}
\end{figure}

\paragraph{Em resumo.}
\begin{itemize}[nosep]
  \item A \textbf{regra dos trapézios} é a mais simples e funciona com qualquer tabela de
        dados, mas converge devagar ($h^2$). É imbatível, porém, para funções periódicas
        integradas num período inteiro, caso em que a expansão de Euler-Maclaurin se anula.
  \item A \textbf{regra 1/3 de Simpson} custa o mesmo e converge com $h^4$; é a escolha padrão
        para dados igualmente espaçados com $n$ par. A \textbf{3/8} completa o caso de $n$ ímpar.
  \item \textbf{Romberg} aproveita as mesmas avaliações dos trapézios e atinge precisão alta
        quando $f$ é suave.
  \item \textbf{Gauss-Legendre} dá a maior precisão por avaliação para funções suaves dadas por
        fórmula, mas exige calcular $f$ em pontos irracionais.
  \item Para funções com singularidades ou mudanças bruscas, todos os métodos perdem
        desempenho; usam-se então métodos \emph{adaptativos}, que refinam a malha só onde é
        necessário (é o que fazem os programas profissionais e o comando \texttt{Integral} do
        GeoGebra).
\end{itemize}

% ==================================================================
\section{Exercícios}

\begin{enumerate}[label=\textbf{\arabic*.},leftmargin=*]
\item Calcule $\int_0^{\pi}\operatorname{sen}x\,dx$ pela regra dos trapézios e pela regra
  $1/3$ de Simpson, ambas com $n=4$. Compare com o valor exato $2$.

\item Verifique que a regra dos trapézios tem grau de exatidão 1 e que a regra $1/3$ de
  Simpson tem grau de exatidão 3, testando-as em $1,x,x^2,x^3,x^4$ no intervalo $[0,2]$ com
  a fórmula simples.

\item Quantos subintervalos são necessários para calcular $\int_0^1 e^{-x^2}\,dx$ com erro
  menor que $10^{-6}$ (a) pela regra dos trapézios; (b) pela regra $1/3$ de Simpson? Use
  $M_2=2$ e $M_4=12$.

\item\label{ex:simpson6} Calcule $S_6$ (regra $1/3$, $n=6$) para $\int_0^1 e^{-x^2}\,dx$ e
  compare com $S^{3/8}_6=0{,}7468380575$.

\item A velocidade de um carro (em m/s) foi registrada a cada 10 segundos:
  \begin{center}
  \begin{tabular}{c|ccccccc}
  $t$ (s) & 0 & 10 & 20 & 30 & 40 & 50 & 60\\ \hline
  $v$ (m/s) & 0 & 8{,}2 & 14{,}6 & 19{,}3 & 22{,}4 & 24{,}1 & 24{,}8
  \end{tabular}
  \end{center}
  Estime a distância percorrida, $\int_0^{60}v(t)\,dt$, pela regra dos trapézios e pelas
  regras $1/3$ e $3/8$ de Simpson.

\item Mostre que, ao integrar as bases de Lagrange de quatro nós igualmente espaçados em
  $[x_0,x_3]$, obtêm-se os pesos $\frac{3h}{8},\frac{9h}{8},\frac{9h}{8},\frac{3h}{8}$ da regra
  $3/8$. (Sugestão: use $x=x_0+sh$, $s\in[0,3]$.)

\item Construa a tabela de Romberg com $K=2$ para $\int_0^1\dfrac{4}{1+x^2}\,dx=\pi$.

\item Use Gauss-Legendre com 2 e com 3 pontos para calcular
  $\int_0^{\pi/2}\operatorname{sen}x\,dx=1$.

\item Deduza a fórmula de Gauss-Legendre com $n=1$ (um nó e um peso) exigindo exatidão para
  $1$ e $t$. Que regra conhecida ela é?

\item \textbf{(GeoGebra)} No aplicativo de Gauss-Legendre, use $np=2$, $f(x)=x^4$ e
  $[a,b]=[0,1]$. Calcule o erro pela fórmula do teorema e compare com o aplicativo.

\item\label{ex:raiz} \textbf{(GeoGebra)} No aplicativo de Romberg, digite $f(x)=\sqrt{x}$ em
  $[0,1]$ (valor exato $\frac23$) e use $K=5$. As razões entre erros consecutivos dos
  trapézios ainda são próximas de 4? O valor $R_{5,5}$ é tão bom quanto no exemplo com
  $e^{-x^2}$? Explique.

\item \textbf{(GeoGebra)} No aplicativo dos trapézios, digite \texttt{sin(2pi x)\^{}2}, isto é,
  $f(x)=\operatorname{sen}^2(2\pi x)$, em $[0,1]$ (valor exato $\frac12$). Compare os resultados
  com $n=2$, $3$, $4$ e $8$ e relacione com a observação sobre funções periódicas da seção de
  comparação.
\end{enumerate}

\subsection*{Respostas}
\begin{enumerate}[label=\textbf{\arabic*.},leftmargin=*,nosep]
\item $T_4\approx1{,}896119$ (erro $0{,}1039$); $S_4\approx2{,}004560$ (erro $0{,}0046$).
\item Trapézios: exata para $1$ e $x$ (dá $2$ e $2$); para $x^2$ dá $4\ne\frac83$.
  Simpson ($h=1$): $\frac13(f(0)+4f(1)+f(2))$ dá $2$, $2$, $\frac83$, $4$, exatos, e para $x^4$
  dá $\frac{20}{3}\ne\frac{32}{5}$.
\item (a) $\frac{h^2}{12}\cdot2\le10^{-6}\Rightarrow h\le0{,}002449\Rightarrow n=409$.
  (b) $\frac{h^4}{180}\cdot12\le10^{-6}\Rightarrow h\le0{,}06223\Rightarrow n\ge16{,}07$, logo $n=18$ (par).
\item $S_6\approx0{,}7468303915$, com erro $\approx6{,}3\times10^{-6}$, cerca de $2{,}2$ vezes menor
  que o da regra $3/8$ com o mesmo $n$ (compare com $\frac{180}{80}=2{,}25$).
\item Trapézios: $1010$ m; Simpson $1/3$: $\approx1017{,}3$ m; Simpson $3/8$: $\approx1017{,}4$ m.
\item Por exemplo, $\int_{x_0}^{x_3}L_0\,dx=h\int_0^3\frac{(s-1)(s-2)(s-3)}{-6}\,ds=\frac{3h}{8}$
  e $\int_{x_0}^{x_3}L_1\,dx=h\int_0^3\frac{s(s-2)(s-3)}{2}\,ds=\frac{9h}{8}$; os outros dois
  seguem por simetria.
\item $R_{0,0}=3$; $R_{1,0}=3{,}1$, $R_{1,1}=3{,}1333333$; $R_{2,0}=3{,}1311765$,
  $R_{2,1}=3{,}1415686$, $R_{2,2}=3{,}1421176$ (erro $5{,}3\times10^{-4}$).
\item $G_2\approx0{,}9984726$ (erro $1{,}5\times10^{-3}$); $G_3\approx1{,}0000081$
  (erro $8{,}1\times10^{-6}$).
\item $w_1=2$ e $w_1t_1=0$, logo $t_1=0$: $\int_a^bf\,dx\approx(b-a)\,f\bigl(\frac{a+b}{2}\bigr)$,
  a regra do ponto médio, com grau de exatidão 1.
\item $G_2=0{,}19\overline{4}$ e o valor exato é $0{,}2$; pelo teorema, o erro é
  $\frac{1\cdot(2!)^4}{5\,(4!)^3}\cdot24=\frac{1}{180}\approx0{,}005556$, exatamente o observado
  (pois $f^{(4)}\equiv24$ é constante).
\item Não: as razões dos trapézios ficam entre $2{,}6$ e $2{,}8$, aproximando-se de $2^{1{,}5}\approx2{,}83$
  (o erro é proporcional a $h^{1{,}5}$), e
  $|I-R_{5,5}|\approx3{,}8\times10^{-4}$, muito pior que no exemplo. Como $f'(x)=\frac{1}{2\sqrt x}$
  não é limitada perto de 0, a expansão de Euler-Maclaurin não vale e a extrapolação não
  elimina os termos certos.
\item Com $n=2$ todos os nós caem nos zeros de $f$ e $T_2=0$ (os nós ``não enxergam'' a
  oscilação); a partir de $n=3$ obtém-se $T_n=\frac12$, exato. Para funções periódicas
  integradas num período completo, os termos da expansão de Euler-Maclaurin se anulam, e a
  regra dos trapézios converge muito rapidamente, desde que haja nós suficientes por período.
\end{enumerate}

% Capítulo de soluções numéricas de EDOs (arquivos usados neste capítulo)
\def\arqEuler{MetodoEuler.ggb}
\def\arqRKdois{RungeKutta2.ggb}
\def\arqRKquatro{RungeKutta4.ggb}
\def\arqOrdem{OrdemMetodosEDO.ggb}
\def\arqSistemas{SistemasEDO_RK4.ggb}
\def\arqTiro{MetodoTiro.ggb}
\def\arqDifFin{DiferencasFinitasPVC.ggb}
\chapter[Soluções numéricas de EDOs]{Soluções numéricas de equações diferenciais ordinárias}\label{cap:edo}

\begin{tcolorbox}[colback=gray!4,colframe=gray!50,title={Como usar os aplicativos deste capítulo},fonttitle=\bfseries]
Ao longo do capítulo, os quadros laranja \emph{Atividade no GeoGebra} trazem um link
\textbf{[abrir o aplicativo]} que abre o arquivo correspondente no GeoGebra Classic:
\begin{itemize}[nosep]
  \item \ggbabrir{\arqEuler}{Método de Euler};
  \item \ggbabrir{\arqRKdois}{Runge-Kutta de 2ª ordem (Heun e ponto médio)};
  \item \ggbabrir{\arqRKquatro}{Runge-Kutta de 4ª ordem};
  \item \ggbabrir{\arqOrdem}{Ordem dos métodos: Euler, Heun e RK4};
  \item \ggbabrir{\arqSistemas}{Sistemas de EDOs com RK4};
  \item \ggbabrir{\arqTiro}{Método do tiro};
  \item \ggbabrir{\arqDifFin}{Diferenças finitas para problemas de contorno}.
\end{itemize}
Para que os links funcionem, mantenha este PDF e a pasta \texttt{ggb} juntos, na mesma
pasta, e abra o PDF num leitor que permita abrir arquivos externos (Adobe Acrobat Reader,
Okular, SumatraPDF, Foxit). Os leitores de PDF embutidos nos navegadores costumam bloquear
esse tipo de link; nesse caso, abra os arquivos da pasta \texttt{ggb} diretamente no GeoGebra.

Os quatro primeiros aplicativos começam com o problema $y'=y-x^2+1$, $y(0)=0{,}5$, que é o
exemplo desenvolvido no texto. Em todos é possível digitar outra equação, mudar a condição
inicial, o passo $h$ e o número de passos $n$, mostrar o campo de direções e usar o botão
\emph{Enquadrar o gráfico}. A janela da direita mostra as contas e tem barras de rolagem
(vertical e horizontal). A curva verde de referência é calculada pelo próprio aplicativo com
RK4 e um passo 40 vezes menor, o que basta para medir o erro dos métodos estudados.
\end{tcolorbox}

% ==================================================================
\section{Problemas de valor inicial}

Uma \emph{equação diferencial ordinária} (EDO) de primeira ordem relaciona uma função
desconhecida $y(x)$ com a sua derivada: $y'=f(x,y)$. Ela descreve como uma grandeza
varia: o crescimento de uma população, o resfriamento de um corpo, a carga de um capacitor, a
concentração de um medicamento no sangue. Uma EDO tem infinitas soluções; para escolher uma,
fixamos o valor de $y$ num ponto.

\begin{definicao}
Um \emph{problema de valor inicial} (PVI) de primeira ordem é
\[
y'=f(x,y),\qquad y(x_0)=y_0,
\]
e uma solução no intervalo $[x_0,b]$ é uma função derivável $y(x)$ que satisfaz a equação em
$[x_0,b]$ e a condição inicial.
\end{definicao}

\begin{teorema}[Existência e unicidade]\label{teo:picard}
Se $f$ é contínua em $D=\{(x,y):x_0\le x\le b,\ y\in\mathbb R\}$ e satisfaz a
\emph{condição de Lipschitz} em $y$,
\[
|f(x,y_1)-f(x,y_2)|\le L\,|y_1-y_2|\qquad\text{para todos }(x,y_1),(x,y_2)\in D,
\]
então o PVI tem uma única solução em $[x_0,b]$. A condição de Lipschitz vale, por exemplo,
quando $\partial f/\partial y$ é contínua e limitada em $D$, com
$L=\max|\partial f/\partial y|$.
\end{teorema}

Os métodos analíticos (variáveis separáveis, fator integrante, coeficientes constantes) só
resolvem uma pequena parte das EDOs que aparecem nas aplicações. Um método numérico não
produz uma fórmula para $y(x)$: ele produz \emph{aproximações} $y_i\approx y(x_i)$ numa malha
de pontos $x_i=x_0+ih$, $i=0,1,\dots,n$, em que $h$ é o \emph{passo}.

\paragraph{Interpretação geométrica: o campo de direções.} A equação $y'=f(x,y)$ diz qual é a
inclinação da solução que passa por cada ponto $(x,y)$ do plano. Desenhando pequenos segmentos
com essas inclinações, obtemos o \emph{campo de direções}; as soluções são as curvas que
seguem o campo (Figura~\ref{fig:campo}). Todos os métodos deste capítulo constroem a solução
aproximada andando passo a passo ``ao longo do campo''.

Usaremos ao longo do capítulo o PVI
\[
y'=y-x^2+1,\qquad y(0)=0{,}5,\qquad x\in[0,2],
\]
cuja solução exata, $y(x)=(x+1)^2-\tfrac12e^x$, nos permite medir os erros. Em particular,
$y(2)=9-\tfrac12e^2=5{,}3054719505\dots$

% ==================================================================
\section{Método de Euler}

\subsection{Dedução}

Pela fórmula de Taylor em torno de $x_i$, com $x_{i+1}=x_i+h$,
\[
y(x_{i+1})=y(x_i)+h\,y'(x_i)+\frac{h^2}{2}y''(\xi_i)
=y(x_i)+h\,f\bigl(x_i,y(x_i)\bigr)+\frac{h^2}{2}y''(\xi_i).
\]
Desprezando o termo em $h^2$ e trocando $y(x_i)$ pela aproximação $y_i$, obtemos o
\emph{método de Euler}:
\[
\boxed{\;y_{i+1}=y_i+h\,f(x_i,y_i),\qquad x_{i+1}=x_i+h.\;}
\]
Geometricamente, em cada passo seguimos a reta tangente à solução que passa por
$(x_i,y_i)$, durante um intervalo de comprimento $h$.

\begin{exemplo}[: o PVI do capítulo pelo método de Euler]
Com $h=0{,}2$ e $f(x,y)=y-x^2+1$:
\begin{align*}
y_1&=0{,}5+0{,}2\,(0{,}5-0+1)=0{,}5+0{,}2\cdot1{,}5=0{,}8,\\
y_2&=0{,}8+0{,}2\,(0{,}8-0{,}04+1)=0{,}8+0{,}2\cdot1{,}76=1{,}152,\\
y_3&=1{,}152+0{,}2\,(1{,}152-0{,}16+1)=1{,}152+0{,}2\cdot1{,}992=1{,}5504.
\end{align*}
Continuando até $x_{10}=2$:
\begin{center}
\renewcommand{\arraystretch}{1.25}
\begin{tabular}{c|c|c|c|c}
$i$ & $x_i$ & $y_i$ & $y(x_i)$ & $|y(x_i)-y_i|$\\ \hline
0 & 0 & $0{,}5$ & $0{,}5$ & 0\\
1 & $0{,}2$ & $0{,}8$ & $0{,}8292986$ & $0{,}0293$\\
2 & $0{,}4$ & $1{,}152$ & $1{,}2140877$ & $0{,}0621$\\
5 & $1{,}0$ & $2{,}4581760$ & $2{,}6408591$ & $0{,}1827$\\
10 & $2{,}0$ & $4{,}8657845$ & $5{,}3054720$ & $0{,}4397$
\end{tabular}
\end{center}
O erro cresce ao longo do intervalo: cada passo comete um pequeno erro e, além disso, parte
de um ponto que já estava errado.
\end{exemplo}

\begin{figure}[h]
\centering
\begin{tikzpicture}
\begin{axis}[width=10.5cm,height=7cm,axis lines=middle,xmin=-0.15,xmax=2.2,ymin=-0.2,ymax=6.2,
  xtick={0,0.5,1,1.5,2},xticklabels={$0$,$0{,}5$,$1$,$1{,}5$,$2$},grid=both,grid style={gray!15},
  xlabel={$x$},ylabel={$y$},unbounded coords=jump,
  legend style={at={(0.02,0.98)},anchor=north west,font=\small}]
\addplot[gray!60,thin,forget plot] coordinates {(-0.055,-0.055)(0.055,0.055)(nan,nan)(-0.045,0.433)(0.045,0.567)(nan,nan)(-0.037,0.926)(0.037,1.074)(nan,nan)(-0.031,1.422)(0.031,1.578)(nan,nan)(-0.027,1.919)(0.027,2.081)(nan,nan)(-0.024,2.418)(0.024,2.582)(nan,nan)(-0.021,2.916)(0.021,3.084)(nan,nan)(-0.019,3.416)(0.019,3.584)(nan,nan)(-0.017,3.915)(0.017,4.085)(nan,nan)(-0.016,4.415)(0.016,4.585)(nan,nan)(-0.014,4.914)(0.014,5.086)(nan,nan)(-0.013,5.414)(0.013,5.586)(nan,nan)(-0.012,5.914)(0.012,6.086)(nan,nan)(0.144,-0.053)(0.256,0.053)(nan,nan)(0.154,0.434)(0.246,0.566)(nan,nan)(0.162,0.926)(0.238,1.074)(nan,nan)(0.168,1.422)(0.232,1.578)(nan,nan)(0.173,1.919)(0.227,2.081)(nan,nan)(0.176,2.418)(0.224,2.582)(nan,nan)(0.179,2.917)(0.221,3.083)(nan,nan)(0.181,3.416)(0.219,3.584)(nan,nan)(0.183,3.915)(0.217,4.085)(nan,nan)(0.184,4.415)(0.216,4.585)(nan,nan)(0.186,4.914)(0.214,5.086)(nan,nan)(0.187,5.414)(0.213,5.586)(nan,nan)(0.188,5.914)(0.212,6.086)(nan,nan)(0.342,-0.049)(0.458,0.049)(nan,nan)(0.352,0.436)(0.448,0.564)(nan,nan)(0.361,0.928)(0.439,1.072)(nan,nan)(0.367,1.423)(0.433,1.577)(nan,nan)(0.372,1.920)(0.428,2.080)(nan,nan)(0.375,2.418)(0.425,2.582)(nan,nan)(0.378,2.917)(0.422,3.083)(nan,nan)(0.381,3.416)(0.419,3.584)(nan,nan)(0.382,3.915)(0.418,4.085)(nan,nan)(0.384,4.415)(0.416,4.585)(nan,nan)(0.385,4.914)(0.415,5.086)(nan,nan)(0.386,5.414)(0.414,5.586)(nan,nan)(0.387,5.914)(0.413,6.086)(nan,nan)(0.538,-0.040)(0.662,0.040)(nan,nan)(0.548,0.441)(0.652,0.559)(nan,nan)(0.558,0.930)(0.642,1.070)(nan,nan)(0.565,1.424)(0.635,1.576)(nan,nan)(0.570,1.921)(0.630,2.079)(nan,nan)(0.574,2.419)(0.626,2.581)(nan,nan)(0.577,2.917)(0.623,3.083)(nan,nan)(0.580,3.416)(0.620,3.584)(nan,nan)(0.582,3.916)(0.618,4.084)(nan,nan)(0.583,4.415)(0.617,4.585)(nan,nan)(0.585,4.915)(0.615,5.085)(nan,nan)(0.586,5.414)(0.614,5.586)(nan,nan)(0.587,5.914)(0.613,6.086)(nan,nan)(0.733,-0.024)(0.867,0.024)(nan,nan)(0.742,0.450)(0.858,0.550)(nan,nan)(0.753,0.936)(0.847,1.064)(nan,nan)(0.761,1.427)(0.839,1.573)(nan,nan)(0.767,1.923)(0.833,2.077)(nan,nan)(0.772,2.420)(0.828,2.580)(nan,nan)(0.776,2.918)(0.824,3.082)(nan,nan)(0.778,3.417)(0.822,3.583)(nan,nan)(0.781,3.916)(0.819,4.084)(nan,nan)(0.783,4.415)(0.817,4.585)(nan,nan)(0.784,4.915)(0.816,5.085)(nan,nan)(0.785,5.414)(0.815,5.586)(nan,nan)(0.787,5.914)(0.813,6.086)(nan,nan)(0.930,0.000)(1.070,0.000)(nan,nan)(0.935,0.468)(1.065,0.532)(nan,nan)(0.945,0.945)(1.055,1.055)(nan,nan)(0.955,1.433)(1.045,1.567)(nan,nan)(0.963,1.926)(1.037,2.074)(nan,nan)(0.969,2.422)(1.031,2.578)(nan,nan)(0.973,2.919)(1.027,3.081)(nan,nan)(0.976,3.418)(1.024,3.582)(nan,nan)(0.979,3.916)(1.021,4.084)(nan,nan)(0.981,4.416)(1.019,4.584)(nan,nan)(0.983,4.915)(1.017,5.085)(nan,nan)(0.984,5.415)(1.016,5.585)(nan,nan)(0.986,5.914)(1.014,6.086)(nan,nan)(1.134,0.029)(1.266,-0.029)(nan,nan)(1.130,0.496)(1.270,0.504)(nan,nan)(1.136,0.964)(1.264,1.036)(nan,nan)(1.147,1.443)(1.253,1.557)(nan,nan)(1.156,1.932)(1.244,2.068)(nan,nan)(1.164,2.425)(1.236,2.575)(nan,nan)(1.169,2.921)(1.231,3.079)(nan,nan)(1.174,3.419)(1.226,3.581)(nan,nan)(1.177,3.917)(1.223,4.083)(nan,nan)(1.179,4.416)(1.221,4.584)(nan,nan)(1.181,4.916)(1.219,5.084)(nan,nan)(1.183,5.415)(1.217,5.585)(nan,nan)(1.185,5.915)(1.215,6.085)(nan,nan)(1.344,0.053)(1.456,-0.053)(nan,nan)(1.334,0.530)(1.466,0.470)(nan,nan)(1.330,0.997)(1.470,1.003)(nan,nan)(1.336,1.465)(1.464,1.535)(nan,nan)(1.346,1.944)(1.454,2.056)(nan,nan)(1.356,2.432)(1.444,2.568)(nan,nan)(1.363,2.925)(1.437,3.075)(nan,nan)(1.369,3.421)(1.431,3.579)(nan,nan)(1.373,3.919)(1.427,4.081)(nan,nan)(1.377,4.417)(1.423,4.583)(nan,nan)(1.379,4.916)(1.421,5.084)(nan,nan)(1.381,5.416)(1.419,5.584)(nan,nan)(1.383,5.915)(1.417,6.085)(nan,nan)(1.556,0.068)(1.644,-0.068)(nan,nan)(1.547,0.557)(1.653,0.443)(nan,nan)(1.536,1.036)(1.664,0.964)(nan,nan)(1.530,1.504)(1.670,1.496)(nan,nan)(1.534,1.971)(1.666,2.029)(nan,nan)(1.544,2.447)(1.656,2.553)(nan,nan)(1.554,2.934)(1.646,3.066)(nan,nan)(1.562,3.426)(1.638,3.574)(nan,nan)(1.568,3.922)(1.632,4.078)(nan,nan)(1.573,4.419)(1.627,4.581)(nan,nan)(1.576,4.918)(1.624,5.082)(nan,nan)(1.579,5.417)(1.621,5.583)(nan,nan)(1.581,5.916)(1.619,6.084)(nan,nan)(1.766,0.076)(1.834,-0.076)(nan,nan)(1.759,0.571)(1.841,0.429)(nan,nan)(1.750,1.062)(1.850,0.938)(nan,nan)(1.740,1.545)(1.860,1.455)(nan,nan)(1.731,2.016)(1.869,1.984)(nan,nan)(1.731,2.482)(1.869,2.518)(nan,nan)(1.740,2.955)(1.860,3.045)(nan,nan)(1.751,3.438)(1.849,3.562)(nan,nan)(1.759,3.929)(1.841,4.071)(nan,nan)(1.766,4.423)(1.834,4.577)(nan,nan)(1.771,4.920)(1.829,5.080)(nan,nan)(1.775,5.418)(1.825,5.582)(nan,nan)(1.778,5.917)(1.822,6.083)(nan,nan)(1.973,0.081)(2.027,-0.081)(nan,nan)(1.969,0.578)(2.031,0.422)(nan,nan)(1.963,1.074)(2.037,0.926)(nan,nan)(1.955,1.567)(2.045,1.433)(nan,nan)(1.945,2.055)(2.055,1.945)(nan,nan)(1.935,2.532)(2.065,2.468)(nan,nan)(1.930,3.000)(2.070,3.000)(nan,nan)(1.935,3.468)(2.065,3.532)(nan,nan)(1.945,3.945)(2.055,4.055)(nan,nan)(1.955,4.433)(2.045,4.567)(nan,nan)(1.963,4.926)(2.037,5.074)(nan,nan)(1.969,5.422)(2.031,5.578)(nan,nan)(1.973,5.919)(2.027,6.081)(nan,nan)};
\addplot[domain=0:2,samples=80,very thick,verde] {(x+1)^2-0.5*exp(x)};
\addplot[thick,vermelho,mark=*,mark size=1.6pt] coordinates {(0.00,0.50000)(0.20,0.80000)(0.40,1.15200)(0.60,1.55040)(0.80,1.98848)(1.00,2.45818)(1.20,2.94981)(1.40,3.45177)(1.60,3.95013)(1.80,4.42815)(2.00,4.86578)};
\legend{solução exata,Euler com $h=0{,}2$}
\end{axis}
\end{tikzpicture}
\caption{Campo de direções de $y'=y-x^2+1$, a solução com $y(0)=0{,}5$ e a poligonal de
Euler com $h=0{,}2$.}
\label{fig:campo}
\end{figure}

\subsection{Erro local e erro global}

Dois tipos de erro aparecem em qualquer método de passo a passo:
\begin{itemize}[nosep]
  \item o \emph{erro local de truncamento} é o erro cometido num único passo, supondo que se
        parte do valor exato $y(x_i)$. No método de Euler, pela fórmula de Taylor acima, ele é
        $\frac{h^2}{2}y''(\xi_i)$, isto é, $O(h^2)$;
  \item o \emph{erro global} $e_i=y(x_i)-y_i$ acumula os erros locais dos $i$ passos e a
        propagação dos erros anteriores.
\end{itemize}
Como o número de passos até um ponto fixo $\bar x$ é $(\bar x-x_0)/h$, somar erros locais
$O(h^2)$ dá um erro global $O(h)$. O teorema a seguir torna isso preciso.

\begin{teorema}[Erro global do método de Euler]\label{teo:erroeuler}
Suponha $f$ nas condições do Teorema~\ref{teo:picard}, com constante de Lipschitz $L$, e
$|y''(x)|\le M$ em $[x_0,b]$. Então as aproximações de Euler satisfazem
\[
|y(x_i)-y_i|\le\frac{hM}{2L}\Bigl(e^{L(x_i-x_0)}-1\Bigr),\qquad i=0,1,\dots,n.
\]
\end{teorema}

\begin{proof}[Ideia da demonstração]
Subtraindo $y_{i+1}=y_i+hf(x_i,y_i)$ de
$y(x_{i+1})=y(x_i)+hf(x_i,y(x_i))+\frac{h^2}{2}y''(\xi_i)$ e usando a condição de
Lipschitz, $|e_{i+1}|\le(1+hL)|e_i|+\frac{h^2M}{2}$. Por indução,
$|e_i|\le\frac{hM}{2L}\bigl((1+hL)^i-1\bigr)$, e $(1+hL)^i\le e^{ihL}=e^{L(x_i-x_0)}$.
\end{proof}

A cota é proporcional a $h$: o método de Euler é de \emph{primeira ordem}. Dividir o passo por
2 divide o erro (aproximadamente) por 2; para ganhar uma casa decimal é preciso usar 10 vezes
mais passos. O fator exponencial mostra também que o erro pode crescer muito em intervalos
longos.

\begin{atividade}{método de Euler}{\arqEuler}
Abra o aplicativo. À esquerda ficam o campo de direções, a poligonal de Euler (vermelha) e a
solução de referência (verde); à direita, a fórmula, o erro em $x_n$ e a tabela com $x_i$,
$y_i$, $f(x_i,y_i)$, o valor de referência e o erro em cada ponto.
\begin{enumerate}[nosep]
  \item Confira os números do exemplo com $h=0{,}2$ e $n=10$.
  \item Use $h=0{,}1$ e $n=20$ (mesmo ponto final $x_n=2$). O erro em $x_n$ caiu
        aproximadamente pela metade?
  \item Marque \emph{Soluções exatas que passam por cada ponto}. Cada passo de Euler leva a
        aproximação para uma curva solução diferente. Use essa figura para explicar a
        diferença entre erro local e erro global.
  \item Digite $f(x,y)=-20y$, com $y_0=1$, $h=0{,}12$ e $n=10$. O que acontece? Repita com
        $h=0{,}04$. (Voltaremos a isso na Seção~\ref{sec:estab}.)
\end{enumerate}
\end{atividade}

% ==================================================================
\section{Métodos de Runge-Kutta}

\subsection{A ideia}

Para aumentar a ordem poderíamos usar mais termos da série de Taylor,
$y_{i+1}=y_i+hy'_i+\frac{h^2}{2}y''_i+\cdots$ (\emph{métodos de Taylor}), mas isso exige
calcular derivadas de $f$: no exemplo, $y''=\frac{d}{dx}(y-x^2+1)=y'-2x=y-x^2+1-2x$. Para
uma $f$ complicada, essas contas são trabalhosas e precisam ser refeitas para cada problema.

Os \emph{métodos de Runge-Kutta} obtêm a mesma ordem usando apenas valores de $f$: em cada
passo, $f$ é avaliada em alguns pontos de teste e as inclinações obtidas são combinadas com
pesos escolhidos para que a fórmula concorde com a série de Taylor até a ordem desejada.
Euler é o método de Runge-Kutta de ordem 1, com uma única inclinação.

\subsection{Métodos de segunda ordem}

Procuramos uma fórmula
\[
y_{i+1}=y_i+h\,(a_1k_1+a_2k_2),\qquad k_1=f(x_i,y_i),\quad k_2=f(x_i+p\,h,\ y_i+p\,h\,k_1).
\]
Expandindo $k_2$ em série de Taylor em duas variáveis e comparando com
$y(x_i+h)=y+hf+\frac{h^2}{2}(f_x+f_yf)+O(h^3)$, obtemos as condições
\[
a_1+a_2=1,\qquad a_2\,p=\tfrac12.
\]
Há infinitas escolhas; as duas mais usadas são:
\begin{itemize}[nosep]
  \item \textbf{método de Heun} (ou \emph{Euler melhorado}), com $a_1=a_2=\frac12$ e $p=1$:
        \[
        k_1=f(x_i,y_i),\quad k_2=f(x_i+h,\ y_i+hk_1),\quad y_{i+1}=y_i+\frac h2(k_1+k_2);
        \]
        primeiro damos um passo de Euler (\emph{previsão}) e depois corrigimos com a média das
        inclinações no início e no fim (\emph{correção}). Se $f$ não depende de $y$, isso é
        exatamente a regra dos trapézios do Capítulo~\ref{cap:integracao};
  \item \textbf{método do ponto médio}, com $a_1=0$, $a_2=1$ e $p=\frac12$:
        \[
        k_1=f(x_i,y_i),\quad k_2=f\bigl(x_i+\tfrac h2,\ y_i+\tfrac h2k_1\bigr),\quad y_{i+1}=y_i+h\,k_2;
        \]
        usamos a inclinação estimada no meio do passo.
\end{itemize}
Ambos têm erro local $O(h^3)$ e erro global $O(h^2)$, com duas avaliações de $f$ por passo.

\begin{exemplo}[: primeiro passo de Heun e do ponto médio]
Com $h=0{,}2$, $k_1=f(0;\,0{,}5)=1{,}5$ nos dois métodos.

\textbf{Heun:} o passo de Euler leva a $(0{,}2;\ 0{,}8)$ e
$k_2=f(0{,}2;\,0{,}8)=0{,}8-0{,}04+1=1{,}76$; logo
$y_1=0{,}5+0{,}1\,(1{,}5+1{,}76)=0{,}826$.

\textbf{Ponto médio:} o meio passo leva a $(0{,}1;\ 0{,}65)$ e
$k_2=f(0{,}1;\,0{,}65)=0{,}65-0{,}01+1=1{,}64$; logo $y_1=0{,}5+0{,}2\cdot1{,}64=0{,}828$.

O valor exato é $y(0{,}2)=0{,}8292986$. Prosseguindo até $x=2$, Heun dá
$y_{10}=5{,}2330546$ (erro $0{,}0724$) e o ponto médio dá $y_{10}=5{,}2903695$ (erro
$0{,}0151$), contra $0{,}4397$ de Euler.
\end{exemplo}

\begin{atividade}{Runge-Kutta de segunda ordem}{\arqRKdois}
O aplicativo mostra, para o passo escolhido no controle $kk$, a inclinação $k_1$ (laranja), o
ponto de teste ($\times$) e a inclinação $k_2$ (roxa). A caixa \emph{Usar o ponto médio}
troca Heun pelo método do ponto médio.
\begin{enumerate}[nosep]
  \item Confira o primeiro passo do exemplo nos dois métodos ($kk=0$).
  \item Marque \emph{Comparar com Euler} e compare as duas poligonais.
  \item Com Heun, use $h=0{,}2$, $n=10$ e depois $h=0{,}1$, $n=20$. O erro em $x_n$ fica
        cerca de 4 vezes menor? Por quê?
  \item Digite $f(x,y)=\cos(x)$, com $x_0=0$ e $y_0=0$. Que regra de integração numérica o
        método de Heun está aplicando? E o método do ponto médio?
\end{enumerate}
\end{atividade}

\subsection{O método de Runge-Kutta de quarta ordem}

O método mais usado na prática combina quatro inclinações por passo:
\[
\boxed{\;
\begin{aligned}
k_1&=f(x_i,\,y_i), & k_2&=f\bigl(x_i+\tfrac h2,\ y_i+\tfrac h2k_1\bigr),\\
k_3&=f\bigl(x_i+\tfrac h2,\ y_i+\tfrac h2k_2\bigr), & k_4&=f(x_i+h,\ y_i+hk_3),\\
y_{i+1}&=y_i+\frac h6\,(k_1+2k_2+2k_3+k_4). & &
\end{aligned}\;}
\]
A inclinação usada é uma média ponderada: uma no início do passo, duas no meio (pesos maiores)
e uma no fim. Quando $f$ só depende de $x$, $k_2=k_3=f(x_i+\frac h2)$ e a fórmula fica
$\frac h6\bigl[f(x_i)+4f(x_i+\frac h2)+f(x_i+h)\bigr]$: é a regra $1/3$ de Simpson no intervalo
$[x_i,x_{i+1}]$. O RK4 tem erro local $O(h^5)$ e erro global $O(h^4)$.

\begin{exemplo}[: o PVI do capítulo pelo RK4]
Primeiro passo, com $h=0{,}2$:
\begin{align*}
k_1&=f(0;\ 0{,}5)=1{,}5,\\
k_2&=f(0{,}1;\ 0{,}5+0{,}1\cdot1{,}5)=f(0{,}1;\ 0{,}65)=1{,}64,\\
k_3&=f(0{,}1;\ 0{,}5+0{,}1\cdot1{,}64)=f(0{,}1;\ 0{,}664)=1{,}654,\\
k_4&=f(0{,}2;\ 0{,}5+0{,}2\cdot1{,}654)=f(0{,}2;\ 0{,}8308)=1{,}7908,\\
y_1&=0{,}5+\frac{0{,}2}{6}\,(1{,}5+2\cdot1{,}64+2\cdot1{,}654+1{,}7908)=0{,}8292933,
\end{align*}
com erro $5{,}3\times10^{-6}$. Ao final, em $x=2$, $y_{10}=5{,}3053630$, com erro
$1{,}09\times10^{-4}$, cerca de 4000 vezes menor que o de Euler com o mesmo passo.
\end{exemplo}

\begin{atividade}{Runge-Kutta de quarta ordem}{\arqRKquatro}
Para o passo escolhido em $kk$, o gráfico mostra as quatro inclinações como pequenos segmentos
desenhados nos pontos em que foram calculadas: $k_1$ no ponto atual, $k_2$ e $k_3$ no meio do
passo e $k_4$ no fim. A tabela traz $k_1,\dots,k_4$ e o erro de cada $y_i$.
\begin{enumerate}[nosep]
  \item Confira os números do primeiro passo ($kk=0$) e o valor final.
  \item Compare os erros com $h=0{,}2$ ($n=10$) e $h=0{,}1$ ($n=20$). Qual é a razão?
  \item Use $f(x,y)=\cos(x)$, $x_0=0$, $y_0=0$, $h=0{,}5$ e $n=4$. Compare $y_4$ com
        $\operatorname{sen}2=\int_0^2\cos x\,dx$ e com a regra $1/3$ de Simpson com 8
        subintervalos do Capítulo~\ref{cap:integracao}.
\end{enumerate}
\end{atividade}

% ==================================================================
\section{Ordem, custo e estabilidade}\label{sec:estab}

\subsection{Ordem e custo}

Dizemos que um método tem \emph{ordem $p$} se o erro global em um ponto fixo é $O(h^p)$. Na
prática, isso significa que, ao dividir $h$ por 2, o erro fica aproximadamente $2^p$ vezes
menor. A tabela mostra os erros em $x=2$ para o problema do capítulo.

\begin{center}
\renewcommand{\arraystretch}{1.3}
\begin{tabular}{c|c|cc|cc|cc}
$n$ & $h$ & erro Euler & razão & erro Heun & razão & erro RK4 & razão\\ \hline
5 & $0{,}4$ & $7{,}46\times10^{-1}$ & & $2{,}63\times10^{-1}$ & & $1{,}64\times10^{-3}$ & \\
10 & $0{,}2$ & $4{,}40\times10^{-1}$ & $1{,}70$ & $7{,}24\times10^{-2}$ & $3{,}63$ & $1{,}09\times10^{-4}$ & $15{,}1$\\
20 & $0{,}1$ & $2{,}42\times10^{-1}$ & $1{,}82$ & $1{,}89\times10^{-2}$ & $3{,}83$ & $6{,}99\times10^{-6}$ & $15{,}6$\\
40 & $0{,}05$ & $1{,}27\times10^{-1}$ & $1{,}90$ & $4{,}82\times10^{-3}$ & $3{,}92$ & $4{,}42\times10^{-7}$ & $15{,}8$\\
80 & $0{,}025$ & $6{,}55\times10^{-2}$ & $1{,}95$ & $1{,}22\times10^{-3}$ & $3{,}96$ & $2{,}78\times10^{-8}$ & $15{,}9$
\end{tabular}
\end{center}

As razões se aproximam de $2$, $4$ e $16$, confirmando as ordens 1, 2 e 4. O custo de um
passo é medido pelo número de avaliações de $f$: 1 para Euler, 2 para Heun e 4 para RK4.
Mesmo assim o RK4 compensa: com $n=10$ (40 avaliações) ele é muito mais preciso que Euler com
$n=80$ (80 avaliações).

\begin{figure}[h]
\centering
\begin{tikzpicture}
\begin{loglogaxis}[width=10cm,height=6.5cm,xlabel={passo $h$},ylabel={erro em $x=2$},
  grid=both,grid style={gray!20},x dir=reverse,xtick={0.025,0.05,0.1,0.2,0.4},xticklabels={$0{,}025$,$0{,}05$,$0{,}1$,$0{,}2$,$0{,}4$},
  legend style={at={(0.02,0.02)},anchor=south west,font=\small}]
\addplot[thick,cL1,mark=*] coordinates {(0.4,0.746)(0.2,0.440)(0.1,0.242)(0.05,0.127)(0.025,0.0655)};
\addplot[thick,cL3,mark=square*] coordinates {(0.4,0.263)(0.2,0.0724)(0.1,0.0189)(0.05,0.00482)(0.025,0.00122)};
\addplot[thick,vermelho,mark=triangle*] coordinates {(0.4,1.64e-3)(0.2,1.09e-4)(0.1,6.99e-6)(0.05,4.42e-7)(0.025,2.78e-8)};
\legend{Euler (ordem 1),Heun (ordem 2),RK4 (ordem 4)}
\end{loglogaxis}
\end{tikzpicture}
\caption{Erro em $x=2$ em função do passo, em escala logarítmica. As inclinações das retas
são as ordens dos métodos.}
\label{fig:ordem}
\end{figure}

\begin{atividade}{ordem dos métodos}{\arqOrdem}
O aplicativo resolve o mesmo PVI por Euler, Heun e RK4 com $n_0$, $2n_0$, $4n_0$ e $8n_0$
passos e mostra os erros e as razões entre erros consecutivos.
\begin{enumerate}[nosep]
  \item Com $n_0=5$, confira a tabela acima.
  \item Digite outro PVI, por exemplo $f(x,y)=-2xy$ com $y_0=1$ e $x_f=1$ (solução $e^{-x^2}$),
        e verifique que as razões continuam próximas de 2, 4 e 16.
  \item Compare o erro do RK4 com $n_0$ passos e o de Euler com $8n_0$ passos (o dobro do custo
        do RK4). Qual método compensa?
\end{enumerate}
\end{atividade}

\subsection{Estabilidade}

Um método convergente pode, mesmo assim, dar resultados absurdos se o passo for grande demais.
O exemplo-modelo é
\[
y'=\lambda y,\qquad y(0)=1,\qquad \lambda<0,
\]
cuja solução $y=e^{\lambda x}$ decai para zero. Pelo método de Euler,
$y_{i+1}=y_i+h\lambda y_i=(1+h\lambda)\,y_i$, logo $y_i=(1+h\lambda)^i$. A aproximação só
decai se $|1+h\lambda|<1$, isto é,
\[
h<\frac{2}{|\lambda|}.
\]
Com $\lambda=-20$ e $h=0{,}12$, o fator é $1+h\lambda=-1{,}4$ e as aproximações oscilam com
amplitude crescente: $-1{,}4$; $1{,}96$; $-2{,}744$; $3{,}84$; $\dots$ Com $h=0{,}04$, o fator é $0{,}2$
e elas decaem corretamente. Esse fenômeno se chama \emph{instabilidade numérica}: não é um
erro de arredondamento, é a própria fórmula que amplifica as perturbações.

Os métodos de Runge-Kutta também têm um limite de estabilidade (para o RK4, $h|\lambda|\lesssim2{,}78$
quando $\lambda$ é real e negativo). Problemas em que convivem escalas de tempo muito diferentes
(variações rápidas que somem logo e variações lentas que interessam) são chamados
\emph{rígidos} (\emph{stiff}); para eles usam-se métodos implícitos, que não têm essa
restrição. Esse assunto vai além deste livro.

% ==================================================================
\section{Sistemas e equações de ordem superior}

\subsection{Sistemas de EDOs}

Muitos modelos envolvem várias grandezas que interagem. Um sistema de duas EDOs de primeira
ordem
\[
\begin{cases}u'=F(x,u,v),& u(x_0)=u_0,\\ v'=G(x,u,v),& v(x_0)=v_0,\end{cases}
\]
pode ser escrito na forma vetorial $\mathbf Y'=\mathbf F(x,\mathbf Y)$, com $\mathbf Y=(u,v)$ e
$\mathbf F=(F,G)$. Todos os métodos anteriores valem \emph{sem mudança}, desde que as contas
sejam feitas com vetores. Para o RK4:
\begin{align*}
\mathbf K_1&=\mathbf F(x_i,\mathbf Y_i), &
\mathbf K_2&=\mathbf F\bigl(x_i+\tfrac h2,\ \mathbf Y_i+\tfrac h2\mathbf K_1\bigr),\\
\mathbf K_3&=\mathbf F\bigl(x_i+\tfrac h2,\ \mathbf Y_i+\tfrac h2\mathbf K_2\bigr), &
\mathbf K_4&=\mathbf F(x_i+h,\ \mathbf Y_i+h\mathbf K_3),
\end{align*}
\[
\mathbf Y_{i+1}=\mathbf Y_i+\frac h6\,(\mathbf K_1+2\mathbf K_2+2\mathbf K_3+\mathbf K_4).
\]
Cuidado: $\mathbf K_2$ só pode ser calculado depois de \emph{todas} as componentes de
$\mathbf K_1$, e assim por diante.

\subsection{Equações de ordem superior}

Uma EDO de segunda ordem $y''=g(x,y,y')$, com $y(x_0)=y_0$ e $y'(x_0)=z_0$, vira um sistema de
primeira ordem com a troca $u=y$, $v=y'$:
\[
u'=v,\qquad v'=g(x,u,v),\qquad u(x_0)=y_0,\quad v(x_0)=z_0.
\]
Por exemplo, o pêndulo simples $\theta''=-\frac gL\operatorname{sen}\theta$ vira
$u'=v$, $v'=-\frac gL\operatorname{sen}u$. O mesmo vale para ordem $m$: introduzimos $y,y',\dots,y^{(m-1)}$
como incógnitas.

\begin{exemplo}[: o modelo predador-presa de Lotka-Volterra]
Seja $u(t)$ uma população de presas e $v(t)$ uma população de predadores (em unidades
convenientes), com
\[
u'=1{,}1\,u-0{,}4\,uv,\qquad v'=0{,}1\,uv-0{,}4\,v,\qquad u(0)=10,\quad v(0)=5.
\]
Sem predadores, as presas cresceriam exponencialmente; sem presas, os predadores morreriam.
Pode-se mostrar que as soluções são \emph{periódicas}: no plano $(u,v)$ (o \emph{plano de
fase}) elas percorrem curvas fechadas. Com $h=0{,}1$, o RK4 reproduz a órbita fechada; o
método de Euler, com o mesmo passo, gera uma espiral que se afasta dela
(Figura~\ref{fig:lv}): em $t=20$, Euler dá $(u,v)\approx(3{,}32;\ 0{,}50)$, enquanto o RK4 dá
$(10{,}21;\ 1{,}33)$. Mesmo com $h=0{,}01$, Euler ainda erra visivelmente.
\end{exemplo}

\begin{figure}[h]
\centering
\begin{tikzpicture}
\begin{axis}[width=10cm,height=6.5cm,xlabel={$u$ (presas)},ylabel={$v$ (predadores)},
  grid=both,grid style={gray!20},xmin=0,xmax=15,ymin=0,ymax=7.5,
  legend style={at={(0.98,0.98)},anchor=north east,font=\small}]
\addplot[thick,vermelho] coordinates {(10.000,5.000)(8.168,5.535)(6.432,5.911)(4.949,6.111)(3.772,6.152)(2.881,6.067)(2.224,5.892)(1.745,5.658)(1.398,5.389)(1.145,5.102)(0.960,4.809)(0.823,4.519)(0.723,4.237)(0.649,3.965)(0.595,3.706)(0.557,3.460)(0.531,3.229)(0.515,3.012)(0.509,2.809)(0.510,2.620)(0.519,2.443)(0.536,2.279)(0.560,2.127)(0.592,1.986)(0.632,1.856)(0.683,1.736)(0.744,1.626)(0.817,1.524)(0.905,1.431)(1.009,1.347)(1.132,1.270)(1.278,1.201)(1.450,1.139)(1.653,1.085)(1.892,1.038)(2.174,0.997)(2.505,0.965)(2.892,0.940)(3.345,0.923)(3.873,0.916)(4.485,0.919)(5.189,0.935)(5.994,0.965)(6.902,1.013)(7.910,1.084)(9.003,1.185)(10.150,1.325)(11.294,1.515)(12.345,1.772)(13.176,2.113)(13.634,2.552)(13.563,3.095)(12.871,3.725)(11.590,4.396)(9.900,5.033)(8.067,5.561)(6.342,5.927)(4.875,6.117)(3.716,6.151)(2.839,6.060)(2.193,5.881)(1.723,5.644)(1.382,5.373)(1.133,5.086)(0.951,4.793)(0.817,4.503)(0.718,4.221)(0.645,3.950)(0.592,3.692)(0.555,3.447)(0.530,3.217)(0.515,3.000)(0.509,2.798)(0.511,2.610)(0.520,2.434)(0.537,2.271)(0.561,2.119)(0.594,1.979)(0.635,1.849)(0.686,1.730)(0.747,1.620)(0.821,1.519)(0.910,1.427)(1.015,1.342)(1.140,1.266)(1.287,1.197)(1.461,1.136)(1.666,1.082)(1.907,1.035)(2.191,0.995)(2.524,0.963)(2.915,0.939)(3.372,0.923)(3.905,0.916)(4.521,0.920)(5.231,0.936)(6.042,0.967)(6.955,1.016)(7.968,1.089)(9.066,1.192)(10.214,1.334)};
\addplot[thick,gray,dashed] coordinates {(10.000,5.000)(8.172,5.570)(6.365,5.991)(4.790,6.227)(3.542,6.284)(2.615,6.196)(1.951,6.003)(1.485,5.745)(1.158,5.449)(0.928,5.137)(0.764,4.822)(0.647,4.512)(0.563,4.213)(0.502,3.927)(0.459,3.656)(0.428,3.401)(0.408,3.162)(0.397,2.939)(0.392,2.731)(0.395,2.537)(0.403,2.358)(0.417,2.191)(0.437,2.037)(0.464,1.895)(0.497,1.764)(0.538,1.642)(0.588,1.531)(0.649,1.429)(0.720,1.335)(0.805,1.250)(0.906,1.172)(1.026,1.101)(1.167,1.037)(1.334,0.980)(1.531,0.929)(1.764,0.885)(2.038,0.847)(2.361,0.815)(2.742,0.790)(3.190,0.772)(3.716,0.762)(4.331,0.760)(5.048,0.767)(5.878,0.787)(6.834,0.820)(7.923,0.872)(9.147,0.947)(10.495,1.053)(11.937,1.203)(13.412,1.411)(14.812,1.700)(15.969,2.099)(16.647,2.641)(16.569,3.354)(15.506,4.235)(13.435,5.222)(10.671,6.178)(7.794,6.933)(5.335,7.375)(3.527,7.500)(2.318,7.379)(1.549,7.101)(1.065,6.739)(0.760,6.338)(0.563,5.927)(0.433,5.523)(0.345,5.133)(0.285,4.763)(0.243,4.415)(0.213,4.089)(0.192,3.784)(0.177,3.501)(0.168,3.238)(0.162,2.995)(0.160,2.769)(0.160,2.561)(0.163,2.368)(0.169,2.190)(0.177,2.025)(0.188,1.873)(0.202,1.734)(0.219,1.604)(0.240,1.486)(0.266,1.376)(0.297,1.276)(0.333,1.183)(0.377,1.098)(0.429,1.020)(0.491,0.949)(0.565,0.884)(0.653,0.825)(0.759,0.771)(0.884,0.722)(1.035,0.678)(1.214,0.639)(1.429,0.605)(1.685,0.575)(1.992,0.549)(2.359,0.528)(2.798,0.512)(3.322,0.501)};
\addplot[only marks,mark=*,azul] coordinates {(10,5)};
\legend{RK4 ($h=0{,}1$),Euler ($h=0{,}1$),{$(u_0,v_0)$}}
\end{axis}
\end{tikzpicture}
\caption{Plano de fase do modelo de Lotka-Volterra para $0\le t\le20$.}
\label{fig:lv}
\end{figure}

\begin{atividade}{sistemas de EDOs com RK4}{\arqSistemas}
O aplicativo resolve um sistema autônomo $u'=F(u,v)$, $v'=G(u,v)$ pelo RK4. As funções $F$ e
$G$ podem ser digitadas; a caixa \emph{Gráficos u(t) e v(t)} troca o plano de fase pelos
gráficos em função do tempo.
\begin{enumerate}[nosep]
  \item Confira os valores do exemplo em $t=20$ ($h=0{,}1$, $n=200$), com e sem \emph{Comparar
        com Euler}.
  \item Nos gráficos $u(t)$ e $v(t)$, estime o período das oscilações. Por que o pico dos
        predadores vem depois do pico das presas?
  \item Digite $F=v$ e $G=-u$ (oscilador harmônico $y''=-y$), com $u_0=1$, $v_0=0$. A órbita
        exata é a circunferência $u^2+v^2=1$. Compare RK4 e Euler com $h=0{,}1$ e $n=200$.
  \item Pêndulo: $F=v$ e $G=-9{,}8\operatorname{sen}(u)$ (digite \texttt{-9.8 sin(u)}), com
        $u_0=1$ e $v_0=0$. Compare com a aproximação linear $G=-9{,}8u$.
\end{enumerate}
\end{atividade}

% ==================================================================
\section{Problemas de valor de contorno}

Num PVI, todas as condições são dadas no mesmo ponto. Em muitos problemas, porém, as condições
aparecem nos dois extremos de um intervalo: a deflexão de uma viga apoiada nas pontas, a
temperatura numa barra com as extremidades mantidas a temperaturas fixas, a forma de um cabo
suspenso.

\begin{definicao}
Um \emph{problema de valor de contorno} (PVC) de segunda ordem é
\[
y''=g(x,y,y'),\qquad a\le x\le b,\qquad y(a)=\alpha,\quad y(b)=\beta.
\]
\end{definicao}

Diferentemente do PVI, um PVC pode não ter solução ou ter infinitas: $y''=-y$,
$y(0)=0$, $y(\pi)=1$ não tem solução (todas as soluções com $y(0)=0$ são $c\operatorname{sen}x$
e valem 0 em $\pi$), enquanto $y''=-y$, $y(0)=0$, $y(\pi)=0$ tem infinitas. Para o caso
linear há um resultado simples:

\begin{teorema}\label{teo:pvclinear}
Se $p$, $q$ e $r$ são contínuas em $[a,b]$ e $q(x)>0$ em $[a,b]$, o PVC linear
\[
y''=p(x)\,y'+q(x)\,y+r(x),\qquad y(a)=\alpha,\quad y(b)=\beta,
\]
tem uma única solução.
\end{teorema}

Estudaremos dois métodos: o \emph{método do tiro}, que transforma o PVC em PVIs (e reaproveita
o RK4), e o \emph{método das diferenças finitas}, que transforma o PVC num sistema linear (e
reaproveita o Capítulo~\ref{cap:sistemas}).

% ==================================================================
\section{Método do tiro}

\subsection{A ideia}

Se conhecêssemos $y'(a)$, o PVC seria um PVI. O método do tiro ``chuta'' esse valor:
para cada $s$, resolvemos
\[
y''=g(x,y,y'),\qquad y(a)=\alpha,\qquad y'(a)=s,
\]
por RK4 (como sistema de primeira ordem) e olhamos onde a solução, que chamamos $y(x;s)$,
chega em $x=b$. Como um canhão que corrige a inclinação depois de cada tiro, ajustamos $s$
até acertar o alvo $\beta$. Em termos do Capítulo~\ref{cap:zeros}, procuramos um zero da função
\[
\phi(s)=y(b;s)-\beta.
\]
Cada avaliação de $\phi$ exige resolver um PVI inteiro, e não temos uma fórmula para
$\phi'(s)$; por isso o \textbf{método da secante} é a escolha natural:
\[
s_{k+1}=s_k-\phi(s_k)\,\frac{s_k-s_{k-1}}{\phi(s_k)-\phi(s_{k-1})}.
\]

\begin{observacao}
Se a EDO é \emph{linear}, $y(b;s)$ é uma função afim de $s$ (combinação de duas soluções, por
superposição). Então $\phi$ é uma reta e a secante encontra o zero exatamente no primeiro
passo: bastam dois tiros e uma combinação.
\end{observacao}

\begin{exemplo}[: um PVC não linear]
Considere
\[
y''=\frac{1}{8}\bigl(32+2x^3-y\,y'\bigr),\qquad 1\le x\le3,\qquad y(1)=17,\quad y(3)=\tfrac{43}{3},
\]
cuja solução exata é $y=x^2+\frac{16}{x}$ (assim, $y'(1)=-14$). Usando RK4 com $h=0{,}1$
(20 passos) em cada tiro, os chutes $s_0=-1{,}4$ e $s_1=s_0+1=-0{,}4$ e a secante:
\begin{center}
\renewcommand{\arraystretch}{1.25}
\begin{tabular}{c|r|r|r}
$k$ & $s_k$ & $y(3;s_k)$ & $\phi(s_k)$\\ \hline
0 & $-1{,}4$ & $20{,}4518393$ & $6{,}12$\\
1 & $-0{,}4$ & $20{,}8582730$ & $6{,}52$\\
2 & $-16{,}4541312$ & $12{,}8171044$ & $-1{,}52$\\
3 & $-13{,}4269919$ & $14{,}6667588$ & $3{,}3\times10^{-1}$\\
4 & $-13{,}9726751$ & $14{,}3495052$ & $1{,}6\times10^{-2}$\\
5 & $-14{,}0004910$ & $14{,}3331574$ & $-1{,}8\times10^{-4}$\\
6 & $-14{,}0001918$ & $14{,}3333334$ & $9{,}2\times10^{-8}$\\
7 & $-14{,}0001919$ & $14{,}3333333$ & $5\times10^{-13}$
\end{tabular}
\end{center}
O valor final de $s$ difere de $-14$ em $2\times10^{-4}$: o que resta é o erro do RK4 com
$h=0{,}1$, e não o da secante. Com esse $s$, o erro máximo da solução em $[1,3]$ é da ordem
de $10^{-4}$.
\end{exemplo}

\begin{figure}[h]
\centering
\begin{tikzpicture}
\begin{axis}[width=10cm,height=6.5cm,xlabel={$x$},ylabel={$y$},grid=both,grid style={gray!20},
  xmin=0.9,xmax=3.1,ymin=9.5,ymax=23,legend style={at={(0.02,0.98)},anchor=north west,font=\small},legend columns=2]
\addplot[thick,gray] coordinates {(1.00,17.0000)(1.10,16.8938)(1.20,16.8469)(1.30,16.8489)(1.40,16.8916)(1.50,16.9684)(1.60,17.0743)(1.70,17.2050)(1.80,17.3576)(1.90,17.5295)(2.00,17.7189)(2.10,17.9246)(2.20,18.1457)(2.30,18.3815)(2.40,18.6319)(2.50,18.8969)(2.60,19.1766)(2.70,19.4715)(2.80,19.7819)(2.90,20.1085)(3.00,20.4518)};
\addplot[thick,cL1] coordinates {(1.00,17.0000)(1.10,16.9839)(1.20,17.0100)(1.30,17.0710)(1.40,17.1613)(1.50,17.2763)(1.60,17.4126)(1.70,17.5675)(1.80,17.7387)(1.90,17.9250)(2.00,18.1251)(2.10,18.3383)(2.20,18.5644)(2.30,18.8032)(2.40,19.0548)(2.50,19.3195)(2.60,19.5977)(2.70,19.8900)(2.80,20.1971)(2.90,20.5196)(3.00,20.8583)};
\addplot[thick,cL3] coordinates {(1.00,17.0000)(1.10,15.5332)(1.20,14.3654)(1.30,13.4307)(1.40,12.6826)(1.50,12.0878)(1.60,11.6216)(1.70,11.2653)(1.80,11.0047)(1.90,10.8287)(2.00,10.7285)(2.10,10.6971)(2.20,10.7288)(2.30,10.8189)(2.40,10.9635)(2.50,11.1593)(2.60,11.4035)(2.70,11.6935)(2.80,12.0271)(2.90,12.4022)(3.00,12.8171)};
\addplot[thick,cL4] coordinates {(1.00,17.0000)(1.10,15.8074)(1.20,14.8685)(1.30,14.1297)(1.40,13.5524)(1.50,13.1084)(1.60,12.7763)(1.70,12.5397)(1.80,12.3860)(1.90,12.3049)(2.00,12.2885)(2.10,12.3300)(2.20,12.4242)(2.30,12.5666)(2.40,12.7535)(2.50,12.9819)(2.60,13.2490)(2.70,13.5528)(2.80,13.8914)(2.90,14.2632)(3.00,14.6668)};
\addplot[very thick,vermelho] coordinates {(1.00,17.0000)(1.10,15.7555)(1.20,14.7734)(1.30,13.9978)(1.40,13.3886)(1.50,12.9167)(1.60,12.5601)(1.70,12.3018)(1.80,12.1289)(1.90,12.0311)(2.00,12.0000)(2.10,12.0291)(2.20,12.1127)(2.30,12.2465)(2.40,12.4267)(2.50,12.6500)(2.60,12.9139)(2.70,13.2159)(2.80,13.5543)(2.90,13.9272)(3.00,14.3333)};
\addplot[only marks,mark=*,mark size=2.5pt] coordinates {(1,17)(3,14.3333)};
\legend{$s_0$,$s_1$,$s_2$,$s_3$,$s_7$}
\end{axis}
\end{tikzpicture}
\caption{Método do tiro: as soluções dos PVIs com $y'(1)=s_k$. O alvo é o ponto $(3,\frac{43}{3})$.}
\label{fig:tiro}
\end{figure}

\begin{atividade}{método do tiro}{\arqTiro}
O aplicativo resolve cada tiro por RK4 (com $z$ representando $y'$ no campo da função) e faz
sete iterações da secante. O controle $kk$ escolhe quantos tiros aparecem no gráfico.
\begin{enumerate}[nosep]
  \item Confira a tabela do exemplo e o erro da solução final em relação à exata.
  \item Mude o chute $s_0$ para $-20$ e para $10$. A secante ainda converge? Em quantas iterações?
  \item Aumente $n$ para 40. O $s$ final fica mais perto de $-14$? Por quê?
  \item PVC linear: digite $g=4(y-x)$, $a=0$, $b=1$, $y(a)=0$, $y(b)=2$, a solução exata
        $y=x+\operatorname{senh}(2x)/\operatorname{senh}(2)$ (digite \texttt{x + sinh(2x)/sinh(2)}) e
        $s_0=0$. Confirme que a secante acerta já em $s_2$.
\end{enumerate}
\end{atividade}

% ==================================================================
\section{Método das diferenças finitas}

\subsection{Aproximação das derivadas}

Dividimos $[a,b]$ em $N+1$ subintervalos de comprimento $h=\frac{b-a}{N+1}$, com nós
$x_i=a+ih$, $i=0,1,\dots,N+1$; os nós $x_1,\dots,x_N$ são os \emph{pontos internos}. Pela
fórmula de Taylor,
\[
y(x_{i\pm1})=y(x_i)\pm h\,y'(x_i)+\frac{h^2}{2}y''(x_i)\pm\frac{h^3}{6}y'''(x_i)+\frac{h^4}{24}y^{(4)}(\xi_\pm).
\]
Somando e subtraindo as duas expansões, obtemos as \emph{diferenças centradas}
\[
y''(x_i)=\frac{y(x_{i+1})-2y(x_i)+y(x_{i-1})}{h^2}+O(h^2),\qquad
y'(x_i)=\frac{y(x_{i+1})-y(x_{i-1})}{2h}+O(h^2).
\]

\subsection{O sistema linear}

Para o PVC linear $y''=p(x)y'+q(x)y+r(x)$, $y(a)=\alpha$, $y(b)=\beta$, trocamos as derivadas
pelas diferenças em cada ponto interno e $y(x_i)$ pela incógnita $w_i$
($w_0=\alpha$, $w_{N+1}=\beta$):
\[
\frac{w_{i+1}-2w_i+w_{i-1}}{h^2}=p_i\,\frac{w_{i+1}-w_{i-1}}{2h}+q_i\,w_i+r_i,
\]
com $p_i=p(x_i)$ etc. Multiplicando por $-h^2$ e agrupando:
\[
\boxed{\;\Bigl(-1-\frac h2p_i\Bigr)w_{i-1}+\bigl(2+h^2q_i\bigr)w_i+\Bigl(-1+\frac h2p_i\Bigr)w_{i+1}=-h^2r_i,\qquad i=1,\dots,N.\;}
\]
Na primeira e na última equação, $w_0=\alpha$ e $w_{N+1}=\beta$ são conhecidos e passam para o
segundo membro. O resultado é um sistema $N\times N$ \emph{tridiagonal}, que pode ser resolvido
por eliminação de Gauss sem trocas de linhas e com custo proporcional a $N$ (não a $N^3$).

\begin{observacao}
Se $q(x)\ge0$ e $h\,|p(x)|\le2$, a matriz é diagonalmente dominante: o sistema tem solução
única e os métodos de Gauss-Jacobi e Gauss-Seidel do Capítulo~\ref{cap:sistemas} convergem.
O erro do método é $O(h^2)$, como o das diferenças centradas: dobrar o número de
subintervalos divide o erro por 4.
\end{observacao}

\begin{exemplo}[: $y''=4(y-x)$, $y(0)=0$, $y(1)=2$]
Aqui $p=0$, $q=4$ e $r=-4x$; a solução exata é $y=x+\dfrac{\operatorname{senh}(2x)}{\operatorname{senh}2}$.
Com $N=3$ pontos internos, $h=0{,}25$ e $2+h^2q=2{,}25$; os segundos membros são
$-h^2r_i=0{,}25\,x_i$, isto é, $0{,}0625$; $0{,}125$; $0{,}1875$, e a última equação recebe
$+\beta=2$:
\[
\begin{pmatrix}2{,}25&-1&0\\-1&2{,}25&-1\\0&-1&2{,}25\end{pmatrix}
\begin{pmatrix}w_1\\w_2\\w_3\end{pmatrix}=
\begin{pmatrix}0{,}0625\\0{,}125\\2{,}1875\end{pmatrix}.
\]
A solução é $w_1=0{,}3951247$, $w_2=0{,}8265306$, $w_3=1{,}3395692$; os valores exatos são
$0{,}3936767$, $0{,}8240271$ e $1{,}3370861$, com erro máximo $2{,}5\times10^{-3}$. Aumentando $N$:
\begin{center}
\renewcommand{\arraystretch}{1.25}
\begin{tabular}{c|c|c}
$N$ & $h$ & erro máximo\\ \hline
3 & $0{,}25$ & $2{,}50\times10^{-3}$\\
4 & $0{,}2$ & $1{,}73\times10^{-3}$\\
9 & $0{,}1$ & $4{,}39\times10^{-4}$\\
19 & $0{,}05$ & $1{,}10\times10^{-4}$
\end{tabular}
\end{center}
De $h=0{,}1$ para $h=0{,}05$, o erro cai pelo fator $4$, como previsto.
\end{exemplo}

\begin{atividade}{diferenças finitas para PVC linear}{\arqDifFin}
O aplicativo monta e resolve o sistema tridiagonal. À direita aparecem a matriz aumentada
(para $N\le8$), a tabela com $x_i$, $w_i$, a solução exata e o erro, e o erro máximo.
\begin{enumerate}[nosep]
  \item Com $N=3$, confira a matriz e a solução do exemplo.
  \item Faça $N=4,9,19$ e confira a tabela de erros.
  \item Digite $p(x)=-2/x$, $q(x)=2/x^2$, $r(x)=\operatorname{sen}(\ln x)/x^2$, $a=1$, $b=2$,
        $y(a)=1$, $y(b)=2$ (desmarque \emph{Mostrar a solução exata}, que se refere ao exemplo
        anterior). Como as entradas da matriz mudam de uma linha para outra?
  \item Compare com o método do tiro aplicado ao mesmo PVC linear (atividade anterior, item 4).
        Qual método dá menor erro com $h=0{,}1$?
\end{enumerate}
\end{atividade}

% ==================================================================
\section{Comparação dos métodos}

\begin{center}
\renewcommand{\arraystretch}{1.4}
\begin{tabular}{l|c|p{3.4cm}|p{5.2cm}}
método & ordem & custo & observações\\ \hline
Euler & 1 & 1 avaliação de $f$ por passo & simples; base para entender os outros; exige $h$ pequeno\\
Heun / ponto médio & 2 & 2 avaliações por passo & previsão e correção\\
RK4 & 4 & 4 avaliações por passo & o método de uso geral mais popular\\
tiro (com RK4) & 4 & um PVI por iteração & reaproveita métodos de PVI; pode divergir\\
diferenças finitas & 2 & um sistema tridiagonal & robusto para PVC lineares
\end{tabular}
\end{center}

\paragraph{Em resumo.}
\begin{itemize}[nosep]
  \item Para \textbf{PVI}, os métodos de Runge-Kutta são de passo único (só usam $y_i$ para
        calcular $y_{i+1}$), fáceis de programar e de mudar de passo; o RK4 é o padrão em
        problemas não rígidos. Programas profissionais usam pares de Runge-Kutta de ordens
        diferentes para estimar o erro e ajustar $h$ automaticamente (passo adaptativo).
  \item A \textbf{estabilidade} limita o passo, independentemente da precisão desejada. Para
        problemas rígidos são necessários métodos implícitos.
  \item Para \textbf{PVC}, o método do tiro aproveita toda a precisão do RK4, mas pode ser
        sensível ao chute inicial e ao crescimento das soluções; as diferenças finitas levam
        a sistemas lineares bem comportados e se estendem naturalmente às equações
        diferenciais parciais.
  \item Tudo depende dos capítulos anteriores: Taylor e erros (Capítulo~\ref{cap:pontoflutuante}),
        secante (Capítulo~\ref{cap:zeros}), sistemas tridiagonais (Capítulo~\ref{cap:sistemas})
        e regras de integração (Capítulo~\ref{cap:integracao}).
\end{itemize}

% ==================================================================
\section{Exercícios}

\begin{enumerate}[label=\textbf{\arabic*.},leftmargin=*]
\item Aplique o método de Euler com $h=0{,}1$ ao PVI $y'=-2xy$, $y(0)=1$, até $x=0{,}5$, e
  compare com a solução exata $y=e^{-x^2}$.

\item Para o mesmo PVI, faça dois passos ($h=0{,}1$) pelo método de Heun e pelo RK4.

\item Mostre que, quando $f(x,y)=g(x)$ não depende de $y$, um passo do RK4 é a regra $1/3$ de
  Simpson e um passo de Heun é a regra dos trapézios no intervalo $[x_i,x_{i+1}]$.

\item Qual é o maior passo $h$ para o qual o método de Euler aplicado a $y'=-50y$, $y(0)=1$,
  produz aproximações que decaem para zero?

\item Escreva como sistema de primeira ordem: (a) $y''+4y'+3y=0$; (b) o pêndulo amortecido
  $\theta''+0{,}5\,\theta'+9{,}8\operatorname{sen}\theta=0$; (c) $y'''=x\,y-y'$.

\item Faça um passo do RK4 com $h=0{,}1$ para o oscilador $u'=v$, $v'=-u$, $u(0)=1$, $v(0)=0$,
  e compare com $(\cos0{,}1;\ -\operatorname{sen}0{,}1)$.

\item Verifique que $y=(x+1)^2-\tfrac12e^x$ é solução do PVI $y'=y-x^2+1$, $y(0)=0{,}5$.

\item Use diferenças finitas com $N=1$ (um único ponto interno) para o PVC $y''=4(y-x)$,
  $y(0)=0$, $y(1)=2$, e compare $w_1$ com $y(0{,}5)$.

\item Para o PVC linear $y''=4(y-x)$, $y(0)=0$, $y(1)=2$, determine o valor exato de $y'(0)$ e
  explique por que o método do tiro com a secante o encontra (a menos do erro do RK4) já em $s_2$.

\item \textbf{(GeoGebra)} No aplicativo de Euler, use $f(x,y)=-20y$, $y_0=1$, $n=10$ e
  $h=0{,}1$. Explique o resultado usando o fator $1+h\lambda$.

\item \textbf{(GeoGebra)} No aplicativo de sistemas, use o modelo de Lotka-Volterra com Euler e
  $h=0{,}1$, $n=200$ e depois $h=0{,}01$, $n=2000$. O ponto final de Euler se aproxima do RK4?
  Com que ordem?

\item \textbf{(GeoGebra)} No aplicativo de diferenças finitas, faça $N=4$, $9$ e $19$ e calcule
  as razões entre os erros máximos. Que ordem elas indicam?
\end{enumerate}

\subsection*{Respostas}
\begin{enumerate}[label=\textbf{\arabic*.},leftmargin=*,nosep]
\item $y_1=1$; $y_2=0{,}98$; $y_3=0{,}9408$; $y_4=0{,}884352$; $y_5=0{,}8136038$; exato
  $e^{-0{,}25}=0{,}7788008$ (erro $0{,}035$).
\item Heun: $y_1=0{,}99$, $y_2=0{,}960696$. RK4: $y_1=0{,}9900498$ (erro $4\times10^{-10}$),
  $y_2=0{,}9607894$ (erro $4\times10^{-9}$); exatos $0{,}9900498$ e $0{,}9607894$.
\item Com $f=g(x)$: $k_1=g(x_i)$, $k_2=k_3=g(x_i+\frac h2)$, $k_4=g(x_i+h)$, logo
  $y_{i+1}-y_i=\frac h6[g(x_i)+4g(x_i+\frac h2)+g(x_{i+1})]$, a regra $1/3$ com espaçamento
  $\frac h2$. Em Heun, $k_1=g(x_i)$ e $k_2=g(x_i+h)$, logo $\frac h2(k_1+k_2)$ é o trapézio.
\item $|1-50h|<1\iff h<0{,}04$.
\item (a) $u'=v$, $v'=-4v-3u$. (b) $u'=v$, $v'=-0{,}5v-9{,}8\operatorname{sen}u$.
  (c) $u'=v$, $v'=w$, $w'=xu-v$.
\item $(0{,}99500417;\ -0{,}09983333)$; exato $(0{,}99500417;\ -0{,}09983342)$.
\item $y'=2(x+1)-\frac12e^x$ e $y-x^2+1=2x+2-\frac12e^x$; além disso, $y(0)=1-\frac12=0{,}5$.
\item $h=0{,}5$: $(2+0{,}25\cdot4)\,w_1=0{,}25\cdot0{,}5\cdot4+0+2$, isto é, $3w_1=2{,}5$, logo
  $w_1=0{,}8333$; o exato é $0{,}8240$ (erro $0{,}0093$).
\item $y'(0)=1+\dfrac{2}{\operatorname{senh}2}\approx1{,}5514411$. A EDO é linear, logo
  $\phi(s)$ é afim e a secante dá o zero exato de $\phi$ no primeiro passo.
\item $1+h\lambda=1-20\cdot0{,}1=-1$: as aproximações alternam entre $-1$ e $1$ e não decaem
  (limite da estabilidade, $h=2/|\lambda|$).
\item Com $h=0{,}1$, Euler termina em $(3{,}32;\ 0{,}50)$; com $h=0{,}01$, em $(9{,}74;\ 1{,}20)$,
  mais perto do RK4 $(10{,}21;\ 1{,}33)$: a distância ao ponto do RK4 cai de cerca de $6{,}9$
  para $0{,}5$. Mesmo com 10 vezes mais passos, Euler continua bem menos preciso que o RK4
  com $h=0{,}1$.
\item As razões ficam perto de $1{,}73\times10^{-3}/4{,}39\times10^{-4}\approx3{,}9$ e
  $4{,}39\times10^{-4}/1{,}10\times10^{-4}\approx4{,}0$ (o passo cai à metade): ordem 2.
\end{enumerate}


\end{document}
